% Généralités sur les circuits électriques domestiques — PCT 3ème — Cours signé OnBuch+
% Programme officiel MINESEC. Compiler avec : tectonic N21-circuits-electriques-domestiques.tex
\newcommand{\DOCMATIERE}{PCT}
\newcommand{\DOCNIVEAU}{3ème}
\newcommand{\DOCMODULE}{Module 4 : Technologie et mécanique}
\newcommand{\DOCLECON}{N21}
\newcommand{\DOCTITRE}{Généralités sur les circuits électriques domestiques}
% OnBuch+ — préambule « cours signés » ENRICHI (compilé avec tectonic / XeLaTeX)
% Système visuel « Pop » d'OnBuch+ : fond crème (jamais blanc), bordures encre
% épaisses, ombres « dures » décalées, titres Archivo Black en capitales.
% Version MATHÉMATIQUES : graphes (pgfplots/tikz), boîtes pédagogiques
% (définition, propriété, méthode, attention, le savais-tu, exercice résolu,
% exercice, TP, bilan), page de garde et signature OnBuch+ ancrée en bas.
%
% Macros attendues dans main.tex AVANT \input{preamble} :
%   \DOCMATIERE  \DOCNIVEAU  \DOCMODULE  \DOCLECON (numéro)  \DOCTITRE
\documentclass[11pt]{article}

\usepackage{fontspec}
\usepackage[french]{babel}
\usepackage[a4paper,margin=2cm,top=3.4cm,bottom=2.8cm,headheight=1.4cm,footskip=1.3cm]{geometry}
\usepackage{amsmath,amssymb,amsfonts}
\usepackage{mathtools} % \xmapsto, \coloneqq... souvent utilisés par les modèles
\usepackage{cancel} % simplifications barrées (bilans de forces)
% \abs{z} (module, valeur absolue) et \norm{u} (norme), utilisés par les modèles
\providecommand{\abs}{}\let\abs\relax\DeclarePairedDelimiter{\abs}{\lvert}{\rvert}
\providecommand{\norm}{}\let\norm\relax\DeclarePairedDelimiter{\norm}{\lVert}{\rVert}
\usepackage[table]{xcolor} % [table] : fournit aussi \rowcolors (alternance de lignes)
\usepackage{pagecolor}
\usepackage{tikz}
\usetikzlibrary{arrows.meta,positioning,calc,shapes.geometric,shapes.misc,shapes.arrows,decorations.pathmorphing,decorations.pathreplacing,decorations.markings,patterns,shadows,mindmap,trees,fit,backgrounds,matrix,intersections,angles,quotes}
\usepackage{pgfplots}
\usepackage[version=4]{mhchem} % équations nucléaires \ce{^{235}_{92}U}
\usepackage[european,siunitx]{circuitikz} % schémas électriques
\pgfplotsset{compat=1.18}
\usepgfplotslibrary{fillbetween,groupplots} % aires entre courbes (calcul intégral)
\usepackage{enumitem}
\usepackage{fancyhdr}
% [most] charge toutes les bibliothèques tcolorbox (listings, minted, raster,
% documentation...) et gonfle inutilement la mémoire TeX sur un document long
% (~300 boîtes) : on ne charge que ce qui est réellement utilisé.
\usepackage[skins,breakable]{tcolorbox}
\usepackage{graphicx}
\usepackage{colortbl}
\usepackage{booktabs}
\usepackage{tabularx}
\usepackage{multirow}
\usepackage{diagbox} % case d'en-tête diagonale des tableaux à double entrée
% Colonnes extensibles centrée / gauche / droite (C, L, R), souvent utilisées par les modèles
\newcolumntype{C}{>{\centering\arraybackslash}X}
\newcolumntype{L}{>{\raggedright\arraybackslash}X}
\newcolumntype{R}{>{\raggedleft\arraybackslash}X}
\usepackage{array}
\usepackage{multicol}
\usepackage{siunitx}
% parse-numbers=false : accepte aussi \SI{2,5 \times 10^{-3}}{...}
\sisetup{locale=FR,output-decimal-marker={,},group-minimum-digits=4,parse-numbers=false}
\DeclareSIUnit\ohm{\ensuremath{\symup{\Omega}}} % U+2126 absent de la police
\DeclareSIUnit\division{div} % réglages d'oscilloscope (ms/div, V/div)
\DeclareSIUnit\tr{tr} % tours (vitesse de rotation, ex. \SI{1500}{\tr\per\minute})
\DeclareSIUnit\molar{M} % concentration molaire (ex. \SI{3}{\molar}, \SI{1}{\micro\molar})
\DeclareSIUnit\an{an} % année (le modèle confond parfois avec \year, un compteur TeX) — ex. \SI{5}{\kilo\meter\per\an}
\DeclareSIUnit\FCFA{FCFA} % franc CFA (ex. \SI{500}{\FCFA\per\kilo\gram})
\DeclareSIUnit\degreeBrix{\textdegree Brix} % teneur en sucre d'un sirop/jus (réfractométrie)
\DeclareSIUnit\month{mois} % le modèle utilise parfois \month, un compteur TeX, comme unité de durée
\DeclareSIUnit\microliter{\micro\liter} % le modèle écrit parfois \microliter en un seul token
\DeclareSIUnit\million{millions} % dénombrement (ex. \SI{15}{\million\per\mL} spermatozoïdes)
\usepackage{qrcode}
% \degree : commande fréquemment utilisée par les modèles pour un angle ou une
% température en dehors de \SI{}{} (non fournie par les paquets chargés).
\providecommand{\degree}{^{\circ}}
% \qed (fin de démonstration), utilisé par les modèles sans amsthm
\providecommand{\qed}{\hfill$\square$}
% \barre : double barre des valeurs interdites dans un tableau de variations (tkz-tab)
\providecommand{\barre}{\ensuremath{\|}}
% environnement proof (amsthm non chargé) utilisé par les modèles
\ifdefined\proof\else\newenvironment{proof}[1][Démonstration]{\par\noindent\textit{#1.}\ }{\qed\par}\fi
\usepackage{lastpage}
\usepackage{hyperref}
\hypersetup{hidelinks}

% ============================================================
% Palette « Pop » OnBuch+ — mêmes hex que la classe P (plus_ui.dart)
% ============================================================
\definecolor{poporange}{HTML}{F59321}
\definecolor{popdark}{HTML}{E07A0C}
\definecolor{popcream}{HTML}{FFF6E8}
\definecolor{popink}{HTML}{1C1714}
\definecolor{poppurple}{HTML}{7A5AE0}
\definecolor{popgreen}{HTML}{2E9E6A}
\definecolor{poppink}{HTML}{E0457B}
\definecolor{popblue}{HTML}{2F7DE1}
\definecolor{popgold}{HTML}{C98A12}
\definecolor{popmuted}{HTML}{8A7F76}
\definecolor{popline}{HTML}{EADFD2}
\definecolor{popgray}{HTML}{8A7F76}
\colorlet{popgrey}{popgray}
% Alias tolérants (le modèle nomme parfois les couleurs différemment)
\colorlet{popwhite}{white}
\colorlet{popblack}{popink}
\colorlet{popred}{poppink}
\colorlet{popyellow}{popgold}
% Teintes claires (fonds de tableaux/encadrés) — le modèle nomme parfois
% les couleurs avec un suffixe L (ex. popblueL) sans les avoir définies.
\colorlet{poporangeL}{poporange!20}
\colorlet{popdarkL}{popdark!20}
\colorlet{poppurpleL}{poppurple!20}
\colorlet{popgreenL}{popgreen!20}
\colorlet{poppinkL}{poppink!20}
\colorlet{popblueL}{popblue!20}
\colorlet{popgoldL}{popgold!20}
\colorlet{popredL}{popred!20}
\colorlet{popgrayL}{popgray!20}
% Teintes claires pour fonds de tableaux / zones
\colorlet{popblueL}{popblue!10!white}
\colorlet{poporangeL}{poporange!14!white}
\colorlet{popgreenL}{popgreen!12!white}
\colorlet{poppurpleL}{poppurple!12!white}
\colorlet{poppinkL}{poppink!12!white}
\colorlet{popgoldL}{popgold!14!white}

\pagecolor{popcream}

% ============================================================
% Typographie
% ============================================================
\defaultfontfeatures{Ligatures=TeX}
\setmainfont{PlusJakartaSans-Medium.ttf}[Path=../../../fonts/,BoldFont=PlusJakartaSans-Bold.ttf]
% Maths en sans-serif (Fira Math) avec chiffres et lettres droites de Plus
% Jakarta, pour que les formules se fondent dans le texte.
\usepackage[math-style=ISO,bold-style=ISO]{unicode-math}
\setmathfont{FiraMath-Regular.otf}[Path=../../../fonts/]
\setmathfont{PlusJakartaSans-Medium.ttf}[Path=../../../fonts/,range={up/num,up/{Latin,latin},\mathup}]
\usepackage{icomma} % virgule décimale sans espace en mode math
% Caractères mathématiques Unicode absents des polices de texte (Plus Jakarta,
% Archivo Black) : rendus par leur équivalent mathématique, en texte comme en
% formule — sinon ils s'affichent en carré vide (ex. « Arithmétique dans ℤ »).
\usepackage{newunicodechar}
\newunicodechar{ℝ}{\ensuremath{\mathbb{R}}}
\newunicodechar{ℤ}{\ensuremath{\mathbb{Z}}}
\newunicodechar{ℂ}{\ensuremath{\mathbb{C}}}
\newunicodechar{ℕ}{\ensuremath{\mathbb{N}}}
\newunicodechar{ℚ}{\ensuremath{\mathbb{Q}}}
\newunicodechar{𝔻}{\ensuremath{\mathbb{D}}}
\newunicodechar{α}{\ensuremath{\alpha}}
\newunicodechar{β}{\ensuremath{\beta}}
\newunicodechar{γ}{\ensuremath{\gamma}}
\newunicodechar{δ}{\ensuremath{\delta}}
\newunicodechar{ε}{\ensuremath{\varepsilon}}
\newunicodechar{θ}{\ensuremath{\theta}}
\newunicodechar{λ}{\ensuremath{\lambda}}
\newunicodechar{μ}{\ensuremath{\mu}}
\newunicodechar{σ}{\ensuremath{\sigma}}
\newunicodechar{τ}{\ensuremath{\tau}}
\newunicodechar{φ}{\ensuremath{\varphi}}
\newunicodechar{ψ}{\ensuremath{\psi}}
\newunicodechar{ω}{\ensuremath{\omega}}
\newunicodechar{Σ}{\ensuremath{\Sigma}}
% majuscules produites par \MakeUppercase dans les titres (α-aminés -> Α)
\newunicodechar{Α}{\ensuremath{\alpha}}
\newunicodechar{Β}{\ensuremath{\beta}}
\newunicodechar{Γ}{\ensuremath{\Gamma}}
\newunicodechar{Δ}{\ensuremath{\Delta}}
\newunicodechar{Ε}{\ensuremath{\varepsilon}}
\newunicodechar{Θ}{\ensuremath{\Theta}}
\newunicodechar{Λ}{\ensuremath{\Lambda}}
\newunicodechar{Μ}{\ensuremath{\mu}}
\newunicodechar{Τ}{\ensuremath{\tau}}
\newunicodechar{Φ}{\ensuremath{\Phi}}
\newunicodechar{Ψ}{\ensuremath{\Psi}}
\newunicodechar{Ω}{\ensuremath{\Omega}}
\newunicodechar{↦}{\ensuremath{\mapsto}}
\newunicodechar{∈}{\ensuremath{\in}}
\newunicodechar{∉}{\ensuremath{\notin}}
\newunicodechar{∧}{\ensuremath{\wedge}}
\newunicodechar{⇔}{\ensuremath{\Leftrightarrow}}
\newunicodechar{⟺}{\ensuremath{\Longleftrightarrow}}
\newunicodechar{⇒}{\ensuremath{\Rightarrow}}
\newunicodechar{⟹}{\ensuremath{\Longrightarrow}}
\newunicodechar{∘}{\ensuremath{\circ}}
\newunicodechar{∪}{\ensuremath{\cup}}
\newunicodechar{∩}{\ensuremath{\cap}}
\newunicodechar{≡}{\ensuremath{\equiv}}
\newunicodechar{⊂}{\ensuremath{\subset}}
\newunicodechar{⊥}{\ensuremath{\perp}}
\newunicodechar{∃}{\ensuremath{\exists}}
\newunicodechar{∀}{\ensuremath{\forall}}
\newunicodechar{∛}{\ensuremath{\sqrt[3]{\,}}}
\newunicodechar{∜}{\ensuremath{\sqrt[4]{\,}}}
\newunicodechar{⃗}{\ensuremath{{}^{\to}}}
\newunicodechar{ⁿ}{\ifmmode^{n}\else\textsuperscript{\ensuremath{n}}\fi}
\newunicodechar{ˣ}{\ifmmode^{x}\else\textsuperscript{\ensuremath{x}}\fi}
\newunicodechar{ᵇ}{\ifmmode^{b}\else\textsuperscript{\ensuremath{b}}\fi}
\newunicodechar{ʳ}{\ifmmode^{r}\else\textsuperscript{\ensuremath{r}}\fi}
\newunicodechar{ᵘ}{\ifmmode^{u}\else\textsuperscript{\ensuremath{u}}\fi}
\newunicodechar{ʸ}{\ifmmode^{y}\else\textsuperscript{\ensuremath{y}}\fi}
\newunicodechar{ᵃ}{\ifmmode^{a}\else\textsuperscript{\ensuremath{a}}\fi}
\newunicodechar{ᵉ}{\ifmmode^{e}\else\textsuperscript{\ensuremath{e}}\fi}
\newunicodechar{ᵏ}{\ifmmode^{k}\else\textsuperscript{\ensuremath{k}}\fi}
\newunicodechar{ᵖ}{\ifmmode^{p}\else\textsuperscript{\ensuremath{p}}\fi}
\newunicodechar{ᴺ}{\ifmmode^{N}\else\textsuperscript{\ensuremath{N}}\fi}
\newunicodechar{ᶜ}{\ifmmode^{c}\else\textsuperscript{\ensuremath{c}}\fi}
\newunicodechar{ʰ}{\ifmmode^{h}\else\textsuperscript{\ensuremath{h}}\fi}
\newunicodechar{ᵠ}{\ifmmode^{q}\else\textsuperscript{\ensuremath{q}}\fi}
\newunicodechar{ₙ}{\ifmmode_{n}\else\textsubscript{\ensuremath{n}}\fi}
\newunicodechar{ₐ}{\ifmmode_{a}\else\textsubscript{\ensuremath{a}}\fi}
\newunicodechar{ᵢ}{\ifmmode_{i}\else\textsubscript{\ensuremath{i}}\fi}
\newunicodechar{ᵣ}{\ifmmode_{r}\else\textsubscript{\ensuremath{r}}\fi}
\newunicodechar{ₖ}{\ifmmode_{k}\else\textsubscript{\ensuremath{k}}\fi}
\newunicodechar{ᵦ}{\ifmmode_{\beta}\else\textsubscript{\ensuremath{\beta}}\fi}
\newunicodechar{ₓ}{\ifmmode_{x}\else\textsubscript{\ensuremath{x}}\fi}
\newfontfamily\popdisplay{ArchivoBlack-Regular.ttf}[Path=../../../fonts/]
\newfontfamily\poplogo{SpaceGrotesk-Bold.ttf}[Path=../../../fonts/]
\newfontfamily\popsemi{PlusJakartaSans-SemiBold.ttf}[Path=../../../fonts/]
\newfontfamily\popxbold{PlusJakartaSans-ExtraBold.ttf}[Path=../../../fonts/]
\newfontfamily\popxboldls{PlusJakartaSans-ExtraBold.ttf}[Path=../../../fonts/,LetterSpace=3.0]

\setlist[enumerate]{leftmargin=1.7em,itemsep=5pt,topsep=4pt}
\setlist[itemize]{leftmargin=1.7em,itemsep=4pt,topsep=4pt,label={\color{poporange}\textbullet}}
\setlength{\parindent}{0pt}
\setlength{\parskip}{5pt}
\renewcommand{\arraystretch}{1.25}

% pgfplots : style Pop par défaut
\pgfplotsset{
  every axis/.append style={axis line style={popink,line width=1pt},tick style={popink},
    grid style={popline,line width=0.5pt},label style={font=\small},tick label style={font=\footnotesize}},
  popcurve/.style={poporange,line width=1.8pt,no markers,smooth},
}
% Même style disponible dans un \draw TikZ simple (hors pgfplots)
\tikzset{popcurve/.style={poporange,line width=1.8pt}}
\tikzset{popgrid/.style={popline,very thin}}
\colorlet{popcurve}{poporange} % le nom du style est parfois utilisé comme couleur

% ============================================================
% Pill / squiggle
% ============================================================
\newcommand{\poppill}[3]{%
  \tikz[baseline=-0.55ex]\node[draw=popink,line width=1.2pt,rounded corners=2.6pt,%
    fill=#2,inner xsep=6.5pt,inner ysep=3pt]{%
    \popxboldls\fontsize{7}{7}\selectfont\color{#3}\MakeUppercase{#1}};%
}
\newcommand{\popsquiggle}[1]{%
  \tikz[baseline=-0.4ex]{
    \draw[#1,line width=2.6pt,line cap=round]
      plot [smooth] coordinates {(0,0.09) (0.34,0.01) (0.68,0.21) (1.02,0.01) (1.36,0.10)};
  }%
}
% Mot-clé surligné (marqueur orange sous le texte)
\newcommand{\cle}[1]{{\popxbold\color{popdark}#1}}

% ============================================================
% En-tête / pied
% ============================================================
\pagestyle{fancy}
\fancyhf{}
\renewcommand{\headrulewidth}{0pt}
\renewcommand{\footrulewidth}{0pt}
\lhead{\raisebox{-0.15ex}{\poplogo\fontsize{13}{13}\selectfont\color{popink}OnBuch\color{poporange}+}}
\rhead{\poppill{\DOCMATIERE\ \textbullet\ \DOCNIVEAU\ \textbullet\ Leçon \DOCLECON}{white}{popink}}
\lfoot{\popxbold\fontsize{7.6}{7.6}\selectfont\color{popmuted}\MakeUppercase{Cours signé OnBuch+}\ \textbullet\ {\popsemi\DOCTITRE}}
\rfoot{\tikz[baseline=-0.6ex]\node[rounded corners=8.5pt,fill=popink,minimum height=17pt,minimum width=17pt,inner xsep=5pt,inner sep=0]{\popxbold\fontsize{8}{8}\selectfont\color{popcream}\thepage\,/\,\pageref*{LastPage}};}

% ============================================================
% Titres
% ============================================================
\newcommand{\chaptitre}[2]{%
  \begin{tcolorbox}[enhanced,colback=poporange,colframe=popink,boxrule=2.6pt,arc=11pt,
      drop shadow=popink,left=15pt,right=15pt,top=12pt,bottom=11pt,before skip=3pt,after skip=18pt]
    \poppill{#2}{popink}{popcream}\par\vspace{9pt}
    {\raggedright\hyphenpenalty=10000\exhyphenpenalty=10000\popdisplay\fontsize{22}{24}\selectfont\color{popcream}\MakeUppercase{#1}\par}
  \end{tcolorbox}
}
% Section numérotée : pastille orange avec le numéro + titre Archivo
\newcounter{popsec}
\newcommand{\coursec}[1]{%
  \stepcounter{popsec}\setcounter{popsub}{0}%
  \par\vspace{16pt}\noindent
  \begin{minipage}{\linewidth}
  \tikz[baseline=-0.7ex]\node[draw=popink,line width=1.6pt,fill=poporange,rounded corners=4pt,minimum size=22pt,inner sep=0]{\popdisplay\fontsize{12}{12}\selectfont\color{popcream}\thepopsec};%
  \hspace{8pt}{\popdisplay\fontsize{15}{17}\selectfont\color{popink}\MakeUppercase{#1}}\par
  \vspace{4pt}\noindent{\color{poporange}\rule{36pt}{3.6pt}}
  \end{minipage}\par\nopagebreak\vspace{8pt}
}
\newcounter{popsub}
\newcommand{\courssub}[1]{%
  \stepcounter{popsub}%
  \par\vspace{9pt}\noindent{\popxbold\fontsize{11.5}{13}\selectfont\color{popink}{\color{poporange}\thepopsec.\thepopsub}\hspace{6pt}#1}\par\nopagebreak\vspace{4pt}
}

% ============================================================
% Boîtes pédagogiques — même recette : fond blanc, bordure épaisse colorée,
% ombre dure encre, badge plein. Titre optionnel [#1].
% ============================================================
\tcbset{popbase/.style={breakable,enhanced,colback=white,boxrule=2.3pt,arc=9pt,
  shadow={3pt}{-3pt}{0pt}{popink},left=14pt,right=14pt,top=11pt,bottom=11pt,before skip=11pt,after skip=15pt,
  fontupper=\popsemi\fontsize{10.6}{14.5}\selectfont\color{popink},
  fontlower=\popsemi\fontsize{10.6}{14.5}\selectfont\color{popink}}}
% Coche dessinée (le glyphe ✓ n'existe pas dans les polices du document)
\newcommand{\popcheck}[1]{\tikz[baseline=-0.5ex]\draw[#1,line width=1.6pt,line cap=round,line join=round](-0.13,0.0)--(-0.04,-0.09)--(0.13,0.1);}
\providecommand{\coche}{\popcheck{popgreen}}
\providecommand{\resultat}[1]{\par\smallskip\noindent\fcolorbox{popgreen}{popgreenL}{\parbox{\dimexpr\linewidth-2\fboxsep-2\fboxrule\relax}{\textbf{Résultat.} #1}}\par\smallskip}
\newcommand{\popboxhead}[3]{\poppill{#1}{#2}{white}\ifx\relax#3\relax\else\hspace{6pt}{\popxbold\fontsize{10.6}{12}\selectfont\color{popink}#3}\fi\par\vspace{7pt}}

\newtcolorbox{aretenir}[1][]{popbase,colframe=popgreen,before upper={\popboxhead{À retenir}{popgreen}{#1}}}
\newtcolorbox{exemplebox}[1][]{popbase,colframe=poporange,before upper={\popboxhead{Exemple}{poporange}{#1}}}
\newtcolorbox{definition}[1][]{popbase,colframe=popblue,before upper={\popboxhead{Définition}{popblue}{#1}}}
\newtcolorbox{propriete}[1][]{popbase,colframe=poppurple,before upper={\popboxhead{Propriété}{poppurple}{#1}}}
\newtcolorbox{methode}[1][]{popbase,colframe=popgold,before upper={\popboxhead{Méthode}{popgold}{#1}}}
\newtcolorbox{attention}[1][]{popbase,colframe=poppink,before upper={\popboxhead{Attention}{poppink}{#1}}}
\newtcolorbox{experience}[1][]{popbase,colframe=popblue,colback=popblueL,before upper={\popboxhead{Expérience}{popblue}{#1}}}
% « Le savais-tu ? » : carte violette pleine (contexte, culture, Cameroun)
\newtcolorbox{savaistu}[1][]{popbase,colback=poppurple,colframe=popink,
  fontupper=\popsemi\fontsize{10.4}{14.2}\selectfont\color{white},
  before upper={\poppill{Le savais-tu\,?}{white}{poppurple}\ifx\relax#1\relax\else\hspace{6pt}{\popxbold\color{white}#1}\fi\par\vspace{7pt}}}
% Exercice résolu : énoncé en haut, \tcblower puis la solution sur fond crème
\newcounter{popexo}\newcounter{popexr}
\newtcolorbox{exoresolu}[1][]{popbase,colframe=poporange,colbacklower=popcream,
  segmentation style={popink,line width=1.2pt,dashed},
  before upper={\stepcounter{popexr}\popboxhead{Exercice résolu \thepopexr}{poporange}{#1}},
  before lower={\poppill{Solution}{popink}{popcream}\par\vspace{6pt}}}
% Exercice (non résolu en place) — la correction est dans la partie Corrigés
\newtcolorbox{exercice}[1][]{popbase,colframe=popink,
  before upper={\stepcounter{popexo}\popboxhead{Exercice \thepopexo}{popink}{#1}}}
% Corrigé
\newtcolorbox{corrige}[1][]{popbase,colframe=popgreen,colback=popgreenL,
  before upper={\popboxhead{Corrigé}{popgreen}{#1}}}
% Objectifs / prérequis / bilan
\newtcolorbox{objectifs}{popbase,colframe=popink,colback=poporangeL,
  before upper={\popboxhead{Objectifs de la leçon}{popink}{}}}
\newtcolorbox{prerequis}{popbase,colframe=popink,colback=popblueL,
  before upper={\popboxhead{Prérequis}{popblue}{}}}
\newtcolorbox{bilan}{popbase,colframe=popink,colback=white,boxrule=2.8pt,
  before upper={\popboxhead{Fiche bilan}{poporange}{L'essentiel à connaître pour l'examen}}}
% Figure encadrée avec légende
\newtcolorbox{popfigure}[1][]{enhanced,breakable=false,colback=white,colframe=popline,boxrule=1.4pt,arc=8pt,
  left=10pt,right=10pt,top=10pt,bottom=8pt,before skip=10pt,after skip=12pt,halign=center,
  lower separated=false,halign lower=center,
  fontlower=\popsemi\fontsize{9}{11.5}\selectfont\color{popmuted},
  #1}
\newcommand{\legende}[1]{\par\vspace{6pt}{\popsemi\fontsize{9}{11.5}\selectfont\color{popmuted}#1}}

% ============================================================
% Signature OnBuch+
% ============================================================
\newcommand{\poplogoinline}[1]{{\poplogo\fontsize{#1}{#1}\selectfont\color{popink}OnBuch\color{poporange}+}}
\newcommand{\signatureonbuch}{%
  \begin{tcolorbox}[enhanced,colback=popink,colframe=popink,boxrule=2.2pt,arc=10pt,
      drop shadow=poporange,left=14pt,right=14pt,top=10pt,bottom=10pt,before skip=0pt,after skip=0pt]
    \tikz[baseline=-0.7ex]\node[circle,fill=popgreen,draw=popcream,line width=1.3pt,minimum size=22pt,inner sep=0]{\popcheck{white}};%
    \hspace{9pt}\begin{minipage}[c]{0.62\linewidth}
      {\popxbold\fontsize{12}{13}\selectfont\color{popcream}Signé\ \textbullet\ {\poplogo OnBuch\color{poporange}+}}\par\vspace{2pt}
      {\raggedright\popsemi\fontsize{8}{10.5}\selectfont\color{popline}\MakeUppercase{Équipe pédagogique OnBuch+}\\\MakeUppercase{Conforme au programme officiel MINESEC}\par}
    \end{minipage}\hfill
    {\poplogo\fontsize{22}{22}\selectfont\color{popcream}OnBuch\color{poporange}+}
  \end{tcolorbox}%
}

% ============================================================
% Page de garde
% #1 titre · #2 matière · #3 niveau · #4 module · #5 n° leçon · #6 durée ·
% #7 description / objectifs (texte ou itemize)
% ============================================================
\newcommand{\pagegarde}[7]{%
  \thispagestyle{empty}
  \newgeometry{a4paper,margin=1.7cm,top=1.6cm,bottom=1.4cm}%
  \begin{tikzpicture}[remember picture,overlay]
    \fill[poporange!18] ([xshift=-3.2cm,yshift=-3.6cm]current page.north east) circle (4.6cm);
    \fill[poppurple!14] ([xshift=2.4cm,yshift=7.6cm]current page.south west) circle (3.8cm);
  \end{tikzpicture}%
  {\poplogo\fontsize{30}{30}\selectfont\color{popink}OnBuch\color{poporange}+}\hfill
  \raisebox{0.6ex}{\poppill{Cours signé \textbullet\ Premium}{popink}{popcream}}\par
  \vspace{1pt}\popsquiggle{poppurple}\par
  \vspace{8pt}
  \begin{tcolorbox}[enhanced,colback=poporange,colframe=popink,boxrule=2.8pt,arc=13pt,
      drop shadow=popink,left=18pt,right=18pt,top=12pt,bottom=13pt,after skip=10pt]
    \poppill{#2 \textbullet\ #3}{popink}{popcream}\hspace{5pt}\poppill{#4}{white}{popink}\par\vspace{8pt}
    {\popdisplay\fontsize{14}{14}\selectfont\color{popink}LEÇON #5}\par\vspace{4pt}
    {\raggedright\hyphenpenalty=10000\exhyphenpenalty=10000\popdisplay\fontsize{25}{27}\selectfont\color{popcream}\MakeUppercase{#1}\par}
    \vspace{8pt}
    {\popxbold\fontsize{9.5}{9.5}\selectfont\color{popink}\MakeUppercase{Durée conseillée : #6}}
  \end{tcolorbox}
  \begin{tcolorbox}[enhanced,colback=white,colframe=popink,boxrule=2.2pt,arc=10pt,
      drop shadow=popink,left=16pt,right=16pt,top=10pt,bottom=10pt,after skip=8pt]
    \poppill{Dans cette leçon}{poppurple}{white}\par\vspace{5pt}
    \popsemi\fontsize{9.3}{12.6}\selectfont\color{popink}#7
  \end{tcolorbox}
  \noindent\begin{minipage}[t]{0.487\linewidth}
    \begin{tcolorbox}[enhanced,colback=popgreen,colframe=popink,boxrule=2.2pt,arc=10pt,
        drop shadow=popink,left=11pt,right=11pt,top=10pt,bottom=10pt,height=3.1cm,valign=center]
      \tcbox[colback=white,colframe=popink,boxrule=1.3pt,arc=5pt,boxsep=0pt,nobeforeafter,
        left=3pt,right=3pt,top=3pt,bottom=3pt,valign=center]{\qrcode[height=1.9cm]{https://play.google.com/store/apps/details?id=com.luvvix.onbuch&hl=fr}}%
      \hspace{8pt}\begin{minipage}[c]{0.5\linewidth}
        {\popdisplay\fontsize{11}{12.5}\selectfont\color{white}\MakeUppercase{Télécharge OnBuch}}\par\vspace{4pt}
        {\raggedright\popsemi\fontsize{8.4}{11}\selectfont\color{white}Gratuit \textbullet\ Google Play\\Tuteur IA, annales, concours, résultats\par}
      \end{minipage}
    \end{tcolorbox}
  \end{minipage}\hfill
  \begin{minipage}[t]{0.487\linewidth}
    \begin{tcolorbox}[enhanced,colback=poppurple,colframe=popink,boxrule=2.2pt,arc=10pt,
        drop shadow=popink,left=11pt,right=11pt,top=10pt,bottom=10pt,height=3.1cm,valign=center]
      \tikz[baseline=-0.7ex]\node[circle,fill=white,draw=popink,line width=1.3pt,minimum size=1.6cm,inner sep=0]{\popdisplay\fontsize{24}{24}\selectfont\color{poppurple}+};%
      \hspace{8pt}\begin{minipage}[c]{0.6\linewidth}
        {\popdisplay\fontsize{11}{12.5}\selectfont\color{white}\MakeUppercase{Passe à OnBuch+}}\par\vspace{4pt}
        {\raggedright\popsemi\fontsize{8.4}{11}\selectfont\color{white}Tous les cours signés, Tuteur IA illimité, fiches PDF — onglet OnBuch+ de l'app\par}
      \end{minipage}
    \end{tcolorbox}
  \end{minipage}
  \vspace{8pt}
  \popcontenu
  \vfill
  \signatureonbuch
  \restoregeometry
  \newpage
}

% Rangée « Ce cours contient » (page de garde)
\newcommand{\poptile}[3]{% #1 couleur #2 gros texte #3 légende
  \begin{tcolorbox}[enhanced,colback=#1,colframe=popink,boxrule=2pt,arc=9pt,shadow={2.5pt}{-2.5pt}{0pt}{popink},
      left=6pt,right=6pt,top=9pt,bottom=9pt,halign=center,valign=center,height=1.75cm,width=\linewidth,nobeforeafter]
    {\popdisplay\fontsize{12.5}{13}\selectfont\color{white}\MakeUppercase{#2}}\par\vspace{3pt}
    {\centering\popsemi\fontsize{7.8}{9.5}\selectfont\color{white}#3\par}
  \end{tcolorbox}}
\newcommand{\popcontenu}{%
  \par\noindent{\popxboldls\fontsize{7.5}{7.5}\selectfont\color{popmuted}CE COURS SIGNÉ CONTIENT}\par\vspace{7pt}
  \noindent\begin{minipage}[t]{0.235\linewidth}\poptile{popblue}{Cours}{complet, illustré, pas à pas}\end{minipage}\hfill
  \begin{minipage}[t]{0.235\linewidth}\poptile{poporange}{Méthodes}{et exercices résolus}\end{minipage}\hfill
  \begin{minipage}[t]{0.235\linewidth}\poptile{poppink}{Exercices}{type examen + corrigés}\end{minipage}\hfill
  \begin{minipage}[t]{0.235\linewidth}\poptile{popgreen}{Fiche bilan}{l'essentiel à retenir}\end{minipage}\par}

% Sommaire
\newtcolorbox{sommaire}{enhanced,colback=white,colframe=popink,boxrule=2.2pt,arc=10pt,
  drop shadow=popink,left=18pt,right=18pt,top=13pt,bottom=13pt,before skip=6pt,after skip=20pt,
  before upper={\poppill{Sommaire}{poppurple}{white}\par\vspace{9pt}\popsemi\fontsize{10.6}{14.5}\selectfont\color{popink}}}

% Signature de fin de cours — poussée tout en bas de la dernière page
\newcommand{\findecours}{%
  \par\vfill\nopagebreak
  \begin{center}\popsquiggle{poporange}\end{center}\vspace{4pt}
  \signatureonbuch
}
\AtBeginDocument{\def\llbracket{\mathopen{[\mkern-3.5mu[}}\def\rrbracket{\mathclose{]\mkern-3.5mu]}}} % intervalles d'entiers (glyphe ⟦ absent de la police)
% Glyphes absents de Fira Math : redéfinis à partir de glyphes disponibles
\AtBeginDocument{%
  \def\setminus{\mathbin{\backslash}}\let\smallsetminus\setminus
  \def\vdots{\mathord{\vbox{\baselineskip3.2pt\lineskiplimit0pt\kern4pt\hbox{.}\hbox{.}\hbox{.}}}}%
  \def\ddots{\mathinner{\mkern1mu\raise6.5pt\hbox{.}\mkern2mu\raise3.6pt\hbox{.}\mkern2mu\raise0.7pt\hbox{.}\mkern1mu}}%
  \def\triangle{\mathord{\scalebox{0.9}{$\symup{\Delta}$}}}%
  \def\nabla{\mathord{\raisebox{\depth}{\scalebox{1}[-1]{$\symup{\Delta}$}}}}%
  \def\checkmark{\popcheck{popdark}}%
  \def\star{\mathbin{\ast}}\def\bigstar{\mathord{\ast}}%
  \def\diamond{\mathbin{\scalebox{0.6}{$\Diamond$}}}%
  \def\bigcup{\mathop{\mathchoice{\raisebox{-0.3em}{\scalebox{1.7}{$\cup$}}}{\scalebox{1.25}{$\cup$}}{\cup}{\cup}}\displaylimits}%
  \def\bigcap{\mathop{\mathchoice{\raisebox{-0.3em}{\scalebox{1.7}{$\cap$}}}{\scalebox{1.25}{$\cap$}}{\cap}{\cap}}\displaylimits}%
  \def\lesseqqgtr{\mathrel{\lessgtr}}\def\bigoplus{\mathop{\oplus}\displaylimits}%
}
\providecommand{\ssi}{si et seulement si} % abréviation fréquente
\AtBeginDocument{\providecommand{\wideparen}{\overparen}} % arc de cercle

%% FIGFIX-BEGIN v5 — correctifs globaux des figures (docs/onbuch-generation/tools/figfix)
\usepackage{adjustbox}\usepackage{etoolbox}\tcbuselibrary{raster}
% toute figure TikZ/pgfplots plus large que la ligne est réduite à la largeur disponible
\BeforeBeginEnvironment{tikzpicture}{\begin{adjustbox}{max width=\linewidth}}
\AfterEndEnvironment{tikzpicture}{\end{adjustbox}}
% légendes pgfplots lisibles par-dessus les courbes
\pgfplotsset{every axis/.append style={legend style={fill=white,fill opacity=0.92,text opacity=1,draw=black!15,inner sep=3pt}}}
% étiquettes d'arêtes (syntaxe edge["..."]) avec fond blanc
\tikzset{every edge quotes/.append style={fill=white,inner sep=1.2pt}}
%% FIGFIX-END
\begin{document}

\pagegarde{Généralités sur les circuits électriques domestiques}{PCT}{3ème}{Module 4 : Technologie et mécanique}{N21}{2 séances (1 h chacune)}{Cette leçon explore l'architecture d'une installation électrique domestique camerounaise, du tableau de répartition aux prises de courant. Elle permet d'identifier les conducteurs (phase, neutre, terre), de comprendre le rôle des protections (disjoncteur, différentiel) et de maîtriser l'utilisation du tournevis testeur en toute sécurité.
\vspace{4pt}
\begin{multicols}{2}\small
\begin{itemize}
\item À la fin de cette leçon, je sais décrire le cheminement de l'énergie électrique du compteur aux appareils récepteurs.
\item Je sais identifier et nommer les éléments du tableau de répartition (disjoncteur d'abonné, différentiel, divisionnaires).
\item Je sais repérer les bornes d'une prise standard (Phase, Neutre, Terre) et respecter le code couleur des fils (NF C 15-100).
\item Je sais expliquer le principe de fonctionnement d'un tournevis testeur et l'utiliser pour identifier la phase.
\item Je sais distinguer les dangers du courant électrique (électrisation, électrocution, incendie) et citer les gestes de premiers secours.
\item Je sais réaliser le schéma normalisé d'un circuit d'éclairage simple (va-et-vient) et d'une prise de courant avec terre.
\end{itemize}\end{multicols}}

\begin{sommaire}
\begin{multicols}{2}
\begin{enumerate}
\item 1. Architecture générale d'une installation domestique
\item 2. Les conducteurs : identification, sections et code couleur
\item 3. La prise de courant 2P+T : bornes, branchement et rôle de la terre
\item 4. Le tournevis testeur : principe, utilisation et limites
\item 5. Sécurité électrique : dangers, protections et premiers secours
\item 6. Lecture de plans, schémas unifilaires et mise en situation
\item Activité d'intégration
\item Méthodes et exercices résolus
\item Exercices
\item Corrigés des exercices
\item Fiche bilan
\end{enumerate}
\end{multicols}
\end{sommaire}

\begin{objectifs}
\begin{itemize}
\item À la fin de cette leçon, je sais décrire le cheminement de l'énergie électrique du compteur aux appareils récepteurs.
\item Je sais identifier et nommer les éléments du tableau de répartition (disjoncteur d'abonné, différentiel, divisionnaires).
\item Je sais repérer les bornes d'une prise standard (Phase, Neutre, Terre) et respecter le code couleur des fils (NF C 15-100).
\item Je sais expliquer le principe de fonctionnement d'un tournevis testeur et l'utiliser pour identifier la phase.
\item Je sais distinguer les dangers du courant électrique (électrisation, électrocution, incendie) et citer les gestes de premiers secours.
\item Je sais réaliser le schéma normalisé d'un circuit d'éclairage simple (va-et-vient) et d'une prise de courant avec terre.
\item Je sais calculer la section minimale d'un conducteur en fonction de l'intensité du disjoncteur de protection (application simplifiée).
\item Je sais appliquer les consignes de sécurité avant toute intervention sur une installation (coupure générale, VAT).
\end{itemize}
\end{objectifs}

\begin{prerequis}
\begin{itemize}
\item Notions de tension, courant, résistance, puissance (P = U × I) et énergie (3ème, début d'année).
\item Loi d'Ohm (U = R × I) et effets du courant sur le corps humain (seuils de perception, tétanisation, fibrillation).
\item Symboles normalisés des dipôles de base (générateur, résistance, interrupteur, lampe, moteur).
\item Règles de base du schéma unifilaire (représentation simplifiée d'un circuit).
\item Unités du Système International : Volt (V), Ampère (A), Watt (W), Ohm (Ω).
\end{itemize}
\end{prerequis}

\newpage

\coursec{Architecture générale d'une installation domestique}

\textbf{Situation d'accroche.} À Douala, dans le quartier Bonabéri, une famille vient d'acheter un réfrigérateur d'occasion. On le branche sur une vieille prise murale à deux trous, héritée d'une installation faite « à la va-vite » il y a vingt ans. Le soir même, le père pose la main sur le châssis métallique de l'appareil pour le déplacer : il reçoit une violente décharge et reste « collé » quelques secondes avant d'être projeté au sol. Heureusement, il s'en sort avec des brûlures légères. Que s'est-il passé ? Un fil dénudé à l'intérieur du réfrigérateur touchait la carcasse métallique : celle-ci s'est retrouvée portée à la tension du secteur, \SI{230}{V}. En touchant le châssis, le père a offert au courant un chemin vers le sol... à travers son propre corps. Deux protections auraient évité ce drame : une \cle{prise de terre} qui dévie le courant de fuite vers le sol, et un \cle{disjoncteur différentiel} qui détecte cette fuite et coupe le circuit en une fraction de seconde.

\textbf{Problématique.} Comment le courant arrive-t-il du réseau ENEO jusqu'à nos lampes, nos prises et nos appareils ? Quels organes protègent, à chaque étape, les personnes et les fils ? C'est toute l'\emph{architecture} de l'installation domestique que nous allons découvrir, du compteur jusqu'au moindre récepteur.

\begin{definition}[Installation électrique domestique]
Une \cle{installation électrique domestique} est l'ensemble des circuits électriques qui alimentent les récepteurs (lampes, prises, appareils) d'un logement, depuis le \cle{compteur} jusqu'aux bornes de raccordement de chaque appareil. Elle comprend le point de livraison, le tableau de répartition, les circuits divisionnaires et les liaisons de terre.
\end{definition}

\courssub{Le point de livraison : compteur et disjoncteur d'abonné}

\textbf{Intuition.} L'électricité est une marchandise que la société ENEO « livre » à ton domicile, comme l'eau est livrée par la CAMWATER. Il faut donc un point d'entrée officiel : c'est le \cle{point de livraison}. Il contient deux appareils essentiels.

\begin{definition}[Compteur et disjoncteur d'abonné]
Le \cle{compteur} est l'appareil qui mesure l'énergie électrique consommée par le logement, exprimée en kilowattheures (\si{kWh}). C'est la propriété du distributeur (ENEO) ; il est plombé et il est interdit d'y toucher.

Le \cle{disjoncteur d'abonné} (ou disjoncteur de branchement), placé juste après le compteur, est l'organe de coupure général de l'installation. Il possède un \cle{calibre} réglé selon l'abonnement souscrit, généralement entre \SI{10}{A} et \SI{60}{A}. Si la consommation dépasse ce calibre, il « saute » et coupe tout le logement.
\end{definition}

\textbf{Monophasé ou triphasé ?} Deux types d'alimentation existent :
\begin{itemize}
\item l'alimentation \cle{monophasée} : un fil de \textbf{phase} et un fil de \textbf{neutre}. La tension entre phase et neutre vaut $U = \SI{230}{V}$. C'est le cas de la grande majorité des logements camerounais ;
\item l'alimentation \cle{triphasée} : trois fils de phase (L1, L2, L3) et un neutre. On la réserve aux abonnés qui ont de gros besoins (atelier, pompe, gros fours de boulangerie).
\end{itemize}

\begin{savaistu}[230 V et 400 V au Cameroun]
Au Cameroun, la tension nominale est de \SI{230}{V} entre phase et neutre. En triphasé, la tension mesurée \emph{entre deux phases} vaut \SI{400}{V}, car $230 \times \sqrt{3} \approx 400$. Les appareils « triphasés » (pompes à eau puissantes, gros fours industriels) se branchent sur les trois phases, plus le neutre et la terre. C'est pourquoi certains quartiers disposent de poteaux portant quatre câbles : trois phases et un neutre.
\end{savaistu}

\begin{attention}[Le disjoncteur d'abonné n'est pas une protection des personnes]
Le disjoncteur d'abonné limite la \emph{puissance} souscrite et protège le réseau contre les gros défauts. Mais son calibre (par exemple \SI{30}{A}) est mille fois supérieur au courant mortel pour un être humain (environ \SI{30}{mA}) : il ne peut donc \textbf{jamais} te sauver d'une électrocution. C'est le rôle du différentiel \SI{30}{mA}, que nous verrons plus bas.
\end{attention}

\courssub{Le tableau de répartition et le disjoncteur différentiel 30 mA}

\textbf{Intuition.} Après le point de livraison, le courant arrive au \cle{tableau de répartition}, souvent fixé près de l'entrée de la maison. C'est le « chef d'orchestre » de l'installation : il reçoit le courant en un seul point et le distribue en plusieurs circuits indépendants, chacun protégé. Imagine un carrefour où un policier (le différentiel) surveille que tout ce qui entre ressort bien, et où chaque rue (chaque circuit) a sa propre barrière (le disjoncteur divisionnaire).

\begin{propriete}[Fonctionnement du disjoncteur différentiel 30 mA]
Le \cle{disjoncteur différentiel} compare en permanence le courant qui part par la phase, $I_{\text{phase}}$, et le courant qui revient par le neutre, $I_{\text{neutre}}$.
\begin{itemize}
\item En fonctionnement normal : $I_{\text{phase}} = I_{\text{neutre}}$, la différence est nulle, tout va bien.
\item En cas de défaut d'isolement (courant de fuite vers la terre, par exemple à travers le corps d'une personne) : une partie du courant ne revient pas par le neutre. Si la différence dépasse \SI{30}{mA}, le différentiel coupe l'alimentation en moins d'une fraction de seconde.
\end{itemize}
La sensibilité de \SI{30}{mA} (soit \SI{0,03}{A}) protège les \textbf{personnes} contre l'électrocution, car un courant inférieur à ce seuil traversant le corps reste généralement non mortel.
\end{propriete}

\begin{experience}[Observer le bouton test d'un différentiel]
\textbf{Matériel :} un tableau électrique domestique équipé d'un interrupteur différentiel \SI{30}{mA} (observation encadrée par un adulte ou le professeur).

\textbf{Protocole :}
\begin{enumerate}
\item Repérer le différentiel : c'est un module large, en tête de rangée, portant la mention « \SI{30}{mA} » et un bouton marqué « T » ou « Test ».
\item Noter quelles lampes et prises sont alimentées.
\item Appuyer sur le bouton « Test ».
\end{enumerate}

\textbf{Observations :} le différentiel déclenche instantanément (son levier bascule) et tous les circuits qu'il protège sont coupés : lampes éteintes, prises hors tension.

\textbf{Interprétation :} le bouton « Test » crée volontairement une petite fuite de courant interne, supérieure à \SI{30}{mA}, entre phase et neutre. Le différentiel la détecte et coupe : cela prouve qu'il fonctionne.

\textbf{Conclusion :} un différentiel doit être testé régulièrement (une fois par mois est recommandé). Un différentiel qui ne déclenche pas au test est dangereux et doit être remplacé.
\end{experience}

\courssub{Les circuits divisionnaires et leurs disjoncteurs}

\begin{definition}[Circuit divisionnaire]
Un \cle{circuit divisionnaire} est une branche de l'installation qui part du tableau et alimente un groupe de récepteurs de même nature (éclairage, prises, four...). Chaque circuit divisionnaire est protégé à son départ par un \cle{disjoncteur divisionnaire magnéto-thermique}.
\end{definition}

\textbf{Comment fonctionne un disjoncteur magnéto-thermique ?} Il porte bien son nom, car il combine deux protections :
\begin{itemize}
\item la protection \textbf{thermique} : une lame bimétallique chauffe et se courbe si le courant dépasse le calibre \emph{pendant longtemps} : c'est la protection contre les \cle{surcharges} (trop d'appareils branchés en même temps) ;
\item la protection \textbf{magnétique} : une bobine crée un champ magnétique qui déclenche \emph{instantanément} si le courant devient très grand : c'est la protection contre les \cle{courts-circuits}.
\end{itemize}
Le disjoncteur divisionnaire protège donc les \textbf{fils} de l'installation contre l'échauffement, cause première des incendies d'origine électrique.

\textbf{Circuits courants et circuits spécialisés.} La norme distingue deux familles de circuits, avec des calibres et des sections de fils adaptés :

\begin{center}
\begin{tabularx}{\linewidth}{|l|X|X|X|}
\hline
\rowcolor{popblueL} \textbf{Circuit} & \textbf{Calibre du disjoncteur} & \textbf{Section des fils} & \textbf{Exemples} \\
\hline
Éclairage & \SI{10}{A} & \SI{1,5}{mm^2} & Lampes, ventilateurs de plafond \\
\hline
Prises de courant & \SI{16}{A} & \SI{2,5}{mm^2} & Télévision, réfrigérateur, fer à repasser \\
\hline
Circuits spécialisés & \SI{20}{A} / \SI{32}{A} & \SI{2,5}{mm^2} / \SI{6}{mm^2} & Four seul, chauffe-eau (\SI{20}{A}) ; Cuisinière (four+plaques), climatiseur puissant (\SI{32}{A}) \\
\hline
\end{tabularx}
\end{center}

\textbf{Pourquoi des fils plus gros pour les gros appareils ?} Un four de \SI{3}{kW} sous \SI{230}{V} absorbe un courant $I = \dfrac{P}{U} = \dfrac{3000}{230} \approx \SI{13}{A}$, et une plaque de cuisson peut dépasser \SI{25}{A}. Un fil trop fin s'échaufferait dangereusement (effet Joule) : d'où la section de \SI{6}{mm^2} et un circuit \emph{dédié}, c'est-à-dire réservé à ce seul appareil.

\begin{attention}[Ne pas confondre différentiel et disjoncteur divisionnaire]
Erreur très fréquente au BEPC !
\begin{itemize}
\item Le \textbf{disjoncteur divisionnaire} (calibre 10, 16, 20 ou \SI{32}{A}) protège les \textbf{fils} contre les surcharges et les courts-circuits.
\item Le \textbf{disjoncteur différentiel} (sensibilité \SI{30}{mA}) protège les \textbf{personnes} contre les fuites de courant vers la terre.
\end{itemize}
Le différentiel \textbf{ne protège pas contre les surcharges} : il ne « voit » que les différences de courant entre phase et neutre. Les deux protections sont complémentaires et toutes les deux indispensables.
\end{attention}

\courssub{La liaison équipotentielle et la prise de terre}

\textbf{Intuition.} Revenons au drame de Bonabéri. Si la carcasse du réfrigérateur avait été reliée à la terre par un fil, le courant de fuite serait parti directement dans le sol (chemin facile, de faible résistance) au lieu d'attendre qu'une main humaine vienne le conduire. Et cette fuite importante aurait fait déclencher le différentiel. Double sécurité !

\begin{definition}[Liaison équipotentielle]
La \cle{liaison équipotentielle principale} consiste à raccorder par des conducteurs de protection (verts et jaunes) toutes les masses métalliques du logement (carcasses d'appareils, tuyauteries, charpentes métalliques) à une \cle{borne de terre} du tableau, elle-même reliée à un \cle{piquet de terre} planté dans le sol.

Dans la salle de bains, pièce où le corps mouillé est beaucoup plus conducteur, on réalise en plus une \cle{liaison équipotentielle locale} reliant entre eux tous les éléments métalliques accessibles (robinetterie, tuyaux, huisseries) afin qu'aucune différence de tension dangereuse ne puisse apparaître entre deux objets qu'on toucherait en même temps.
\end{definition}

\begin{aretenir}[Le duo de sécurité]
Prise de terre + différentiel \SI{30}{mA} = protection efficace des personnes. La terre offre au courant de fuite un chemin vers le sol ; le différentiel détecte cette fuite et coupe le circuit. L'un sans l'autre, la protection est incomplète.
\end{aretenir}

\courssub{Schémas unifilaire et multifilaire}

Pour représenter une installation, l'électricien utilise deux types de schémas :
\begin{itemize}
\item le \cle{schéma multifilaire} : on dessine \textbf{tous les fils} (phase, neutre, terre) séparément. Il est précis mais vite surchargé ;
\item le \cle{schéma unifilaire} : on représente l'ensemble des conducteurs d'un circuit par \textbf{un seul trait}, et les appareillages (compteur, disjoncteurs) par des symboles normalisés. C'est le schéma de synthèse, idéal pour comprendre l'architecture générale.
\end{itemize}

\begin{popfigure}
\begin{center}
\begin{tikzpicture}[font=\small, >=Stealth,
  equip/.style={draw=popdark, fill=popblueL, rounded corners=2pt, minimum height=0.9cm, align=center},
  div/.style={draw=popdark, fill=popgreenL, rounded corners=2pt, minimum height=0.8cm, align=center},
  rec/.style={draw=popdark, fill=poporangeL, rounded corners=2pt, minimum height=0.7cm, align=center}]
% Chaîne principale
\node[equip, align=center] (compteur) at (0,0){Compteur\\ ENEO};
\node[equip, right=1.1cm of compteur, align=center] (dab){Disjoncteur\\ d'abonné\\ \SI{30}{A}};
\node[equip, right=1.1cm of dab, fill=poppinkL, align=center] (diff){Différentiel\\ \SI{30}{mA}};
\draw[thick, popdark] (compteur) -- (dab);
\draw[thick, popdark] (dab) -- (diff);
% Arrivée réseau
\draw[thick, popdark] (-2.2,0) -- (compteur) node[midway, above]{Réseau \SI{230}{V}};
% Répartiteur
\node[circle, fill=popdark, inner sep=2pt] (bus) at ($(diff.east)+(1.0,0)$) {};
\draw[thick, popdark] (diff) -- (bus);
% Trois départs
\node[div] (d1) at ($(bus)+(2.4,2.0)$) {Disj. \SI{10}{A}\\ Éclairage};
\node[div] (d2) at ($(bus)+(2.4,0)$) {Disj. \SI{16}{A}\\ Prises};
\node[div] (d3) at ($(bus)+(2.4,-2.0)$) {Disj. \SI{32}{A}\\ Cuisinière};
\draw[thick, popdark] (bus) |- (d1.west);
\draw[thick, popdark] (bus) -- (d2.west);
\draw[thick, popdark] (bus) |- (d3.west);
% Récepteurs
\node[rec] (r1) at ($(d1)+(3.0,0)$) {Lampes\\ \SI{1,5}{mm^2}};
\node[rec] (r2) at ($(d2)+(3.0,0)$) {Prises 2P+T\\ \SI{2,5}{mm^2}};
\node[rec] (r3) at ($(d3)+(3.0,0)$) {Cuisinière\\ \SI{6}{mm^2}};
\draw[thick, popdark] (d1) -- (r1);
\draw[thick, popdark] (d2) -- (r2);
\draw[thick, popdark] (d3) -- (r3);
% Terre
\node[equip, fill=popgoldL] (terre) at ($(d3)+(0,-1.6)$) {Borne de terre};
\draw[thick, popgreen!70!black] (terre.south) -- ++(0,-0.5);
\draw[thick, popgreen!70!black] ($(terre.south)+(0,-0.5)+(-0.3,0)$) -- ++(0.6,0);
\draw[thick, popgreen!70!black] ($(terre.south)+(0,-0.62)+(-0.2,0)$) -- ++(0.4,0);
\draw[thick, popgreen!70!black] ($(terre.south)+(0,-0.74)+(-0.1,0)$) -- ++(0.2,0);
\draw[thick, popgreen!70!black, dashed] (r2.south) |- (terre.east);
\draw[thick, popgreen!70!black, dashed] (r3.south) -- (terre.north);
\end{tikzpicture}
\end{center}
\legende{Schéma unifilaire d'une installation domestique : le courant du réseau traverse le compteur, le disjoncteur d'abonné, puis le différentiel \SI{30}{mA}, avant d'être réparti vers trois circuits divisionnaires (éclairage \SI{10}{A}, prises \SI{16}{A}, cuisinière \SI{32}{A}). Les traits en pointillés verts représentent les conducteurs de protection reliant les masses des récepteurs à la borne de terre, elle-même reliée au piquet de terre (symbole en bas).}
\end{popfigure}

\begin{popfigure}
\begin{center}
\begin{tikzpicture}[font=\footnotesize]
% Coffret
\draw[fill=popcream, draw=popdark, thick, rounded corners=4pt] (0,0) rectangle (10.5,5.2);
\draw[fill=popblueL, draw=popdark] (0.3,4.6) rectangle (10.2,5.0);
\node at (5.25,4.8) {\textbf{Tableau de répartition — vue de face}};
% Rangée 1 : différentiel + divisionnaires
\draw[fill=poppinkL, draw=popdark] (0.5,2.6) rectangle (2.1,4.2);
\node[align=center] at (1.3,3.6) {Différentiel\\ \SI{30}{mA}};
\draw[fill=white, draw=popdark] (1.1,2.75) rectangle (1.5,3.1);
\node at (1.3,2.92) {\tiny T};
\foreach \x/\t in {2.5/{Éclair.\\ \SI{10}{A}}, 3.9/{Prises\\ \SI{16}{A}}, 5.3/{Cuisinière\\ \SI{32}{A}}, 6.7/{Chauffe-eau\\ \SI{20}{A}}, 8.1/{Clim\\ \SI{20}{A}}}{
  \draw[fill=popgreenL, draw=popdark] (\x,2.6) rectangle (\x+1.1,4.2);
  \node[align=center] at (\x+0.55,3.4) {\t};
  \draw[fill=white, draw=popdark] (\x+0.35,2.75) rectangle (\x+0.75,3.15);
}
% Peigne
\draw[fill=popgold, draw=popdark] (0.5,2.35) rectangle (9.2,2.55);
\node at (5.0,2.2) {Peigne d'alimentation (phase + neutre)};
% Rangée 2 : borniers
\draw[fill=popblue!40, draw=popdark] (0.5,1.3) rectangle (4.0,1.9);
\node at (2.25,1.6) {Bornier de phase (rouge)};
\draw[fill=popblueL, draw=popdark] (4.3,1.3) rectangle (7.8,1.9);
\node at (6.05,1.6) {Bornier de neutre (bleu)};
\draw[fill=popgreenL, draw=popdark] (8.1,1.3) rectangle (10.0,1.9);
\node[align=center] at (9.05,1.6) {\tiny Borne\\ \tiny de terre};
% Étiquettes
\draw[fill=white, draw=popmuted] (0.5,0.4) rectangle (10.0,1.0);
\node[align=left, anchor=west] at (0.7,0.7) {\tiny Repérage : D1 = chambres, D2 = prises salon, D3 = cuisine (cuisinière), D4 = chauffe-eau, D5 = climatiseur};
\end{tikzpicture}
\end{center}
\legende{Tableau modulaire type : en tête de rangée, le différentiel \SI{30}{mA} (rose) ; à sa suite, les disjoncteurs divisionnaires modulaires (verts) alimentés par un peigne d'alimentation (doré) qui distribue phase et neutre sans fils volants ; en bas, les borniers de phase, de neutre et de terre, et l'étiquette de repérage des circuits.}
\end{popfigure}

\begin{methode}[Identifier un circuit dans le tableau]
Tu ne sais pas quel disjoncteur protège la prise du salon ? Voici la démarche :
\begin{enumerate}
\item Couper (mettre sur « 0 ») le disjoncteur divisionnaire suspect.
\item Vérifier quelle prise ou quelle lampe ne fonctionne plus (tester avec une lampe ou un appareil branché).
\item Étiqueter clairement le disjoncteur (par exemple : « prises salon »), puis le réenclencher.
\end{enumerate}
Un tableau bien étiqueté fait gagner un temps précieux en cas de panne... et évite de couper le mauvais circuit pendant une réparation !
\end{methode}

\begin{exoresolu}[Choisir le bon disjoncteur]
\textbf{Énoncé.} Dans une maison de Yaoundé, on installe un chauffe-eau de puissance $P = \SI{2,2}{kW}$ sous la tension $U = \SI{230}{V}$. (a) Calculer le courant absorbé par le chauffe-eau. (b) Choisir parmi \SI{10}{A}, \SI{16}{A} et \SI{20}{A} le calibre du disjoncteur divisionnaire adapté. (c) Ce disjoncteur protège-t-il les personnes contre l'électrocution ? Justifier.
\tcblower
\textbf{Solution.} (a) La puissance électrique s'écrit $P = U \times I$, donc :
\begin{align*}
I &= \frac{P}{U} = \frac{2200}{230} \approx \SI{9,6}{A}.
\end{align*}
(b) Le calibre doit être supérieur au courant de fonctionnement : on choisit un disjoncteur de \SI{16}{A} (le calibre \SI{10}{A} serait trop juste et risquerait de déclencher à chaud). On câblera ce circuit en fils de \SI{2,5}{mm^2}.

(c) \textbf{Non.} Le disjoncteur divisionnaire de \SI{16}{A} protège les \emph{fils} contre les surcharges et les courts-circuits. Or un courant de fuite mortel pour une personne peut valoir seulement \SI{50}{mA}, valeur mille fois trop faible pour le faire déclencher. La protection des personnes est assurée par le \textbf{disjoncteur différentiel \SI{30}{mA}} placé en tête du groupe de circuits, associé à la prise de terre.
\end{exoresolu}

\begin{exercice}[À toi de jouer]
Une famille possède un abonnement dont le disjoncteur d'abonné est réglé à \SI{30}{A}. Un soir, fonctionnent en même temps : une climatisation ($P_1 = \SI{2,5}{kW}$), un fer à repasser ($P_2 = \SI{1,2}{kW}$), un réfrigérateur ($P_3 = \SI{0,3}{kW}$) et dix lampes de \SI{60}{W} chacune. La tension du réseau est \SI{230}{V}.
\begin{enumerate}
\item Calculer la puissance totale consommée.
\item Calculer le courant total appelé. Le disjoncteur d'abonné va-t-il déclencher ?
\item Proposer une solution si le disjoncteur déclenche souvent.
\end{enumerate}
\end{exercice}

\begin{corrige}[Exercice « À toi de jouer »]
1. Puissance des lampes : $10 \times 60 = \SI{600}{W} = \SI{0,6}{kW}$. Puissance totale :
\[ P = 2{,}5 + 1{,}2 + 0{,}3 + 0{,}6 = \SI{4,6}{kW}. \]
2. Courant total : $I = \dfrac{P}{U} = \dfrac{4600}{230} = \SI{20}{A}$. Comme $20 < 30$, le disjoncteur d'abonné \textbf{ne déclenche pas} : l'installation tient la charge.
3. Si les déclenchements sont fréquents (par exemple en ajoutant un four), on peut soit \emph{délester} (ne pas faire fonctionner en même temps les gros appareils), soit souscrire auprès d'ENEO un abonnement de puissance supérieure, ce qui fera monter le calibre du disjoncteur d'abonné.
\end{corrige}

\begin{aretenir}[L'essentiel de la section 1]
\begin{itemize}
\item Le courant suit le chemin : \textbf{réseau ENEO} $\rightarrow$ \textbf{compteur} $\rightarrow$ \textbf{disjoncteur d'abonné} (calibre 10 à \SI{60}{A}) $\rightarrow$ \textbf{différentiel \SI{30}{mA}} $\rightarrow$ \textbf{disjoncteurs divisionnaires} $\rightarrow$ récepteurs.
\item Le différentiel \SI{30}{mA} protège les \textbf{personnes} ; les disjoncteurs divisionnaires magnéto-thermiques protègent les \textbf{fils} (surcharges et courts-circuits).
\item Circuits courants : éclairage \SI{10}{A} / \SI{1,5}{mm^2}, prises \SI{16}{A} / \SI{2,5}{mm^2}. Circuits spécialisés : four/chauffe-eau \SI{20}{A} / \SI{2,5}{mm^2}, cuisinière/clim \SI{32}{A} / \SI{6}{mm^2}.
\item La liaison équipotentielle relie toutes les masses métalliques à la terre ; elle est renforcée localement dans la salle de bains.
\item Le schéma unifilaire résume l'installation avec un trait par circuit et des symboles normalisés.
\end{itemize}
\end{aretenir}

\coursec{Les conducteurs : identification, sections et code couleur}

\textbf{Situation d'accroche.} Tu ouvres une gaine électrique dans le grenier et tu vois trois fils : un bleu, un rouge, un vert-jaune. Lequel est la phase ? Lequel est le neutre ? Lequel est la terre ? Si tu les inverses en branchant une prise, l'appareil fonctionnera peut-être... mais la protection ne jouera plus son rôle. Connaître le code couleur, c'est la base de la sécurité.

\begin{definition}[Conducteur et isolant]
Un \cle{conducteur} (cuivre, aluminium) laisse passer le courant électrique. Un \cle{isolant} (PVC, caoutchouc, porcelaine) ne le laisse pas passer. Dans une installation, les fils de cuivre sont gainés d'isolant coloré pour être identifiés.
\end{definition}

\begin{propriete}[Code couleur normalisé (NF C 15-100)]
En France et au Cameroun, le code couleur des isolants est \textbf{strictement obligatoire} :
\begin{itemize}
\item \textbf{Bleu} : \cle{Neutre} (N). C'est le conducteur de retour du courant.
\item \textbf{Vert-Jaune} (rayé) : \cle{Terre} / Conducteur de protection (PE). \textbf{Jamais utilisé pour autre chose.}
\item \textbf{Rouge, Noir, Marron, Gris, Blanc} : \cle{Phase} (L). Ce sont les conducteurs « actifs » qui apportent l'énergie. \textbf{Interdiction formelle d'utiliser du bleu ou du vert-jaune pour la phase.}
\end{itemize}
\end{propriete}

\begin{attention}[Piège mortel : l'inversion Phase/Neutre]
Sur une prise, si tu inverses le bleu (neutre) et le rouge (phase), la lampe s'allume quand même. Mais l'interrupteur coupe le \emph{neutre} au lieu de la \emph{phase} : le filament de la lampe reste sous \SI{230}{V} même éteint ! Risque d'électrocution en changeant l'ampoule. \textbf{Respecte toujours le code couleur.}
\end{attention}

\courssub{Section des conducteurs et intensité admissible}

La \cle{section} d'un fil (en \si{mm^2}) détermine l'intensité maximale qu'il peut traverser sans s'échauffer excessivement. Le tableau ci-dessous donne les valeurs usuelles pour des fils en cuivre, gainés, posés en gaine (méthode de pose C) :

\begin{center}
\begin{tabularx}{\linewidth}{|c|X|X|}
\hline
\rowcolor{popblueL} \textbf{Section (\si{mm^2})} & \textbf{Intensité max. admise (A)} & \textbf{Usage courant} \\
\hline
1,5 & 10 à 13 & Éclairage (protégé par disj. 10 A) \\
\hline
2,5 & 16 à 20 & Prises 16 A, circuits spécialisés 20 A (chauffe-eau, four seul) \\
\hline
6 & 32 à 40 & Cuisinière, plaque de cuisson, climatiseur puissant (protégé par disj. 32 A) \\
\hline
10 & 50 à 60 & Arrivée générale, départ disjoncteur d'abonné \\
\hline
\end{tabularx}
\end{center}

\begin{aretenir}[Règle d'or : Section $\ge$ Calibre]
Le disjoncteur divisionnaire (calibre $I_n$) protège le fil. La section du fil doit être choisie pour supporter \textbf{au moins} $I_n$.
\begin{itemize}
\item Disj. \SI{10}{A} $\rightarrow$ Fil \SI{1,5}{mm^2} mini.
\item Disj. \SI{16}{A} ou \SI{20}{A} $\rightarrow$ Fil \SI{2,5}{mm^2} mini.
\item Disj. \SI{32}{A} $\rightarrow$ Fil \SI{6}{mm^2} mini.
\end{itemize}
\end{aretenir}

\courssub{Fils rigides vs fils souples}

\begin{itemize}
\item \textbf{Fil rigide (H07V-U)} : une seule âme de cuivre. Utilisé pour le câblage fixe dans les gaines (murs, plafonds). Ne se plie pas facilement, tient bien en place.
\item \textbf{Fil souple (H07V-R ou H07V-K)} : plusieurs brins fins. Utilisé pour les raccordements dans le tableau, les cordons d'alimentation d'appareils, les fils de terre. Se plie facilement mais \textbf{ne doit pas être mis sous vis directement} sans embout (risque de coupure des brins).
\end{itemize}

\begin{experience}[Reconnaître et préparer des fils]
\textbf{Matériel :} chutes de fils rigides (1,5 ; 2,5 ; 6 \si{mm^2}) et souples de même section, pince à dénuder, embouts à sertir, pinces à sertir, domino ou connecteur automatique (Wago).

\textbf{Protocole :}
\begin{enumerate}
\item Identifier la section en lisant l'inscription sur la gaine (ex: \texttt{2,5 mm²}).
\item Dénuder \SI{10}{mm} à \SI{12}{mm} d'isolant sans entailler le cuivre.
\item Sur fil souple : enficher un embout, sertir fermement.
\item Raccorder : phase sur phase (domino ou Wago), neutre sur neutre, terre sur terre.
\item Vérifier la tenue mécanique (tirer doucement) et l'absence de cuivre nu apparent.
\end{enumerate}

\textbf{Observation :} Un fil rigide se visse directement sous la vis d'une prise ou d'un domino. Un fil souple \textbf{doit} avoir un embout serti.
\end{experience}

\begin{exoresolu}[Choisir le fil pour un circuit]
\textbf{Énoncé.} On veut ajouter un circuit pour un climatiseur de \SI{3,5}{kW} en \SI{230}{V}. (a) Calculer le courant. (b) Choisir le calibre du disjoncteur et la section du fil. (c) Quelle couleur pour la phase, le neutre, la terre ?
\tcblower
\textbf{Solution.} (a) $I = P/U = 3500/230 \approx \SI{15,2}{A}$.
(b) Courant $> \SI{16}{A}$ ? Non, $15,2 < 16$. Mais pour un moteur (démarrage fort), on préfère un disjoncteur \SI{20}{A} (courbe D) ou \SI{16}{A} courbe C avec marge. Section fil : \SI{2,5}{mm^2} (supporte 20 A). Circuit dédié obligatoire.
(c) Phase : \textbf{Rouge} (ou Noir/Marron). Neutre : \textbf{Bleu}. Terre : \textbf{Vert-Jaune}.
\end{exoresolu}

\coursec{La prise de courant 2P+T : bornes, branchement et rôle de la terre}

\textbf{Situation d'accroche.} Tu as acheté une rallonge « 3 fils » pour ton ordinateur portable. Tu regardes la prise mâle : trois broches. La prise femelle au mur : trois trous. Pourquoi trois ? Les deux trous ronds (ou ovales) sont pour le courant « qui va et qui vient » (Phase et Neutre). Le troisième, souvent métallique et plus haut (ou au milieu), est la \cle{terre}. C'est elle qui te protège si l'ordinateur a un défaut d'isolement.

\begin{definition}[Prise 2P+T (Deux Pôles + Terre)]
C'est la prise standard française (type E) utilisée au Cameroun. Elle comporte trois bornes de raccordement à l'arrière :
\begin{itemize}
\item \textbf{Borne L (Phase)} : à \textbf{droite} quand on regarde la prise de face (trous vers soi). Bornier souvent marqué \texttt{L} ou \texttt{P}. Fil \textbf{Rouge/Noir/Marron}.
\item \textbf{Borne N (Neutre)} : à \textbf{gauche} quand on regarde la prise de face. Bornier souvent marqué \texttt{N}. Fil \textbf{Bleu}.
\item \textbf{Borne Terre (symbole $\perp$ ou $\downarrow$)} : en haut (ou au centre) de la prise, souvent métallique, reliée à la pince de terre de la prise. Fil \textbf{Vert-Jaune}.
\end{itemize}
\end{definition}

\begin{popfigure}
\begin{center}
\begin{tikzpicture}[font=\small, scale=1.2]
% Prise vue de face (trous)
\draw[fill=popcream, draw=popdark, thick, rounded corners=2pt] (-1.5,-1) rectangle (1.5,1.5);
% Trous
\draw[fill=popdark] (-0.6,0) circle (0.18); % Neutre gauche
\draw[fill=popdark] (0.6,0) circle (0.18);  % Phase droite
\draw[fill=popgray, draw=popdark] (0,0.7) circle (0.15); % Terre haut (métallique)
% Flèches et labels
\draw[popblue, thick, <-] (-0.6,-1.3) -- (-0.6,-0.5) node[below] {Neutre (N) \textbf{Bleu}};
\draw[poporange, thick, <-] (0.6,-1.3) -- (0.6,-0.5) node[below] {Phase (L) \textbf{Rouge/Noir}};
\draw[popgreen!70!black, thick, <-] (0,1.5) -- (0,1.0) node[above] {Terre \textbf{Vert-Jaune}};
% Légende orientation
\node[align=center, text width=3cm] at (0,-2.2) {\textbf{Vue de face (trous vers vous)} \\ Phase à DROITE, Neutre à GAUCHE, Terre en HAUT};
\end{tikzpicture}
\end{center}
\legende{Prise 2P+T vue de face : disposition des trous et code couleur des fils à raccorder à l'arrière. La terre est la broche saillante sur la prise mâle (le cordon) et le trou supérieur sur la prise femelle (le mur).}
\end{popfigure}

\begin{methode}[Brancher une prise 2P+T (côté câble)]
\begin{enumerate}
\item \textbf{Couper le courant} au disjoncteur divisionnaire concerné (vérifier avec tournevis testeur ou VAT).
\item Dénuder la gaine extérieure sur \SI{4}{cm} environ. Repérer les 3 fils.
\item Dénuder chaque fil sur \SI{10}{mm} (rigide) ou mettre embout + sertir (souple).
\item \textbf{Raccorder :} Bleu $\rightarrow$ borne N (Gauche), Rouge/Noir $\rightarrow$ borne L (Droite), Vert-Jaune $\rightarrow$ borne Terre (Haut/Centre).
\item Serrer fermement les vis (ou enficher dans Wago). Aucun fil de cuivre nu ne doit dépasser.
\item Refermer la prise, revisser le cache. Rétablir le courant.
\end{enumerate}
\end{methode}

\begin{attention}[Erreurs fatales sur prise 2P+T]
\begin{itemize}
\item \textbf{Oublier la terre} (ne pas connecter le vert-jaune) : l'appareil fonctionne, mais en cas de défaut, le châssis devient mortel et le différentiel ne déclenchera pas (pas de chemin de fuite).
\item \textbf{Inverser Phase et Neutre} : interrupteur sur le neutre, danger à l'entretien.
\item \textbf{Mettre la phase sur la borne terre} : \textbf{TRÈS GRAVE}. Toute la liaison équipotentielle (robinets, châssis, piquet de terre) passe sous \SI{230}{V}. Électrocution garantie au premier contact robinet/évier.
\end{itemize}
\end{attention}

\courssub{Le rôle de la terre : la protection par mise à la terre}

\begin{propriete}[Fonctionnement de la protection par terre]
\begin{enumerate}
\item Un défaut d'isolement survient : un fil de phase touche la carcasse métallique d'un appareil (classe I).
\item La carcasse est reliée à la borne terre de la prise par le fil vert-jaune du cordon.
\item Le courant de fuite traverse : Phase $\rightarrow$ Carcasse $\rightarrow$ Fil vert-jaune $\rightarrow$ Borne de terre du tableau $\rightarrow$ Piquet de terre $\rightarrow$ Sol.
\item Ce courant de fuite crée un déséquilibre détecté par le \textbf{différentiel \SI{30}{mA}} qui coupe l'alimentation.
\end{enumerate}
\textbf{Condition :} L'appareil doit être de \textbf{Classe I} (symbole double carré $\square\square$ barré ? Non, symbole \textbf{Classe I} = pas de double carré, borne de terre obligatoire). Les appareils \textbf{Classe II} (double isolation, symbole $\square\square$) n'ont pas de borne de terre (prise 2P sans T ou fiche plate).
\end{propriete}

\begin{savaistu}[La prise « Shuko » (type F) et l'adaptateur]
Au Cameroun, on trouve aussi des prises « Shuko » (allemandes, type F) : la terre est sur les deux bornes latérales métalliques, pas au centre. Si tu branches une fiche française (terre centrale) dans une Shuko \textbf{sans adaptateur correct}, la terre n'est \textbf{pas connectée} ! Danger. Utilise toujours un adaptateur muni de la pince de terre latérale.
\end{savaistu}

\coursec{Le tournevis testeur : principe, utilisation et limites}

\textbf{Situation d'accroche.} Tu veux changer un interrupteur. Tu as coupé le disjoncteur, mais tu veux être \emph{sûr} qu'il n'y a plus de tension sur les fils avant de toucher le cuivre. Tu sors ton tournevis testeur (celui avec une petite lampe néon dans le manche transparent). Tu touches la vis de la borne : la lampe ne s'allume pas. Es-tu en sécurité ? Pas forcément...

\begin{definition}[Tournevis testeur (détecteur de tension unipolaire)]
C'est un tournevis dont le manche transparent contient un \cle{néon} (lampe à décharge) en série avec une \cle{résistance de haute valeur} (environ \SI{1}{M\Omega}). Le contact métallique au bout du manche (pommeau) sert de seconde électrode via le corps de l'utilisateur.
\end{definition}

\begin{propriete}[Principe de fonctionnement]
\begin{enumerate}
\item Tu touches la pointe du tournevis sur un conducteur (borne, fil dénudé).
\item Si ce conducteur est la \textbf{Phase} (sous tension \SI{230}{V} par rapport à la terre) :
\item Le courant passe : Pointe $\rightarrow$ Néon $\rightarrow$ Résistance $\rightarrow$ Ton corps $\rightarrow$ Tes pieds $\rightarrow$ Sol (Terre) $\rightarrow$ Neutre du transformateur ENEO (relié à la terre).
\item Le courant est très faible ($I \approx \frac{230}{10^6} = \SI{0,23}{mA}$) : \textbf{non dangereux}, mais suffisant pour allumer le néon.
\item Si le conducteur est le \textbf{Neutre} (potentiel $\approx 0$ V par rapport à la terre) ou la \textbf{Terre} : pas de différence de tension $\rightarrow$ pas de courant $\rightarrow$ néon \textbf{éteint}.
\end{enumerate}
\end{propriete}

\begin{methode}[Utiliser un tournevis testeur en sécurité]
\begin{enumerate}
\item \textbf{Vérifier l'outil :} Néons intact, pas de fissure sur le manche, isolation de la lame parfaite. Tester sur une prise connue sous tension (phase) : le néon doit s'allumer.
\item \textbf{Position :} Tenir le tournevis par le manche isolant. \textbf{Toucher le contact métallique du pommeau avec ton pouce ou l'index} (contact corps-terre indispensable).
\item \textbf{Test :} Toucher la pointe sur la borne ou le fil à tester.
\item \textbf{Lecture :} Néon allumé = \textbf{Phase sous tension}. Néon éteint = Neutre, Terre, \textbf{ou Phase coupée... ou pile HS / contact corps manquant / isolation chaussures !}
\end{enumerate}
\end{methode}

\begin{attention}[Limites dangereuses du tournevis testeur — \textbf{À connaître pour le BEPC}]
\begin{itemize}
\item \textbf{Ne prouve PAS l'absence de tension.} Néon éteint $\neq$ « pas de danger ». Causes possibles : tu as des semelles isolantes (bottes caoutchouc, tapis), le néon est grillé, la résistance interne est coupée, tu ne touches pas le pommeau, le neutre est « flottant » (coupé en amont mais remonté par les récepteurs allumés).
\item \textbf{Ne mesure PAS la tension.} Il ne dit pas si c'est 230 V, 400 V ou 12 V.
\item \textbf{Ne détecte pas le neutre coupé sous tension.} Si le neutre est coupé en amont d'une lampe allumée, le fil neutre \emph{après} la lampe est à \SI{230}{V} (remonté par le filament). Le testeur s'allume sur le neutre ! $\rightarrow$ Piège mortel.
\item \textbf{Interdit pour la vérification d'absence de tension (VAT) avant intervention.} La norme exige un \textbf{Contrôleur de Tension Normé (VAT)} ou un multimètre (catégorie III/IV) pour valider l'absence de tension (consignation).
\end{itemize}
\end{attention}

\begin{aretenir}[Tournevis testeur : usage restreint]
Utile pour \textbf{repérer la phase} sur une prise ou un bornier \emph{hors tension de travail} (installation neuve, dépannage léger). \textbf{Interdit} pour valider une mise hors tension avant intervention (consignation). Dans ce cas : VAT ou Multimètre obligatoire.
\end{aretenir}

\coursec{Sécurité électrique : dangers, protections et premiers secours}

\textbf{Situation d'accroche.} Ton cousin branche un ventilateur sur une rallonge dénudée. Il pose le pied sur le fil nu. Il crie, se raidit, ne peut plus lâcher le fil. Son cœur s'emballe. Que faire ? \textbf{Ne le touche pas !} Tu deviendrais la deuxième victime. Coupe le courant. Appelle les secours.

\begin{definition}[Électrisation vs Électrocution]
\begin{itemize}
\item \textbf{Électrisation :} Passage d'un courant électrique dans le corps. Conséquences : brûlures (entrées/sorties), tétanisation (crampes musculaires involontaires), troubles du rythme cardiaque. La victime survit.
\item \textbf{Électrocution :} Électrisation \textbf{mortelle}. Arrêt cardiaque (fibrillation ventriculaire) ou arrêt respiratoire.
\end{itemize}
\end{definition}

\begin{propriete}[Effets du courant alternatif 50 Hz sur le corps humain]
\begin{center}
\begin{tabularx}{\linewidth}{|c|X|}
\hline
\rowcolor{popblueL} \textbf{Intensité (mA)} & \textbf{Effet physiologique} \\
\hline
0,5 -- 1 & Seuil de perception (picotement) \\
\hline
10 -- 15 & \textbf{Seuil de non-lâcher (tétanisation)} : impossible d'ouvrir la main, on reste « collé » \\
\hline
20 -- 25 & Difficultés respiratoires (tétanisation muscles respiratoires) \\
\hline
30 & \textbf{Seuil de fibrillation ventriculaire} (risque mortel majeur) $\rightarrow$ \textbf{Justification du \SI{30}{mA} du différentiel} \\
\hline
50 -- 100 & Fibrillation ventriculaire probable (cœur bat de façon chaotique, plus de circulation sanguine) \\
\hline
> 100 & Arrêt cardiaque, brûlures internes graves \\
\hline
\end{tabularx}
\end{center}
\end{propriete}

\begin{aretenir}[Facteurs aggravants]
La gravité dépend de : \textbf{Intensité} (tension / résistance du corps), \textbf{Durée} du contact, \textbf{Trajet} (main-pied traverse le cœur ; main-main aussi), \textbf{Etat de la peau} (mouillée = résistance divisée par 100 !), \textbf{Fréquence} (50 Hz plus dangereux que continu).
\end{aretenir}

\courssub{Les protections collectives et individuelles}

\begin{itemize}
\item \textbf{Protection par mise à la terre + Différentiel \SI{30}{mA}} (Classe I) : Coupe avant fibrillation.
\item \textbf{Double isolation (Classe II, symbole $\square\square$)} : Pas de pièce métallique accessible reliée à l'intérieur. Pas de terre nécessaire. (Ex: perceuse plastique, sèche-cheveux).
\item \textbf{Très Basse Tension de Sécurité (TBTS)} : \SI{12}{V} ou \SI{24}{V} (salle de bain, piscine, chantier). Alimentée par transformateur de séparation (isolement galvanique).
\item \textbf{Séparation électrique} : Un seul récepteur sur un transformateur d'isolement (poste à souder, lampe de chantier).
\item \textbf{Dispositifs de protection individuelle (EPI) :} Gants isolants, bottes isolantes, tapis isolant, visière (risque arc électrique). \textbf{Uniquement pour électriciens qualifiés (H0/B0 minimum).}
\end{itemize}

\courssub{Conduite à tenir en cas d'accident électrique : PROTÉGER - ALERTER - SECOURIR}

\begin{methode}[Premiers secours face à une électrisation]
\begin{enumerate}
\item \textbf{PROTÉGER (Toi d'abord) :} \textbf{NE TOUCHE PAS LA VICTIME} tant qu'elle est en contact avec le courant. Risque de « passage de relais » (le courant passe par toi).
\item \textbf{COUPER LE COURANT :} Disjoncteur général (disjoncteur d'abonné) $\rightarrow$ Position « 0 » / OFF. Si inaccessible : Dégager la victime avec un objet \textbf{SEC ET ISOLANT} (balai en bois, chaise en plastique, corde sèche, vêtement sec). \textbf{JAMAIS d'objet métallique, JAMAIS mains nues.}
\item \textbf{ALERTER :} Appeler le \textbf{112} (urgence européenne, fonctionne au Cameroun) ou \textbf{118} (pompiers) ou \textbf{117} (police). Dire : « Accident électrique, adresse, nombre de victimes, état (conscient, respire, brûlé) ».
\item \textbf{SECOURIR (si courant coupé) :}
    \begin{itemize}
    \item Victime \textbf{consciente et qui respire} : Allonger, couvrir, surveiller, rassurer. Ne pas donner à boire/manger. Surveiller les brûlures (refroidir \SI{15}{min} à l'eau tiède \SIrange{15}{25}{\celsius} si récentes).
    \item Victime \textbf{inconsciente mais qui respire} : Position Latérale de Sécurité (PLS). Surveiller la respiration.
    \item Victime \textbf{ne respire pas / pas de pouls} : \textbf{Massage cardiaque (RCP)} immédiatement : 30 compressions thoraciques / 2 insufflations (bouche-à-bouche ou masque) si formé. Sinon \textbf{compressions seules} (100-120/min) jusqu'à l'arrivée des secours ou reprise de la respiration. Défibrillateur (DAE) si disponible : suivre les instructions vocales.
    \end{itemize}
\end{enumerate}
\end{methode}

\begin{savaistu}[Brûlures électriques : particularité]
Une brûlure électrique a souvent une \textbf{porte d'entrée} (main) et une \textbf{porte de sortie} (pied). L'intérieur (muscles, tendons, nerfs, vaisseaux, os) est souvent bien plus détruit que la peau en surface. \textbf{Toute victime d'électrisation doit être vue par un médecin}, même si elle « va bien » : risque de troubles cardiaques retardés, nécrose interne, rhabdomyolyse (reins).
\end{savaistu}

\coursec{Lecture de plans, schémas unifilaires et mise en situation}

\textbf{Problématique.} Comment passer du schéma unifilaire du tableau (vu en section 1) au câblage réel d'un circuit d'éclairage ou de prises ? Comment lire le plan d'un électricien pour vérifier une installation ?

\courssub{Symboles normalisés (extrait NF C 03-013 / IEC 60617)}

\begin{center}
\begin{tabularx}{\linewidth}{|c|X|c|}
\hline
\rowcolor{popblueL} \textbf{Appareil} & \textbf{Description} & \textbf{Symbole schématique} \\
\hline
Interrupteur (allumage simple) & Coupe la phase vers la lampe & \tikz[baseline=-0.5ex]{\draw[thick] (0,0) -- (0.5,0); \draw[thick] (0.8,0.3) -- (0.5,0) -- (0.8,-0.3);} \\
\hline
Va-et-vient (2 inters pour 1 lampe) & Deux positions, un commun & \tikz[baseline=-0.5ex]{\draw[thick] (0,0) -- (0.5,0); \draw[thick] (0.5,0) -- (0.8,0.3); \draw[thick] (0.5,0) -- (0.8,-0.3); \draw[fill=white] (0.5,0) circle (0.05);} \\
\hline
Prise 2P+T & 3 bornes L, N, Terre & \tikz[baseline=-0.5ex]{\draw[thick] (0,0) rectangle (1,0.6); \draw[thick] (0.5,0.6) -- (0.5,0.9); \draw[fill=popdark] (0.5,0.9) circle (0.05);} \\
\hline
Lampe / Récepteur & Résistance / Charge & \tikz[baseline=-0.5ex]{\draw[thick] (0,0) circle (0.3); \draw[thick] (-0.3,0) -- (-0.6,0); \draw[thick] (0.3,0) -- (0.6,0); \node at (0,0) {$\times$};} \\
\hline
Disjoncteur divisionnaire & Protection magnéto-thermique & \tikz[baseline=-0.5ex]{\draw[thick] (0,0) -- (1,0); \draw[thick] (0.5,-0.3) rectangle (0.5,0.3);} \\
\hline
Différentiel 30 mA & Protection personnes & \tikz[baseline=-0.5ex]{\draw[thick] (0,0) -- (1,0); \draw[thick] (0.5,-0.3) rectangle (0.5,0.3); \node at (0.5,0) {30 mA};} \\
\hline
\end{tabularx}
\end{center}

\courssub{Schéma multifilaire : circuit d'éclairage simple (allumage)}

C'est le schéma de câblage réel. On voit les 3 fils (Phase, Neutre, Terre) partout.

\begin{popfigure}
\begin{center}
\begin{tikzpicture}[font=\small, >=Stealth, node distance=1.2cm and 1.5cm]
% Nœuds
\node (disj) {Disj. \SI{10}{A}};
\node[right=of disj] (inter) {Interrupteur};
\node[right=of inter] (lampe) {Lampe};
\node[below=1.5cm of disj] (bornes) {Borniers Tableau};
% Fils Phase (Rouge)
\draw[poporange, very thick] (disj) -- node[above] {L (Rouge)} (inter);
\draw[poporange, very thick] (inter) -- node[above] {L (Rouge)} (lampe);
% Fils Neutre (Bleu) - retour direct depuis bornier
\draw[popblue, very thick] (bornes) -- node[left] {N (Bleu)} ++(0,0.5) -| (lampe);
% Fils Terre (Vert-Jaune) - vers masse lampe si Classe I
\draw[popgreen!70!black, very thick, dashed] (bornes) -- node[left] {PE (V-J)} ++(0,-0.5) -| (lampe);
% Connexions bornier
\draw[poporange, very thick] (disj) |- (bornes);
\draw[popblue, very thick] (bornes) -- ++(-1,0);
\draw[popgreen!70!black, very thick, dashed] (bornes) -- ++(-1,0);
% Légendes bornier
\node[left=0.2cm of bornes, align=right] {Phase\\Neutre\\Terre};
\end{tikzpicture}
\end{center}
\legende{Schéma multifilaire d'un allumage simple. La phase (orange) est coupée par l'interrupteur. Le neutre (bleu) et la terre (vert-jaune pointillé) vont directement du bornier du tableau à la lampe. La terre n'est connectée que si la lampe est en métal (Classe I).}
\end{popfigure}

\courssub{Schéma multifilaire : circuit de prises 2P+T}

\begin{popfigure}
\begin{center}
\begin{tikzpicture}[font=\small, >=Stealth, node distance=1.5cm and 2cm]
% Tableau
\node (disj) {Disj. \SI{16}{A}};
\node[below=1cm of disj] (bornes) {Borniers};
% Prises en dérivation (en parallèle)
\node[right=of disj] (p1) {Prise 1};
\node[right=of p1] (p2) {Prise 2};
\node[right=of p2] (p3) {Prise 3};
% Phase
\draw[poporange, very thick] (disj) -- (p1) -- (p2) -- (p3);
% Neutre
\draw[popblue, very thick] (bornes) -- ++(0,0.5) -| (p1) (p1) -- (p2) -- (p3);
% Terre
\draw[popgreen!70!black, very thick, dashed] (bornes) -- ++(0,-0.5) -| (p1) (p1) -- (p2) -- (p3);
% Connexions bornier
\draw[poporange, very thick] (disj) |- (bornes);
\draw[popblue, very thick] (bornes) -- ++(-1,0);
\draw[popgreen!70!black, very thick, dashed] (bornes) -- ++(-1,0);
\node[left=0.2cm of bornes, align=right] {Phase\\Neutre\\Terre};
\end{tikzpicture}
\end{center}
\legende{Schéma multifilaire d'un circuit de prises (dérivation en parallèle). Phase, Neutre et Terre sont distribués de prise en prise (pontage sur les bornes). Maximum 8 prises sur un circuit \SI{16}{A} (NF C 15-100).}
\end{popfigure}

\courssub{Du plan au câblage : exercice de mise en situation}

\begin{exoresolu}[Lire un plan d'installation]
\textbf{Énoncé.} Sur le plan d'un appartement (schéma unifilaire ci-contre), on voit : Disjoncteur d'abonné 45 A $\rightarrow$ Différentiel 30 mA $\rightarrow$ 3 divisionnaires : D1 (10 A), D2 (16 A), D3 (20 A).
Le plan d'implantation montre : D1 alimente 5 points lumineux au salon et 3 en chambre. D2 alimente 6 prises au salon, 4 en cuisine, 2 en chambre. D3 alimente un chauffe-eau 3 kW dans la salle de bain.
\begin{enumerate}
\item Vérifier la conformité du nombre de points d'éclairage et de prises par circuit (Norme : max 8 points d'éclairage / circuit 10 A ; max 8 prises / circuit 16 A ; 1 seul appareil sur circuit spécialisé 20 A).
\item Le chauffe-eau 3 kW est-il bien protégé par du 20 A ? Calculer $I$.
\item Quel est le risque si on ajoute 3 prises sur le circuit D2 ?
\end{enumerate}
\tcblower
\textbf{Solution.}
1. D1 : 5+3 = 8 points lumineux. \textbf{Conforme} (limite 8). D2 : 6+4+2 = 12 prises. \textbf{NON CONFORME} (limite 8 prises sur circuit 16 A). Il faut créer un 2ème circuit prises (ex: D4 16 A) pour la cuisine ou les chambres.
2. $I = P/U = 3000/230 \approx \SI{13}{A}$. Disjoncteur \SI{20}{A} OK (13 < 20). Fil \SI{2,5}{mm^2} OK (supporte 20 A). Circuit dédié obligatoire. \textbf{Conforme}.
3. Risque : Surcharge. Si on branche gros appareils (fer, bouilloire, micro-ondes) sur les 15 prises, le courant peut dépasser \SI{16}{A} sans court-circuit. Le disjoncteur 16 A finira par déclencher (thermique) mais le fil \SI{2,5}{mm^2} peut chauffer anormalement avant le déclenchement si la surcharge est modérée mais prolongée. Risque d'incendie.
\end{exoresolu}

\begin{exercice}[Synthèse : Le tableau de M. Nkodo]
M. Nkodo à Yaoundé a un abonnement 45 A. Son tableau comporte : Différentiel 30 mA, puis 4 divisionnaires :
D1 : 10 A (Éclairage total 8 lampes LED 10 W).
D2 : 16 A (Prises salon/chambres : TV, déco, chargeurs).
D3 : 20 A (Four encastrable 2,5 kW).
D4 : 16 A (Machine à laver + Sèche-linge branchés sur une double prise).
\begin{enumerate}
\item Calculer le courant du four (D3). Le disjoncteur 20 A et le fil 2,5 mm² sont-ils adaptés ?
\item Calculer le courant maxi possible sur D4 si Machine à laver = 2,2 kW et Sèche-linge = 2,0 kW (fonctionnant ensemble). Que va-t-il se passer ?
\item Proposer une modification conforme pour D4.
\item Sur le schéma unifilaire, où placer le différentiel 30 mA par rapport aux divisionnaires ? Pourquoi ?
\end{enumerate}
\end{exercice}

\begin{corrige}[Exercice « Le tableau de M. Nkodo »]
1. Four : $I = 2500/230 \approx \SI{10,9}{A}$. Disj 20 A > 10,9 A : OK. Fil 2,5 mm² (supporte 20 A) : OK. Circuit dédié : OK.
2. D4 : $P_{tot} = 2,2 + 2,0 = \SI{4,2}{kW}$. $I = 4200/230 \approx \SI{18,3}{A}$.
   Le disjoncteur D4 est calibre \SI{16}{A}. $18,3 > 16$ $\rightarrow$ \textbf{Surcharge}. Le disjoncteur D4 va déclencher (protection thermique) au bout de quelques minutes. La machine à laver et le sèche-linge s'arrêtent. Le fil 2,5 mm² supporte 20 A, donc pas d'échauffement dangereux immédiat, mais coupure de service agaçante.
3. Modifier : Créer un circuit dédié \SI{20}{A} (fil 2,5 mm²) pour le sèche-linge (ou la machine). Ou passer D4 en \SI{20}{A} \textbf{SI} le fil est en \SI{2,5}{mm^2} (mais 18,3 A est proche de 20 A, marge faible). Mieux : 2 circuits distincts 16 A ou 20 A.
4. Le différentiel 30 mA est placé \textbf{en amont} de tous les divisionnaires (juste après le disjoncteur d'abonné). Il protège ainsi \textbf{tous} les circuits divisionnaires en aval. S'il était en aval d'un divisionnaire, ce circuit ne serait pas protégé contre les fuites vers la terre.
\end{corrige}

\begin{aretenir}[L'essentiel de la leçon complète (Séances 29-30)]
\begin{itemize}
\item \textbf{Architecture :} Réseau $\rightarrow$ Compteur $\rightarrow$ Disj. Abonné (Puissance) $\rightarrow$ Différentiel 30 mA (Personnes) $\rightarrow$ Disj. Divisionnaires (Fils) $\rightarrow$ Récepteurs.
\item \textbf{Conducteurs :} Bleu = Neutre, Vert-Jaune = Terre, Rouge/Noir/Marron = Phase. Sections : 1,5 (10A), 2,5 (16/20A), 6 (32A). Rigide en gaine, souple + embout en tableau/prise.
\item \textbf{Prise 2P+T :} Phase à DROITE, Neutre à GAUCHE, Terre en HAUT. Terre obligatoire pour appareils Classe I (métalliques).
\item \textbf{Tournevis testeur :} Détecte la Phase (néon allumé). \textbf{Ne prouve PAS l'absence de tension.} Interdit pour consignation (utiliser VAT/Multimètre).
\item \textbf{Sécurité :} Électrisation (blessure) vs Électrocution (mort). Seuil danger : 30 mA (fibrillation). Protection : Terre + Différentiel 30 mA + Double isolation (Classe II) + TBTS.
\item \textbf{Premiers secours :} 1. Ne pas toucher victime. 2. Couper courant (Disj. Général). 3. Alerter (112/118). 4. Secourir (PLS / Massage cardiaque / DAE).
\item \textbf{Schémas :} Unifilaire (architecture, 1 trait/circuit) vs Multifilaire (câblage réel, 3 fils). Symboles normalisés. Lecture de plans : vérifier calibres, sections, nombre de points par circuit.
\end{itemize}
\end{aretenir}

\coursec{Les conducteurs : identification, sections et code couleur}

Dans la section précédente, nous avons découvert l'architecture générale d'une installation domestique : le courant arrive du réseau ENEO, traverse le compteur et le disjoncteur général, puis se répartit au tableau de répartition vers les différents circuits de la maison (éclairage, prises, cuisson, chauffe-eau). Mais entre le tableau et la lampe de ta chambre, entre le tableau et la prise du salon, il y a des \cle{fils}. Ces fils ne sont pas tous pareils : ils n'ont ni la même couleur, ni la même épaisseur, ni le même rôle. C'est ce que nous allons étudier maintenant. Comprendre les conducteurs, c'est comprendre le \emph{système sanguin} de l'installation électrique : chaque « vaisseau » a sa fonction, son calibre et sa couleur, et aucune confusion n'est permise.

\courssub{Les trois conducteurs actifs : Phase, Neutre et Terre}

Ouvre une prise de courant débranchée (hors tension !) chez toi, ou observe le chantier d'une maison en construction dans ton quartier : tu verras toujours \textbf{trois fils} arriver dans chaque boîte. Pourquoi trois, alors qu'une lampe n'a besoin que de deux fils pour s'allumer ?

\begin{experience}[Observation : démonter une fiche de prise (hors tension)]
\textbf{Matériel :} une vieille fiche de prise hors d'usage, un tournevis, une installation \textbf{coupée au disjoncteur général}.

\textbf{Protocole :}
\begin{enumerate}
\item Vérifier que le disjoncteur général est sur OFF.
\item Dévisser le capot de la fiche.
\item Observer les fils connectés aux bornes et noter leurs couleurs.
\end{enumerate}

\textbf{Observations :} on trouve trois fils de couleurs différentes : un fil \textbf{bleu}, un fil \textbf{rouge} (ou brun, ou noir) et un fil \textbf{vert et jaune} (bicolore).

\textbf{Interprétation :} chaque couleur correspond à un rôle précis et normalisé. Le courant « arrive » par un fil, « repart » par un autre, et le troisième est un gardien de sécurité.

\textbf{Conclusion :} une installation domestique utilise trois conducteurs : la \cle{Phase} (L), le \cle{Neutre} (N) et la \cle{Terre} ou conducteur de protection (PE).
\end{experience}

\begin{definition}[Les trois conducteurs d'une installation domestique]
\begin{itemize}
\item La \cle{Phase} (repère \textbf{L}, de l'anglais \emph{Line}) : c'est le fil qui \textbf{amène} le courant depuis le tableau vers l'appareil. C'est le fil « dangereux » : il porte une tension de \SI{230}{V} par rapport à la terre.
\item Le \cle{Neutre} (repère \textbf{N}) : c'est le fil de \textbf{retour} du courant vers le transformateur ENEO. En régime normal, il est parcouru par le même courant que la phase.
\item La \cle{Terre} ou conducteur de \textbf{protection} (repère \textbf{PE}, \emph{Protective Earth}) : il relie les carcasses métalliques des appareils (frigo, fer à repasser, machine à laver) à la prise de terre plantée dans le sol.
\end{itemize}
\end{definition}

\begin{propriete}[Rôles respectifs du Neutre et de la Terre]
Le \textbf{Neutre} est un conducteur \cle{actif} : il transporte le courant en fonctionnement normal du circuit. Sans neutre, la lampe ne s'allume pas.

La \textbf{Terre (PE)} est un conducteur de \cle{protection} : elle ne transporte du courant \textbf{qu'en cas de défaut} d'isolement (par exemple quand un fil dénudé touche la carcasse métallique du frigo). Elle évacue alors ce courant de défaut vers le sol, ce qui fait déclencher le disjoncteur différentiel et protège les personnes.
\end{propriete}

\begin{attention}[Ne jamais confondre Neutre et Terre]
Même si le neutre et la terre sont tous deux reliés au sol quelque part dans le réseau, ils ne sont \textbf{jamais interchangeables}. Brancher la carcasse d'un appareil sur le neutre au lieu de la terre supprime la protection différentielle : en cas de défaut, la carcasse reste sous tension et peut électrocuter l'utilisateur. C'est une erreur \textbf{mortelle}, fréquente dans les installations « bricolées ».
\end{attention}

\courssub{Le code couleur obligatoire (norme NF C 15-100)}

Imagine un électricien de Douala appelé à réparer une panne dans une maison de Yaoundé, câblée par un autre artisan dix ans plus tôt. Si chacun choisissait ses couleurs au hasard, l'intervention serait une partie de roulette russe. C'est pourquoi la norme \cle{NF C 15-100} (appliquée au Cameroun comme en France) impose un \textbf{code couleur strict}.

\begin{aretenir}[Code couleur des conducteurs]
\begin{center}
\begin{tabularx}{\linewidth}{|l|X|X|}\hline
\rowcolor{popblueL} \textbf{Conducteur} & \textbf{Couleur(s) autorisée(s)} & \textbf{Couleurs interdites} \\ \hline
Phase (L) & Rouge, brun, noir (ou toute autre couleur sauf bleu et vert/jaune) & Bleu clair, vert/jaune \\ \hline
Neutre (N) & \textbf{Bleu clair} (obligatoire) & Toutes les autres \\ \hline
Terre (PE) & \textbf{Vert/jaune bicolore} (obligatoire et réservé) & Toutes les autres \\ \hline
\end{tabularx}
\end{center}
Le vert/jaune est \textbf{réservé exclusivement} à la terre : aucun autre conducteur n'a le droit de l'utiliser.
\end{aretenir}

\begin{attention}[Danger mortel : les confusions de couleurs]
\begin{itemize}
\item Ne jamais utiliser un fil \textbf{vert/jaune} comme phase ou comme neutre.
\item Ne jamais utiliser un fil \textbf{bleu} comme phase.
\end{itemize}
Pourquoi ? Parce qu'un jour, quelqu'un (toi, un électricien, un locataire) ouvrira cette boîte en se fiant aux couleurs. S'il touche un fil bleu en le croyant neutre alors que c'est une phase, il recevra \SI{230}{V}. Le code couleur n'est pas une décoration : c'est un \textbf{contrat de sécurité} entre tous ceux qui touchent à l'installation, aujourd'hui et dans trente ans.
\end{attention}

\courssub{La section des conducteurs : pourquoi l'épaisseur compte}

Fais l'expérience suivante : fais passer le même débit d'eau dans un tuyau d'arrosage fin puis dans un tuyau large. Dans le tuyau fin, l'eau « force », le tuyau chauffe et le jet est faible. Pour l'électricité, c'est pareil : un fil trop fin pour le courant qu'on lui impose \textbf{chauffe} (effet Joule), sa gaine fond, et c'est l'incendie. Les incendies d'origine électrique sont l'une des premières causes de feux d'habitation.

\begin{definition}[Section nominale]
La \cle{section nominale} d'un conducteur est la \textbf{surface de sa coupe transversale} (la partie métallique, sans l'isolant), exprimée en millimètres carrés ($\mathrm{mm}^2$). Elle détermine l'\textbf{intensité admissible} : c'est-à-dire le courant maximal que le fil peut transporter en permanence \textbf{sans échauffement excessif}. Plus la section est grande, plus le courant admissible est élevé.
\end{definition}

\begin{aretenir}[Correspondance Section / Calibre du disjoncteur / Usage]
\begin{center}
\begin{tabularx}{\linewidth}{|l|l|X|}\hline
\rowcolor{popblueL} \textbf{Section (cuivre)} & \textbf{Calibre max du disjoncteur ($I_n$)} & \textbf{Usage typique} \\ \hline
\SI{1,5}{mm^2} & \SI{10}{A} & Circuits d'éclairage (lampes, plafonniers) \\ \hline
\SI{2,5}{mm^2} & \SI{16}{A} & Circuits de prises de courant ordinaires \\ \hline
\SI{6}{mm^2} & \SI{32}{A} & Circuits spécialisés (cuisinière, plaque, chauffe-eau puissant) \\ \hline
\SI{10}{mm^2} & \SI{40}{A} à \SI{50}{A} & Départ général, alimentation du tableau \\ \hline
\end{tabularx}
\end{center}
Règle d'or : le \textbf{disjoncteur protège le fil}. Son calibre ne doit jamais dépasser le courant admissible de la section qu'il protège.
\end{aretenir}

\begin{methode}[Choisir la section d'un circuit]
\begin{enumerate}
\item \textbf{Identifier le calibre} du disjoncteur de protection prévu ($I_n$), en fonction de l'usage du circuit (éclairage, prises, circuit spécialisé).
\item \textbf{Choisir la section standard} correspondante dans le tableau (exemple : $I_n = \SI{16}{A}$ pour des prises $\Rightarrow$ $S = \SI{2,5}{mm^2}$).
\item \textbf{Vérifier la chute de tension} si la longueur du circuit dépasse \SI{20}{m} à \SI{30}{m} (grande maison, dépendance, portail électrique). Si la chute dépasse la limite, augmenter la section.
\end{enumerate}
\end{methode}

\begin{popfigure}
\begin{center}
\begin{tikzpicture}[scale=1.1]
% Gaine extérieure
\fill[popmuted!40] (0,0) circle (2.0);
\draw[very thick, popdark] (0,0) circle (2.0);
\fill[popcream] (0,0) circle (1.85);
% Conducteur Phase (rouge) en haut
\fill[poporange!85!black] (-0.62,0.55) circle (0.62);
\draw[thick, popdark] (-0.62,0.55) circle (0.62);
\fill[popgold!80!poporange] (-0.62,0.55) circle (0.38);
% Conducteur Neutre (bleu) en bas à droite
\fill[popblue] (0.62,0.55) circle (0.62);
\draw[thick, popdark] (0.62,0.55) circle (0.62);
\fill[popgold!80!poporange] (0.62,0.55) circle (0.38);
% Conducteur Terre (vert/jaune) en bas
\fill[popgreen] (0,-0.75) circle (0.62);
\draw[thick, popdark] (0,-0.75) circle (0.62);
\fill[popgold!80!poporange] (0,-0.75) circle (0.38);
% Annotations
\draw[->, thick, popdark] (-1.4,1.6) -- (-0.95,1.0);
\node[anchor=south, popdark, font=\small\bfseries] at (-1.4,1.6) {Phase (L) : rouge};
\draw[->, thick, popdark] (1.6,1.6) -- (1.0,1.0);
\node[anchor=south, popdark, font=\small\bfseries] at (1.6,1.6) {Neutre (N) : bleu};
\draw[->, thick, popdark] (1.5,-1.5) -- (0.55,-1.15);
\node[anchor=north, popdark, font=\small\bfseries, align=center] at (1.5,-1.5) {Terre (PE) :\\ vert/jaune};
\draw[->, thick, popdark] (-2.6,-1.4) -- (-1.9,-1.1);
\node[anchor=north, popdark, font=\small\bfseries, align=center] at (-2.6,-1.4) {Gaine\\ extérieure};
\draw[->, thick, popdark] (-2.6,0.9) -- (-1.15,0.6);
\node[anchor=east, popdark, font=\small\bfseries, align=center] at (-2.6,0.9) {Isolant\\ (PVC)};
\draw[->, thick, popdark] (2.7,-0.3) -- (0.95,-0.62);
\node[anchor=west, popdark, font=\small\bfseries, align=center] at (2.7,-0.3) {Âme en cuivre\\ $S = \SI{2,5}{mm^2}$};
\end{tikzpicture}
\end{center}
\legende{Coupe transversale d'un câble \textbf{3G2,5} : trois conducteurs (G = dont un conducteur de terre vert/jaune) de section \SI{2,5}{mm^2} chacun, chacun entouré de son isolant coloré, le tout protégé par une gaine extérieure.}
\end{popfigure}

\courssub{Gainage et repérage des conducteurs}

Un câble ne se laisse jamais nu dans un mur ou sous un toit. Il est toujours protégé mécaniquement.

\begin{definition}[Gainage]
Le \cle{gainage} est la protection mécanique des conducteurs, qui ajoute aussi une isolation supplémentaire. On utilise :
\begin{itemize}
\item les \textbf{tubes IRL} (rigides, posés en apparent, par exemple le long d'un mur d'atelier) ;
\item les \textbf{gaines ICTA} (souples et annelées, encastrées dans les murs avant le crépissage) ;
\item les \textbf{goulottes} (moulures à couvercle clipsable, en apparent, faciles à ouvrir pour ajouter un fil).
\end{itemize}
\end{definition}

Dans une boîte de dérivation ou au tableau, on trouve souvent plusieurs fils de même couleur (plusieurs phases rouges, par exemple). Comment savoir quel fil va vers la chambre et lequel va vers le salon ?

\begin{aretenir}[Repérage aux extrémités]
Chaque conducteur doit être \textbf{repéré à ses deux extrémités} par des \textbf{bagues colorées} ou des \textbf{étiquettes} portant les repères \textbf{L}, \textbf{N}, \textbf{PE} (et le nom du circuit : « prises salon », « éclairage chambre 1 »). Ce repérage est \textbf{obligatoire} dans les boîtes de dérivation et au tableau de répartition. Sans lui, le moindre dépannage devient long, hasardeux et dangereux.
\end{aretenir}

\courssub{La chute de tension : quand le fil est trop long}

Même un fil de bonne section possède une petite résistance. Sur une grande longueur (alimenter une dépendance au fond de la cour, un forage, un portail à \SI{50}{m} de la maison), cette résistance provoque une perte de tension : c'est la \cle{chute de tension}. Si elle est trop grande, les lampes éclairent faiblement, les moteurs peinent et chauffent.

\begin{propriete}[Formule simplifiée de la chute de tension]
Pour un circuit monophasé en cuivre, la chute de tension vaut :
\[ \Delta U = \frac{2 \times \rho \times L \times I}{S} \]
avec :
\begin{itemize}
\item $\rho = \SI{0,0175}{\Omega \cdot mm^2/m}$ : résistivité du cuivre ;
\item $L$ : longueur du câble en mètres (le facteur 2 tient compte de l'aller \emph{et} du retour du courant) ;
\item $I$ : intensité du courant en ampères ;
\item $S$ : section du conducteur en $\mathrm{mm}^2$.
\end{itemize}
Limites admises par la norme : $\Delta U \leq 3\,\%$ de $U_0$ pour l'\textbf{éclairage} et $\Delta U \leq 5\,\%$ de $U_0$ pour les \textbf{autres usages} (prises, moteurs), avec $U_0 = \SI{230}{V}$.
\end{propriete}

\begin{exemplebox}[Lecture de la formule]
La formule se lit naturellement : la chute de tension \textbf{augmente} avec la longueur $L$ et le courant $I$, et \textbf{diminue} quand la section $S$ augmente. Donc pour un circuit long ou très chargé, on prend un fil plus gros. C'est exactement le même raisonnement que pour le tuyau d'eau : plus le trajet est long, plus le tuyau doit être large.
\end{exemplebox}

\begin{exoresolu}[Vérification d'un circuit de prises]
\textbf{Énoncé.} Un circuit de prises en cuivre de section $S = \SI{2,5}{mm^2}$ alimente une charge de \SI{2000}{W} située à $L = \SI{25}{m}$ du tableau, sous $U_0 = \SI{230}{V}$. Calculer l'intensité $I$, puis la chute de tension $\Delta U$ en volts et en pourcentage. Conclure.

\tcblower
\textbf{Solution détaillée.}

\textbf{Étape 1 : calcul de l'intensité.}
\[ I = \frac{P}{U} = \frac{2000}{230} \approx \SI{8,7}{A} \]

\textbf{Étape 2 : application de la formule.}
\begin{align*}
\Delta U &= \frac{2 \times \rho \times L \times I}{S} \\
&= \frac{2 \times 0{,}0175 \times 25 \times 8{,}7}{2{,}5} \\
&\approx \SI{3,0}{V}
\end{align*}

\textbf{Étape 3 : conversion en pourcentage.}
\[ \frac{\Delta U}{U_0} = \frac{3{,}0}{230} \approx 0{,}013 = 1{,}3\,\% \]

\textbf{Conclusion.} $1{,}3\,\% < 5\,\%$ : la chute de tension est \textbf{conforme} pour un circuit de prises. La section de \SI{2,5}{mm^2} est validée.
\end{exoresolu}

\begin{popfigure}
\begin{center}
\begin{tikzpicture}
\begin{axis}[
width=0.95\linewidth, height=7.2cm,
xlabel={Longueur $L$ (m)},
ylabel={Chute de tension $\Delta U$ (\% de \SI{230}{V})},
xmin=0, xmax=60, ymin=0, ymax=8,
xtick={0,10,20,30,40,50,60},
ytick={0,1,2,3,4,5,6,7,8},
grid=both, grid style={popline!40},
legend style={at={(0.03,0.97)}, anchor=north west, font=\small, draw=popline, fill=popcream},
]
% Seuils
\addplot[domain=0:60, samples=2, dashed, thick, popgreen] {3};
\addlegendentry{Seuil 3\,\% (éclairage)}
\addplot[domain=0:60, samples=2, dashed, thick, poporange] {5};
\addlegendentry{Seuil 5\,\% (autres usages)}
% S=1.5 mm2, I=10 A : % = 2*0.0175*L*10/1.5/230*100 = 0.05072*L
\addplot[domain=0:60, samples=2, very thick, popblue] {0.0507*x};
\addlegendentry{$S=\SI{1,5}{mm^2}$, $I=\SI{10}{A}$}
% S=2.5 mm2, I=16 A : 2*0.0175*16/2.5/230*100 = 0.09739*L
\addplot[domain=0:60, samples=2, very thick, poppink] {0.0974*x};
\addlegendentry{$S=\SI{2,5}{mm^2}$, $I=\SI{16}{A}$}
% S=6 mm2, I=32 A : 2*0.0175*32/6/230*100 = 0.08116*L
\addplot[domain=0:60, samples=2, very thick, poppurple] {0.0812*x};
\addlegendentry{$S=\SI{6}{mm^2}$, $I=\SI{32}{A}$}
\end{axis}
\end{tikzpicture}
\end{center}
\legende{Chute de tension $\Delta U$ (en \% de \SI{230}{V}) en fonction de la longueur $L$ du circuit, pour trois couples section/courant. Lecture : le circuit prises \SI{2,5}{mm^2} à \SI{16}{A} (courbe rose) franchit le seuil de 5\,\% vers \SI{51}{m} ; au-delà, il faut passer à \SI{4}{mm^2} ou \SI{6}{mm^2}.}
\end{popfigure}

\begin{attention}[Pièges fréquents sur la chute de tension]
\begin{itemize}
\item \textbf{Oublier le facteur 2} : le courant fait l'aller-retour (phase + neutre), la longueur de fil est donc $2L$.
\item \textbf{Mélanger les unités} : $\rho$ est en $\Omega \cdot \mathrm{mm}^2/\mathrm{m}$, donc $L$ en mètres et $S$ en $\mathrm{mm}^2$ obligatoirement.
\item \textbf{Se tromper de seuil} : 3\,\% pour l'éclairage, 5\,\% pour les autres usages. Au BEPC, cite toujours le seuil utilisé pour justifier ta conclusion.
\item Croire qu'une chute de tension conforme dispense de vérifier le calibre du disjoncteur : les deux vérifications (échauffement \emph{et} chute de tension) sont indépendantes et toutes deux obligatoires.
\end{itemize}
\end{attention}

\begin{exercice}[Circuit d'éclairage d'une grande villa]
Un circuit d'éclairage en cuivre de section \SI{1,5}{mm^2}, protégé par un disjoncteur \SI{10}{A}, alimente des lampes situées à \SI{40}{m} du tableau. Le courant réel est de \SI{6}{A}.
\begin{enumerate}
\item Calculer la chute de tension $\Delta U$ en volts.
\item L'exprimer en pourcentage de \SI{230}{V}.
\item Le seuil pour l'éclairage est de 3\,\%. Le circuit est-il conforme ? Sinon, proposer une solution.
\end{enumerate}
\end{exercice}

\begin{corrige}[Exercice : circuit d'éclairage]
\begin{enumerate}
\item $\Delta U = \dfrac{2 \times 0{,}0175 \times 40 \times 6}{1{,}5} = \dfrac{8{,}4}{1{,}5} = \SI{5,6}{V}$.
\item $\dfrac{5{,}6}{230} \approx 0{,}0243 = 2{,}4\,\%$.
\item $2{,}4\,\% < 3\,\%$ : le circuit est \textbf{conforme}, mais de justesse. Si la longueur dépassait environ \SI{50}{m}, il faudrait passer à une section de \SI{2,5}{mm^2} (la chute tomberait alors à $\dfrac{2 \times 0{,}0175 \times 50 \times 6}{2{,}5} = \SI{4,2}{V}$, soit 1,8\,\%).
\end{enumerate}
\end{corrige}

\begin{savaistu}[Pourquoi pas de fils d'aluminium dans nos maisons ?]
L'aluminium conduit bien le courant et coûte moins cher que le cuivre : ENEO l'utilise d'ailleurs pour les lignes aériennes du réseau. Pourtant, la norme NF C 15-100 l'\textbf{interdit pour les circuits intérieurs fixes} des habitations, en France comme au Cameroun. Deux raisons : l'aluminium est \textbf{fragile au serrage} (il s'écrase et se « détend » sous les vis des bornes, ce qui crée des mauvais contacts qui chauffent) et il \textbf{s'oxyde} en surface, cette couche d'oxyde étant isolante et provoquant des échauffements dangereux. À section égale, il admet en outre moins de courant : il faudrait des sections environ 1,6 fois plus grandes. Le cuivre reste donc le roi de nos installations domestiques.
\end{savaistu}

\begin{aretenir}[L'essentiel de la section 2]
\begin{itemize}
\item Trois conducteurs : \textbf{Phase (L)} qui amène le courant, \textbf{Neutre (N)} qui le ramène, \textbf{Terre (PE)} qui protège les personnes.
\item Code couleur obligatoire : phase en rouge/brun/noir, neutre en \textbf{bleu clair}, terre en \textbf{vert/jaune} réservé.
\item La \textbf{section} (en $\mathrm{mm}^2$) fixe le courant admissible : \SI{1,5}{mm^2}/\SI{10}{A} (éclairage), \SI{2,5}{mm^2}/\SI{16}{A} (prises), \SI{6}{mm^2}/\SI{32}{A} (spécialisés), \SI{10}{mm^2}/\SI{40}{A}--\SI{50}{A} (départ général).
\item Les câbles sont protégés par tubes IRL, gaines ICTA ou goulottes, et repérés L, N, PE aux extrémités.
\item Chute de tension : $\Delta U = \dfrac{2 \rho L I}{S}$, limitée à 3\,\% (éclairage) ou 5\,\% (autres usages) de \SI{230}{V}.
\end{itemize}
\end{aretenir}

Maintenant que tu sais reconnaître chaque conducteur, choisir sa section et vérifier sa chute de tension, il est temps de voir comment ces trois fils aboutissent concrètement sur une prise de courant : c'est l'objet de la section suivante, consacrée à la prise 2P+T et au rôle précis de ses bornes.

\coursec{La prise de courant 2P+T : bornes, branchement et rôle de la terre}

Dans la section précédente, nous avons appris à identifier les conducteurs d'une installation grâce au code couleur : bleu pour le neutre, rouge ou noir pour la phase, vert-jaune pour la terre. Mais à quoi servent exactement ces trois fils une fois qu'ils arrivent au mur de ta chambre ? Ils aboutissent à un organe que tu utilises chaque jour sans y penser : la \cle{prise de courant}. C'est elle qui fait le lien entre l'installation fixe de la maison et tes appareils mobiles (télévision, fer à repasser, chargeur de téléphone, groupe électrogène). Dans cette section, nous allons ouvrir une prise, comprendre le rôle de chacune de ses bornes, découvrir pourquoi la fameuse broche de terre peut te sauver la vie, et apprendre à brancher une prise dans les règles de l'art.

\courssub{Description de la prise 2P+T (norme NF C 61-314, type E)}

Observe la prise murale de ton salon : tu vois deux trous cylindriques et, entre eux, une petite tige métallique qui dépasse. Cette prise est dite \cle{2P+T}, c'est-à-dire « deux pôles plus terre ». C'est le modèle normalisé en France et au Cameroun (norme NF C 61-314, aussi appelée « type E »).

\begin{definition}[Prise de courant 2P+T]
Une \cle{prise de courant 2P+T} est un connecteur mural comportant trois contacts :
\begin{itemize}
\item le pôle \cle{Phase} (repère \textbf{L}), qui amène le courant depuis le réseau ;
\item le pôle \cle{Neutre} (repère \textbf{N}), qui assure le retour du courant ;
\item le contact de \cle{Terre} (repère \textbf{PE} ou symbole $\downarrow$), réalisé sous forme d'une \textbf{broche mâle} solidaire de la prise, qui s'insère dans la fiche de l'appareil.
\end{itemize}
Elle délivre une tension efficace de \SI{220}{V} à \SI{240}{V}, à la fréquence de \SI{50}{Hz}.
\end{definition}

\begin{attention}[Ne pas confondre prise et fiche]
La \textbf{prise} est la partie \emph{femelle}, fixée au mur. La \textbf{fiche} est la partie \emph{mâle}, au bout du cordon de l'appareil. Sur la prise type E, seule la terre est une broche mâle : c'est une particularité que nous expliquerons plus loin.
\end{attention}

En retournant une prise démontée (installation hors tension !), on découvre à l'arrière trois bornes de raccordement, clairement repérées :

\begin{center}
\begin{tabularx}{\linewidth}{|l|X|X|}\hline
\rowcolor{popblueL} \textbf{Borne} & \textbf{Position (vue de face)} & \textbf{Repérage constructeur} \\ \hline
Phase (L) & Trou de \textbf{droite} & Lettre L, souvent marque rouge \\ \hline
Neutre (N) & Trou de \textbf{gauche} & Lettre N, souvent marque bleue \\ \hline
Terre (PE) & Broche centrale, en \textbf{haut} & Marque verte/jaune ou symbole de terre \\ \hline
\end{tabularx}
\end{center}

\begin{aretenir}[Le sens de montage : la terre en haut]
La convention impose de poser la prise \textbf{broche de terre vers le haut}. Pourquoi ? Si la fiche est à moitié débranchée et qu'un objet métallique (règle, couteau, lame de scie) tombe par-dessus, il touchera d'abord la broche de terre — qui n'est jamais dangereuse — avant de pouvoir toucher les pôles actifs. C'est un détail de conception qui évite des courts-circuits et des électrisations.
\end{aretenir}

\begin{popfigure}
\begin{center}
\begin{tikzpicture}[scale=1.0, every node/.style={font=\small}]
% ===== VUE DE FACE =====
\begin{scope}
\draw[fill=popcream, thick, popdark] (0,0) circle (1.7);
\draw[fill=white, thick, popdark] (0,0) circle (1.45);
% broche terre en haut
\draw[fill=popgreen, thick, popdark] (0,1.15) circle (0.16);
% trous L et N
\draw[fill=popink!80, thick] (0.55,-0.25) circle (0.14);
\draw[fill=popink!80, thick] (-0.55,-0.25) circle (0.14);
% ergots de fixation (petits trous ronds)
\draw[fill=popmuted!40, popdark] (0.55,-1.05) circle (0.07);
\draw[fill=popmuted!40, popdark] (-0.55,-1.05) circle (0.07);
% annotations
\draw[->, popgreen, thick] (0,1.35) -- (0,2.1) node[above, popgreen!60!black]{\textbf{Terre (PE) en haut}};
\draw[->, poporange, thick] (0.75,-0.25) -- (1.9,-0.25) node[right, poporange!70!black]{\textbf{L : Phase (droite)}};
\draw[->, popblue, thick] (-0.75,-0.25) -- (-1.9,-0.25) node[left, popblue!70!black]{\textbf{N : Neutre (gauche)}};
\node[below, popdark] at (0,-1.9) {\textbf{Vue de face}};
\end{scope}
% ===== VUE ARRIERE =====
\begin{scope}[xshift=8.2cm]
\draw[fill=popcream, thick, popdark, rounded corners=6pt] (-2,-1.7) rectangle (2,1.7);
% bornes de raccordement (rectangles avec vis)
\draw[fill=popblue!30, thick, popblue] (-1.2,-0.5) rectangle (-0.5,0.3);
\draw[fill=poporange!30, thick, poporange] (0.5,-0.5) rectangle (1.2,0.3);
\draw[fill=popgreen!30, thick, popgreen!60!black] (-0.35,0.9) rectangle (0.35,1.5);
% vis de serrage
\draw[popdark] (-0.85,-0.1) circle (0.09);
\draw[popdark] (0.85,-0.1) circle (0.09);
\draw[popdark] (0,1.2) circle (0.09);
% fils sortant vers le bas
\draw[line width=2.2pt, popblue] (-0.85,-0.5) -- (-0.85,-1.7);
\draw[line width=2.2pt, poporange] (0.85,-0.5) -- (0.85,-1.7);
\draw[line width=2.2pt, popgreen!70!popgold] (0,0.9) -- (0,0.3) -- (0,-1.7);
% serre-câble
\draw[thick, popdark, fill=popmuted!25] (-0.6,-1.7) rectangle (0.6,-1.35);
\node[popblue!70!black, below] at (-0.85,-1.75) {\textbf{Bleu $\to$ N}};
\node[poporange!70!black, below] at (0.85,-1.75) {\textbf{Rouge/Noir $\to$ L}};
\node[popgreen!50!black, right] at (0.4,1.2) {\textbf{Vert-Jaune $\to$ Terre}};
\node[popdark, left] at (-0.65,-1.5) {Serre-câble};
\node[below, popdark] at (0,-2.35) {\textbf{Vue arrière (raccordement)}};
\end{scope}
\end{tikzpicture}
\legende{Prise 2P+T type E. À gauche : vue de face, la broche de terre est placée en haut, la phase (L) à droite et le neutre (N) à gauche. À droite : vue arrière avec les trois bornes de raccordement (vis apparentes), le code couleur des conducteurs et le serre-câble qui immobilise la gaine.}
\end{center}
\end{popfigure}

\courssub{Le rôle de la terre : un fil qui protège les personnes}

\courssub{La notion de masse}

Pour comprendre la terre, il faut d'abord définir la \cle{masse} d'un appareil.

\begin{definition}[Masse d'un appareil électrique]
La \cle{masse} est toute partie conductrice d'un appareil \textbf{accessible au toucher} (carcasse métallique d'un lave-linge, d'un four, d'un réfrigérateur, châssis d'un groupe électrogène), \textbf{normalement hors tension}, mais susceptible de se trouver accidentellement sous tension en cas de \textbf{défaut d'isolement} (par exemple si un fil de phase dénudé touche la carcasse).
\end{definition}

Imagine la scène : à l'intérieur de ton lave-linge, l'isolant d'un fil de phase s'est usé par vibration et le fil nu touche la carcasse métallique. Sans protection, toute la carcasse se retrouve portée à \SI{220}{V}. Si tu la touches, le courant traverse ton corps vers le sol : c'est l'\cle{électrisation}, potentiellement mortelle.

\courssub{Le conducteur de protection (PE)}

Le fil vert-jaune, appelé \cle{conducteur de protection} (PE), relie la carcasse de l'appareil (via la fiche et la prise) à un \textbf{piquet de terre} planté dans le sol de la cour ou du jardin. La masse est ainsi maintenue au potentiel de la terre, c'est-à-dire \SI{0}{V}.

En cas de défaut phase-masse, un \textbf{courant de défaut} $I_d$ s'établit : phase $\to$ carcasse $\to$ fil de terre $\to$ piquet $\to$ sol $\to$ neutre du transformateur ENEO. Le \textbf{disjoncteur différentiel} (\SI{30}{mA}), placé en tête de l'installation, compare en permanence le courant qui part par la phase et celui qui revient par le neutre. Dès qu'il détecte une différence supérieure à \SI{30}{mA} — signe qu'une partie du courant « fuit » vers la terre — il \textbf{coupe automatiquement} le circuit en une fraction de seconde (moins de \SI{40}{ms}), avant que le courant ne devienne dangereux pour le corps humain.

\begin{popfigure}
\begin{center}
\begin{tikzpicture}[scale=0.95, every node/.style={font=\small}]
% Disjoncteur différentiel
\draw[thick, popdark, fill=popblueL] (0,3.2) rectangle (1.4,4.2);
\node[popdark] at (0.7,3.7) {\textbf{Diff.}};
\node[popdark, below] at (0.7,3.15) {\SI{30}{mA}};
% Arrivée réseau (Phase + Neutre)
\draw[line width=1.6pt, poporange] (-2.2,3.95) -- (0,3.95);
\draw[line width=1.6pt, popblue] (-2.2,3.45) -- (0,3.45);
\node[left, poporange!70!black] at (-2.2,3.95) {Phase};
\node[left, popblue!70!black] at (-2.2,3.45) {Neutre};
% Lignes vers appareil
\draw[line width=1.6pt, poporange] (1.4,3.95) -- (5.2,3.95);
\draw[line width=1.6pt, popblue] (1.4,3.45) -- (5.2,3.45);
% Appareil Classe I (Lave-linge)
\draw[thick, popdark, fill=popmuted!20, rounded corners=4pt] (5.2,2.6) rectangle (7.4,4.6);
\node[popdark] at (6.3,4.25) {Lave-linge};
\node[popdark] at (6.3,2.85) {(Classe I)};
% Défaut Phase -> Masse
\draw[line width=1.6pt, poporange] (5.6,3.95) -- (5.6,3.55);
\fill[poppink] (5.6,3.55) circle (0.09);
\node[poppink!70!black, left] at (5.5,3.7) {Défaut};
% Fil de terre (PE) vers le bas
\draw[line width=1.8pt, popgreen!70!popgold] (6.3,2.6) -- (6.3,0.6);
\draw[->, line width=1.4pt, popgreen!50!black] (6.3,1.6) -- (6.3,0.7);
\node[popgreen!50!black, right] at (6.4,1.6) {$I_d$ : courant de défaut};
% Sol (terre)
\draw[thick, popdark] (4.6,0.6) -- (8.0,0.6);
\foreach \x in {4.8,5.2,5.6,6.0,6.4,6.8,7.2,7.6}{\draw[popdark] (\x,0.6) -- (\x-0.18,0.38);}
% Piquet de terre
\draw[line width=2.2pt, popgold!80!black] (6.3,0.6) -- (6.3,-0.3);
\draw[popgold!80!black] (6.15,-0.3) -- (6.45,-0.3);
\draw[popgold!80!black] (6.2,-0.42) -- (6.4,-0.42);
\draw[popgold!80!black] (6.25,-0.54) -- (6.35,-0.54);
\node[popgold!60!black, right] at (6.5,-0.35) {Piquet de terre};
% Retour vers neutre transfo par le sol (flèche pointillée)
\draw[->, dashed, thick, poppurple] (6.0,-0.7) -- (-1.0,-0.7) -- (-1.0,2.6) node[midway, left, poppurple!70!black, align=center]{Retour par le sol\\vers le neutre ENEO};
% Déclenchement différentiel
\draw[->, dashed, thick, poppink] (0.7,4.2) -- (0.7,4.9) node[above, poppink!70!black]{Déclenchement !};
\end{tikzpicture}
\legende{Schéma de principe du défaut d'isolement en régime TT. Le fil de phase touche la carcasse du lave-linge : le courant de défaut $I_d$ s'écoule par le conducteur de protection (vert-jaune) vers le piquet de terre, puis revient par le sol au neutre du transformateur. Le différentiel \SI{30}{mA} détecte cette fuite ($I_{\text{phase}} \neq I_{\text{neutre}}$) et coupe le circuit instantanément.}
\end{center}
\end{popfigure}

\begin{propriete}[Le régime TT et l'obligation du différentiel]
Au Cameroun comme en France, les installations domestiques fonctionnent en \cle{régime TT} : la prise de terre de l'installation (le piquet planté chez toi) est \textbf{distincte} de la prise de terre du neutre du transformateur ENEO. Dans ce régime, la boucle de défaut passe par le sol, dont la résistance n'est pas négligeable : le courant de défaut est donc trop faible pour faire sauter un simple fusible ou déclencher un disjoncteur magnétothermique classique. C'est pourquoi le \textbf{disjoncteur différentiel \SI{30}{mA} est OBLIGATOIRE} : c'est lui, et lui seul, qui protège les personnes contre les contacts indirects.
\end{propriete}

\begin{exemplebox}[Pourquoi mesurer la résistance de terre ?]
La tension de contact $U_c$ à laquelle est portée la carcasse au moment du défaut (au seuil de déclenchement du différentiel) vaut $U_c = R_{\text{terre}} \times I_{\Delta n}$ avec $I_{\Delta n} = \SI{30}{mA}$. La norme impose $R_{\text{terre}} \leq \SI{100}{\ohm}$ pour un différentiel \SI{30}{mA}.
\begin{itemize}
\item Si $R_{\text{terre}} = \SI{50}{\ohm}$ : $U_c = \num{0,03} \times 50 = \SI{1,5}{V}$ : sans danger, et le différentiel coupe aussitôt.
\item Si $R_{\text{terre}} = \SI{500}{\ohm}$ (piquet de mauvaise qualité, sol sec en saison sèche) : $U_c = \num{0,03} \times 500 = \SI{15}{V}$, et surtout le courant de défaut réel peut rester inférieur à \SI{30}{mA} : le différentiel ne déclenche pas, la carcasse reste sous tension \textbf{de façon permanente} : danger mortel !
\end{itemize}
D'où l'importance de faire \textbf{mesurer la résistance de terre} de son installation par un électricien qualifié lors de la réception des travaux et périodiquement.
\end{exemplebox}

\courssub{Appareils de classe I et de classe II}

Tous les appareils n'ont pas besoin de la terre. On distingue deux grandes classes de protection contre les chocs électriques :

\begin{center}
\begin{tabularx}{\linewidth}{|l|X|X|}\hline
\rowcolor{popblueL} \textbf{Critère} & \textbf{Classe I} & \textbf{Classe II} \\ \hline
Protection contre les contacts indirects & Carcasse métallique reliée à la terre (fil vert-jaune) & Double isolation renforcée (pas de pièce métallique accessible reliée à l'intérieur) \\ \hline
Fiche & Avec contact de terre (3 broches) & Sans terre (2 broches seulement) \\ \hline
Symbole normalisé & Symbole de terre : $\downarrow$ & Double carré : $\square$ dans $\square$ \\ \hline
Exemples courants & Lave-linge, four, fer à repasser, réfrigérateur, unité centrale PC, pompe à eau, groupe électrogène & Téléviseur, chargeur de téléphone, perceuse sans fil, sèche-cheveux, lampe sur pied \\ \hline
\end{tabularx}
\end{center}

\begin{attention}[Deux erreurs fatales à ne jamais commettre]
\begin{enumerate}
\item \textbf{Inverser phase et neutre} sur une prise : l'appareil fonctionnera normalement (courant alternatif), donc personne ne s'en aperçoit... mais l'interrupteur, qui est câblé sur la phase, ne coupera plus que le neutre. La lampe ou l'appareil restera sous tension même « éteint » : danger mortel lors du changement d'une ampoule ou d'une réparation.
\item \textbf{Utiliser le fil de terre comme neutre} (bricolage fréquent quand il manque un fil bleu) : la terre n'est pas prévue pour véhiculer le courant de service. Toutes les carcasses reliées à cette « terre » se retrouvent alors sous tension. La terre ne doit \textbf{JAMAIS} servir de neutre !
\end{enumerate}
\end{attention}

\courssub{Branchement pratique d'une prise murale}

\begin{methode}[Raccorder une prise 2P+T en sept étapes]
\begin{enumerate}
\item \textbf{Couper} le disjoncteur divisionnaire du circuit concerné (ou le disjoncteur général) et \textbf{vérifier l'absence de tension (VAT)} avec le tournevis testeur sur les fils à raccorder.
\item \textbf{Dénuder} chaque fil sur \num{10} à \num{12}{mm} avec la pince à dénuder, sans entailler les brins de cuivre.
\item \textbf{Torsader} légèrement les brins de chaque extrémité pour éviter qu'un brin ne s'échappe et ne crée un court-circuit.
\item \textbf{Connecter} rigoureusement : fil \textbf{bleu} sur la borne \textbf{N} (gauche), fil \textbf{rouge ou noir} sur la borne \textbf{L} (droite), fil \textbf{vert-jaune} sur la borne de \textbf{terre} (centre, en haut).
\item \textbf{Serrer} fermement les vis (couple conseillé : \num{0,8} à \SI{1,2}{N.m}), puis tirer doucement sur chaque fil pour vérifier qu'il est bien prisonnier. Aucun fil nu ne doit dépasser de la borne.
\item \textbf{Fixer} le serre-câble sur la gaine (pas sur les fils nus), clipser le mécanisme dans le boîtier d'encastrement et poser la plaque de finition, \textbf{broche de terre en haut}.
\item \textbf{Remettre} sous tension et \textbf{tester} avec le tournevis testeur : il doit s'allumer \textbf{uniquement} sur le trou de droite (phase).
\end{enumerate}
\end{methode}

\begin{experience}[Vérification de la polarité d'une prise avec un tournevis testeur]
\textbf{Matériel :} Une prise montée (ou le circuit de la section 3.3), un tournevis testeur (type néon), un multimètre (optionnel).
\textbf{Protocole :}
\begin{enumerate}
\item Mettre la prise sous tension (disjoncteur enclenché).
\item Insérer la pointe du tournevis testeur dans le trou de \textbf{droite} (phase) : le néon doit s'allumer (faiblement).
\item Insérer la pointe dans le trou de \textbf{gauche} (neutre) : le néon doit rester \textbf{éteint}.
\item Toucher la broche de terre (haut) : le néon doit rester \textbf{éteint}.
\end{enumerate}
\textbf{Observations :} Si le néon s'allume à gauche ou sur la terre, il y a une inversion de fils ou un défaut d'isolement. Couper immédiatement le disjoncteur et revérifier le câblage.
\textbf{Interprétation :} Le tournevis testeur met en évidence la présence de la tension par rapport à la terre (ton corps fait office de capacité vers la terre). Il ne mesure pas la valeur de la tension, il indique seulement la présence de la phase.
\end{experience}

\begin{exoresolu}[Identifier et corriger un mauvais raccordement]
Énoncé. Un élève bricoleur raccorde une prise ainsi : fil bleu sur la borne de droite, fil rouge sur la borne de gauche, fil vert-jaune sur la terre. Il teste : le tournevis testeur s'allume sur le trou de gauche. Son père branche une lampe : elle s'allume normalement. Le père affirme : « Tout est correct puisque ça marche. » Qu'en penses-tu ? Justifie et propose une correction.
\tcblower
Solution. Le raccordement est \textbf{incorrect} : la phase (fil rouge) est sur la borne de gauche et le neutre (fil bleu) sur celle de droite, c'est l'inverse de la convention NF C 61-314. Le testeur le confirme : il s'allume à gauche, alors qu'il devrait s'allumer à droite. Certes, la lampe fonctionne, car une lampe est un récepteur passif qui fonctionne quel que soit le sens du courant alternatif. Mais le danger est caché : tout interrupteur ou tout appareil câblé en amont en respectant la convention coupera le neutre au lieu de la phase. Lors d'un changement d'ampoule ou d'une réparation, on pourra toucher une partie encore portée à \SI{220}{V} en croyant le circuit hors tension. Correction : couper le disjoncteur, vérifier l'absence de tension (VAT), puis remettre le fil rouge sur la borne L (droite) et le fil bleu sur la borne N (gauche), resserrer les vis et refaire le test au tournevis testeur.
\end{exoresolu}

\begin{savaistu}[La broche de terre : une spécialité franco-camerounaise]
La broche de terre mâle qui dépasse de la prise est une particularité du système français (type E), hérité au Cameroun. En Allemagne et dans de nombreux pays européens, on utilise le type F dit « Schuko » : la terre est assurée par deux languettes métalliques latérales, sur les côtés de la fiche. Il existe aujourd'hui des fiches hybrides (type E/F), compatibles avec les deux systèmes : elles possèdent à la fois un trou pour la broche française et des languettes latérales pour le système allemand. Regarde la fiche de ton ordinateur portable ou de ton chargeur de téléphone : il y a de fortes chances qu'elle soit de ce type hybride !
\end{savaistu}

\begin{aretenir}[L'essentiel de la section]
\begin{itemize}
\item Une prise 2P+T comporte trois contacts : phase (L, à droite), neutre (N, à gauche) et terre (broche mâle, en haut).
\item Le conducteur de protection vert-jaune relie les masses métalliques des appareils de classe I au piquet de terre ; en cas de défaut, le différentiel \SI{30}{mA} coupe le circuit.
\item En régime TT (Cameroun), le différentiel \SI{30}{mA} est obligatoire et la résistance de terre doit être inférieure à \SI{100}{\ohm}.
\item On n'inverse jamais phase et neutre, et la terre ne sert jamais de neutre.
\item Le branchement se fait hors tension, en respectant le code couleur, et se vérifie au tournevis testeur (phase à droite uniquement).
\end{itemize}
\end{aretenir}

\begin{exercice}[Entraînement : La prise de courant]
\begin{enumerate}
\item Sur une prise 2P+T type E vue de face, indique la position (gauche, droite, haut, bas) de la phase, du neutre et de la terre.
\item Un appareil de classe II (ex: chargeur téléphone) a une fiche à deux broches. Pourquoi n'a-t-il pas besoin de la broche de terre ?
\item Lors du branchement d'une prise, on dénude les fils sur \num{10} à \num{12}{mm}. Que risque-t-on si on les dénude sur \num{20}{mm} ?
\item Un électricien mesure une résistance de terre de \SI{150}{\ohm} sur une installation protégée par un différentiel \SI{30}{mA}. Cette installation est-elle conforme ? Justifie par un calcul de la tension de contact maximale au seuil de déclenchement.
\item Complète : En régime TT, la boucle de défaut passe par le \_\_\_\_\_\_\_\_\_. C'est pour cela que le \_\_\_\_\_\_\_\_\_ \SI{30}{mA} est indispensable.
\end{enumerate}
\end{exercice}

\begin{corrige}[Exercice : La prise de courant]
\begin{enumerate}
\item Phase : trou de \textbf{droite}. Neutre : trou de \textbf{gauche}. Terre : \textbf{broche en haut}.
\item L'appareil de classe II possède une \textbf{double isolation renforcée}. Aucune pièce métallique accessible n'est reliée aux parties actives internes. Il n'y a donc pas de masse à mettre à la terre.
\item Si on dénude trop long (\num{20}{mm}), du cuivre nu dépassera de la borne de raccordement. Risque de \textbf{court-circuit} entre phase et neutre (ou phase et terre) à l'intérieur de la prise, ou risque de contact direct avec un fil nu sous tension.
\item $U_c = R_{\text{terre}} \times I_{\Delta n} = 150 \times \num{0,03} = \SI{4,5}{V}$. Bien que \SI{4,5}{V} soit une tension de contact faible (non dangereuse), la norme impose $R_{\text{terre}} \leq \SI{100}{\ohm}$ pour garantir le déclenchement du différentiel même pour des courants de défaut inférieurs au seuil nominal. Avec \SI{150}{\ohm}, le courant de défaut réel pourrait être $< \SI{30}{mA}$ et le différentiel ne pas déclencher. \textbf{L'installation n'est pas conforme.}
\item La boucle de défaut passe par le \textbf{sol}. C'est pour cela que le \textbf{disjoncteur différentiel} \SI{30}{mA} est indispensable.
\end{enumerate}
\end{corrige}

\coursec{Le tournevis testeur : principe, utilisation et limites}

Dans la section précédente, nous avons appris à reconnaître les trois bornes d'une prise 2P+T : la \cle{phase}, le \cle{neutre} et la \cle{terre}. Mais voici un problème très concret : chez toi, face à une prise murale, comment savoir \emph{lequel} des deux orifices est la phase ? Le courant ne se voit pas, ne se sent pas, ne sent rien. Il faut un instrument. L'outil le plus répandu dans les ateliers et les maisons camerounaises est le fameux \og tournevis testeur \fg{}, que les électriciens appellent aussi \og cherche-phase \fg{}. Dans cette section, nous allons comprendre \emph{comment} il fonctionne, \emph{comment} l'utiliser correctement, et surtout \emph{ce qu'il ne faut jamais lui demander} : c'est une question de vie ou de mort.

\courssub{Qu'est-ce qu'un tournevis testeur ?}

\begin{definition}[Tournevis testeur]
Le \cle{tournevis testeur} (ou \og cherche-phase \fg{}) est un petit outil de la taille d'un tournevis, qui permet de \textbf{détecter la présence d'une tension alternative} (généralement entre \SI{100}{V} et \SI{500}{V}) sur un conducteur. Il contient une petite lampe au \cle{néon} (ou une LED) montée en série avec une \textbf{très forte résistance} (environ \SI{1}{M\Omega}), le tout logé dans un manche transparent. Le circuit se referme à travers le \textbf{corps de l'utilisateur} et sa liaison capacitive avec la terre.
\end{definition}

Ouvrons mentalement l'outil. De la pointe vers le manche, on trouve dans l'ordre :
\begin{enumerate}
\item la \textbf{pointe métallique} (la lame du tournevis), que l'on met en contact avec le conducteur à tester ;
\item la \textbf{lampe néon}, petit tube de verre contenant du gaz néon qui s'allume en orange-rouge dès qu'une tension d'environ \SI{60}{V} à \SI{80}{V} apparaît à ses bornes ;
\item la \textbf{résistance de sécurité} d'environ \SI{1}{M\Omega} (un million d'ohms !), cachée dans le manche ;
\item le \textbf{bouton métallique} (ou anneau) à l'extrémité du manche, sur lequel on pose le doigt.
\end{enumerate}

\begin{experience}[Observation : le testeur ne s'allume que si on touche le bouton]
\textbf{Matériel :} un tournevis testeur, une prise sous tension, un adulte encadrant.\par
\textbf{Protocole :} (1) Insérer la pointe du testeur dans l'orifice phase d'une prise \emph{sans toucher} le bouton métallique du manche. (2) Recommencer en posant fermement le pouce sur le bouton métallique.\par
\textbf{Observations :} dans le cas (1), le néon reste éteint ; dans le cas (2), le néon s'allume faiblement.\par
\textbf{Interprétation :} le néon ne peut s'allumer que si un courant, même minuscule, le traverse. Ce courant doit donc passer \emph{par le corps de l'utilisateur}. Sans contact du doigt, le circuit est ouvert : rien ne s'allume.\par
\textbf{Conclusion :} le corps humain fait partie du circuit électrique du tournevis testeur.
\end{experience}

\courssub{Principe de fonctionnement : le courant passe par ton corps !}

Voici l'idée qui surprend toujours : quand le néon s'allume, \textbf{un courant électrique traverse ton corps}. Rassure-toi tout de suite : il est si faible que tu ne le sens même pas. Suivons son trajet complet :

\begin{center}
Pointe $\rightarrow$ Néon $\rightarrow$ Résistance \SI{1}{M\Omega} $\rightarrow$ Doigt $\rightarrow$ Corps $\rightarrow$ Capacité corps/terre $\rightarrow$ Terre $\rightarrow$ Neutre du transformateur.
\end{center}

Deux remarques essentielles :
\begin{itemize}
\item Le corps humain n'est pas relié à la terre par un fil, mais il forme avec le sol un petit \textbf{condensateur} : on parle de \cle{capacité parasitaire} (ou capacitive) corps-terre. En courant alternatif à \SI{50}{Hz}, cette capacité laisse passer un courant minuscule, suffisant pour le néon.
\item La boucle se referme au niveau du \textbf{transformateur ENEO} du quartier, où le neutre est relié à la terre (liaison neutre-terre). C'est exactement le même principe que la prise de terre étudiée à la section 3.
\end{itemize}

\begin{popfigure}
\begin{center}
\begin{tikzpicture}[font=\small, >=Stealth]
% Pointe
\draw[fill=popmuted!40] (0,-0.12) rectangle (1.6,0.12);
\draw[fill=popmuted!60] (1.6,-0.25) rectangle (2.1,0.25);
\node[below] at (0.8,-0.15) {Pointe};
% Néon
\draw[fill=poporangeL] (2.1,-0.3) rectangle (3.1,0.3);
\node at (2.6,0) {Néon};
% Résistance
\draw[fill=poppinkL] (3.1,-0.3) rectangle (5.1,0.3);
\node at (4.1,0) {$R \approx \SI{1}{M\Omega}$};
% Bouton + doigt
\draw[fill=popmuted!60] (5.1,-0.35) rectangle (5.5,0.35);
\draw[fill=popgoldL] (5.5,-0.5) rectangle (6.3,0.5);
\node at (5.9,0) {Doigt};
% Corps
\draw[fill=popgold!30, rounded corners] (6.3,-0.9) rectangle (7.5,0.9);
\node at (6.9,0) {Corps};
% Capacité corps-terre
\draw[thick] (7.5,0) -- (8.1,0);
\draw[thick] (8.1,-0.4) -- (8.1,0.4);
\draw[thick] (8.3,-0.4) -- (8.3,0.4);
\node[above] at (8.2,0.45) {$C_{\text{corps/terre}}$};
\draw[thick] (8.3,0) -- (8.9,0);
% Terre
\draw[thick] (8.9,0) -- (8.9,-0.5);
\draw[thick] (8.6,-0.5) -- (9.2,-0.5);
\draw[thick] (8.7,-0.65) -- (9.1,-0.65);
\draw[thick] (8.8,-0.8) -- (9.0,-0.8);
\node[right] at (9.2,-0.6) {Terre};
% Retour neutre
\draw[thick, popblue] (8.9,-0.9) -- (8.9,-1.6) -- (0.8,-1.6) -- (0.8,-0.3);
\node[popblue, below] at (4.8,-1.6) {Retour vers le neutre du transformateur (liaison neutre-terre)};
% Flèche courant
\draw[->, very thick, poppink] (0.3,0.7) -- (5.3,0.7) node[midway, above] {Courant $I \approx \num{0,2}$ mA};
% Phase
\node[popink, above] at (0.2,0.15) {Phase \SI{230}{V}};
\end{tikzpicture}
\end{center}
\legende{Schéma équivalent du circuit électrique d'un tournevis testeur : le courant part de la phase, traverse le néon, la résistance de \SI{1}{M\Omega}, le doigt et le corps de l'utilisateur, puis rejoint la terre par la capacité corps-terre, avant de revenir au neutre du transformateur.}
\end{popfigure}

\begin{propriete}[Sécurité intrinsèque du testeur]
Grâce à la résistance série d'environ \SI{1}{M\Omega}, le courant qui traverse l'utilisateur est limité à \textbf{moins de \SI{1}{mA}}, c'est-à-dire en dessous du \cle{seuil de perception} (environ \SI{1}{mA}) et très loin du seuil de non-lâcher (\SI{10}{mA}). En conditions normales d'utilisation sur le réseau \SI{230}{V}, le tournevis testeur est donc \textbf{sans danger pour un utilisateur sain}. Le néon ne s'allume que si la tension dépasse environ \SI{100}{V}.
\end{propriete}

\begin{exoresolu}[Calcul du courant qui traverse l'utilisateur]
\textbf{Énoncé.} Un élève utilise un tournevis testeur sur la phase d'une prise \SI{230}{V}. La résistance interne du testeur vaut $R = \SI{1}{M\Omega}$. La résistance du corps et l'impédance de la capacité corps-terre valent ensemble environ $R_c = \SI{500}{k\Omega}$. (1) Calculer le courant $I$ traversant le corps. (2) Comparer au seuil de non-lâcher de \SI{10}{mA} et conclure.
\tcblower
\textbf{Solution.} (1) Les résistances sont en série, donc la résistance totale vaut :
\[ R_{\text{tot}} = R + R_c = \num{1000000} + \num{500000} = \num{1500000}\ \Omega = \SI{1,5}{M\Omega}. \]
D'après la loi d'Ohm :
\[ I = \frac{U}{R_{\text{tot}}} = \frac{230}{\num{1500000}} \approx \num{1,5e-4}\ \text{A} = \SI{0,15}{mA}. \]
(2) Comparaison : $\num{0,15}\ \text{mA} \ll 10\ \text{mA}$. Le courant est environ \textbf{65 fois plus faible} que le seuil de non-lâcher, et même inférieur au seuil de perception (\SI{1}{mA}) : on ne sent rien. \textbf{Conclusion :} en conditions normales, l'utilisation du tournevis testeur sur \SI{230}{V} est sans danger.
\end{exoresolu}

\courssub{Utilisation correcte : identifier la phase}

\begin{methode}[Identifier la phase sur une prise inconnue]
\begin{enumerate}
\item \textbf{Vérifier le testeur} sur une prise dont on connaît la phase (ou une prise sûre, sous tension) : le néon doit s'allumer. Si ce n'est pas le cas, le testeur est hors service : ne pas l'utiliser.
\item \textbf{Tenir le manche} en posant fermement le pouce ou l'index sur le \textbf{bouton métallique} (ou l'anneau) : c'est ce contact qui referme le circuit à travers le corps.
\item \textbf{Insérer la pointe} dans chacun des deux orifices de la prise à tester, l'un après l'autre.
\item L'orifice qui \textbf{allume le néon} est la \cle{phase}. Celui qui laisse le néon éteint est le \cle{neutre}.
\item La \textbf{terre} se reconnaît sans test : c'est la broche métallique saillante (ou les clips latéraux) de la prise 2P+T.
\end{enumerate}
\end{methode}

\begin{attention}[Gestes interdits]
\begin{itemize}
\item Ne jamais toucher la \textbf{lame métallique} du testeur avec les doigts pendant le test : le courant ne serait alors plus limité par la résistance interne !
\item Ne jamais utiliser le testeur comme un vrai tournevis pour forcer sur des vis : le manche transparent est fragile et contient le néon.
\item Ne jamais tester une installation humide ou avec les mains mouillées.
\end{itemize}
\end{attention}

\courssub{Les limites majeures du tournevis testeur}

C'est le point le plus important de cette leçon. Le tournevis testeur est un bon \emph{détecteur de présence} de tension, mais un très mauvais \emph{détecteur d'absence}. Pourquoi ? Parce que le néon peut rester éteint alors que la phase est bien présente :

\begin{itemize}
\item le \textbf{néon est grillé} ou la résistance interne est coupée (testeur tombé par terre) ;
\item l'utilisateur porte des \textbf{chaussures à semelles isolantes} ou est sur un sol très sec : la capacité corps-terre devient trop faible, le courant ne passe plus, le néon reste éteint ;
\item le \textbf{neutre est coupé} en amont : la boucle ne se referme plus, mais la phase reste dangereuse ;
\item la tension est \textbf{inférieure à environ \SI{100}{V}} : le néon ne s'allume pas, pourtant une tension de \SI{50}{V} peut déjà être dangereuse en milieu humide.
\end{itemize}

Dans tous ces cas, le testeur affiche un \textbf{faux négatif} : il dit \og pas de tension \fg{} alors que le conducteur est sous tension. De plus, le testeur \textbf{ne mesure pas} la tension : il ne peut pas distinguer \SI{230}{V} de \SI{400}{V}.

\begin{attention}[DANGER MORTEL : jamais de VAT au tournevis testeur]
Il est \textbf{formellement interdit} d'utiliser un tournevis testeur pour effectuer la \cle{Vérification d'Absence de Tension} (\cle{VAT}) avant d'intervenir sur une installation. Un faux négatif (néon grillé, semelles isolantes, neutre coupé) peut faire croire que le circuit est hors tension et provoquer une \textbf{électrocution mortelle}. La VAT réglementaire ne peut être réalisée \textbf{qu'avec un contrôleur de tension (VAT) normé IEC\_61243-3}, de catégorie CAT III ou CAT IV, à deux pointes, fonctionnant sans pile pour la détection de tension (signalisation LED et sonore).
\end{attention}

\begin{definition}[VAT et contrôleur de tension]
La \cle{VAT} (Vérification d'Absence de Tension) est l'opération de sécurité qui consiste à s'assurer qu'un circuit est bien hors tension avant d'y travailler. Elle s'effectue avec un \cle{contrôleur de tension} normé : appareil à \textbf{deux pointes de touche} que l'on place entre phase et neutre, puis entre phase et terre. Il signale la tension par LED et par un \textbf{bip sonore}, et fonctionne même avec ses piles déchargées pour la fonction détection.
\end{definition}

\begin{popfigure}
\begin{center}
\begin{tikzpicture}[font=\small]
% Tournevis testeur
\draw[fill=popblueL, rounded corners=2pt] (0,-0.35) rectangle (2.6,0.35);
\draw[fill=popmuted!50] (-1.4,-0.1) rectangle (0,0.1);
\draw[fill=popmuted!70] (2.6,-0.28) rectangle (3.0,0.28);
\draw[fill=poporange] (1.0,-0.18) rectangle (1.7,0.18);
\node at (1.35,0) {\tiny néon};
\node[below=2pt] at (0.8,-0.4) {Tournevis testeur};
\node[below=14pt, align=center, text width=4.2cm] at (0.8,-0.4) {\footnotesize 1 pointe -- détecte la phase\\ \footnotesize corps dans le circuit};
% Contrôleur VAT
\begin{scope}[xshift=8cm]
\draw[fill=poppinkL, rounded corners=4pt] (-0.9,-0.9) rectangle (0.9,0.9);
\node at (0,0.45) {\footnotesize VAT};
\draw[fill=popgreenL] (-0.5,0.05) rectangle (0.5,0.3);
\node at (0,0.17) {\tiny LED + bip};
\draw[thick] (-0.4,-0.9) -- (-0.9,-2.0);
\draw[thick] (0.4,-0.9) -- (0.9,-2.0);
\draw[fill=popmuted!70] (-1.0,-2.2) rectangle (-0.8,-1.9);
\draw[fill=popmuted!70] (0.8,-2.2) rectangle (1.0,-1.9);
\node[below=2pt] at (0,-2.2) {Contrôleur VAT (IEC\_61243-3)};
\node[below=14pt, align=center, text width=4.6cm] at (0,-2.2) {\footnotesize 2 pointes -- CAT III / IV\\ \footnotesize seul outil autorisé pour la VAT};
\end{scope}
\end{tikzpicture}
\end{center}
\legende{Comparaison : à gauche, le tournevis testeur (une pointe, néon dans le manche transparent) ; à droite, le contrôleur de tension VAT (deux pointes, signalisation LED et sonore, catégories de sécurité CAT III/IV).}
\end{popfigure}

\courssub{Entretien et vérification du testeur}

Un testeur en panne est un testeur dangereux. Deux règles d'or :

\begin{enumerate}
\item \textbf{Avant chaque utilisation}, vérifier le néon sur une prise \emph{connue} sous tension : le néon doit s'allumer franchement. C'est le même réflexe que le contrôleur VAT, que l'on teste avant et après chaque mesure.
\item \textbf{Au moindre doute} (chute, manche fissuré, éclairage faible ou intermittent), remplacer l'outil : un tournevis testeur coûte quelques centaines de francs CFA, une vie n'a pas de prix.
\end{enumerate}

\begin{savaistu}[Et les testeurs \og bipolaires \fg{} ?]
Il existe des testeurs \textbf{bipolaires} (à deux pointes, avec pile) qui affichent la tension par une échelle de LEDs : \SI{12}{V}, \SI{24}{V}, \SI{230}{V}, \SI{400}{V}. Ils sont très utiles pour un diagnostic rapide (par exemple vérifier si une rallonge alimente bien la maison pendant une coupure ENEO). Mais attention : s'ils ne portent pas la norme IEC\_61243-3 et la catégorie CAT III/IV, ce ne sont \textbf{toujours pas} des appareils de VAT réglementaires. Seul un contrôleur normé garantit ta sécurité avant une intervention.
\end{savaistu}

\begin{aretenir}[L'essentiel de la section]
\begin{itemize}
\item Le tournevis testeur contient un \textbf{néon en série avec une résistance d'environ \SI{1}{M\Omega}} ; le circuit se referme par le corps de l'utilisateur et la capacité corps-terre.
\item Le courant traversant le corps est inférieur à \SI{1}{mA} : sans danger en usage normal sur \SI{230}{V}.
\item Il \textbf{détecte la phase} (néon allumé) mais ne mesure pas la tension et peut donner des faux négatifs.
\item \textbf{Jamais de VAT au tournevis testeur} : la vérification d'absence de tension exige un contrôleur de tension normé IEC\_61243-3, CAT III/IV.
\item Vérifier le testeur sur une prise connue avant chaque utilisation.
\end{itemize}
\end{aretenir}

\begin{exercice}[À toi de jouer]
Un électricien veut intervenir sur une prise de la maison. Il teste l'orifice phase avec son tournevis testeur : le néon reste éteint. Il conclut : \og pas de tension, je peux travailler \fg{}. (1) Citer trois causes possibles de l'extinction du néon alors que la phase est présente. (2) Sa conclusion est-elle sûre ? (3) Quel appareil aurait-il dû utiliser, et selon quelle norme ?
\end{exercice}

\begin{corrige}[Exercice 1]
(1) Le néon peut être grillé ; l'électricien peut porter des chaussures à semelles isolantes qui coupent la liaison capacitive corps-terre ; le neutre peut être coupé en amont, empêchant la boucle de se refermer. (2) Non : le néon éteint ne prouve pas l'absence de tension (faux négatif possible). (3) Il aurait dû réaliser une VAT avec un contrôleur de tension normé IEC\_61243-3, de catégorie CAT III ou CAT IV, à deux pointes, en le vérifiant avant et après la mesure.
\end{corrige}

\coursec{Sécurité électrique : dangers, protections et premiers secours}

Dans la section précédente, nous avons appris à utiliser le tournevis testeur pour repérer la phase sur une prise. Mais ce petit outil a une limite fondamentale : il t'indique qu'un danger existe, il ne te protège pas de ce danger. Or le réseau ENEO distribue du \SI{230}{V} en courant alternatif \SI{50}{Hz}, une énergie invisible, silencieuse et inodore. Chaque année, des accidents domestiques — fils dénudés, multiprises surchargées, appareils utilisés les pieds mouillés — provoquent des brûlures graves, voire des décès. Cette section est peut-être la plus importante de toute ta scolarité en électricité : elle peut un jour sauver ta vie ou celle d'un proche.

\courssub{Effets du courant électrique sur le corps humain}

\textbf{Situation de départ.} Pendant la saison des pluies à Douala, un élève touche la carcasse métallique d'un ventilateur défectueux en sortant de la douche. Il ressent une violente secousse et ne peut pas lâcher l'appareil. Pourquoi le courant l'a-t-il « collé » au ventilateur ? Pourquoi le fait d'être mouillé a-t-il tant aggravé la situation ?

\begin{definition}[Électrisation, électrocution, brûlure électrique]
L'\cle{électrisation} est le passage d'un courant électrique à travers le corps humain, avec des conséquences cliniques (douleur, contraction musculaire, troubles cardiaques, brûlures). L'\cle{électrocution} est une électrisation \emph{mortelle}. La \cle{brûlure électrique} est une lésion thermique provoquée par l'effet Joule ($P = R \cdot I^{2}$) dans les tissus ; elle est souvent profonde, avec un point d'entrée et un point de sortie du courant visibles sur la peau.
\end{definition}

\textbf{Ce n'est pas la tension qui tue, c'est le courant.} La tension $U$ est la « pression » qui pousse le courant ; mais la grandeur qui détermine la gravité de l'accident est l'\cle{intensité} $I$ qui traverse le corps, ainsi que la \emph{durée} du passage. D'après la loi d'Ohm :
\[ I = \frac{U}{R_{\text{corps}}} \]
Plus la résistance du corps est faible, plus le courant est dangereux.

\begin{propriete}[Résistance du corps humain]
La résistance électrique du corps humain varie énormément selon l'état de la peau :
\begin{itemize}
\item peau \textbf{sèche} et intacte : $R \approx \SI{100}{k\Omega}$ ;
\item peau \textbf{mouillée} ou blessée : $R \approx \SI{1}{k\Omega}$, soit 100 fois moins !
\end{itemize}
C'est pourquoi la salle de bain et la piscine sont des lieux à \emph{danger extrême} : la norme y impose un zonage (volumes 0, 1, 2) et l'usage obligatoire de la très basse tension de sécurité (TBTS \SI{12}{V}) à proximité de l'eau.
\end{propriete}

\textbf{Les seuils physiologiques (courant alternatif \SI{50}{Hz}).} Les normes internationales (CEI 60479-1) ont établi expérimentalement les effets du courant en fonction de son intensité et de sa durée :

\begin{center}
\begin{tabularx}{\linewidth}{|l|X|}
\hline
\rowcolor{popblueL} \textbf{Intensité} & \textbf{Effet sur le corps} \\
\hline
\SI{0,5}{mA} & \textbf{Perception} : léger picotement, seuil de sensation. \\
\hline
\SI{10}{mA} & \textbf{Réaction / refus} : contraction musculaire, on peut encore lâcher prise. \\
\hline
15 à \SI{20}{mA} & \textbf{Tétanisation} : les muscles se contractent et ne répondent plus ; impossible de lâcher l'objet. Si le courant passe par le thorax, les muscles respiratoires sont bloqués. \\
\hline
50 à \SI{100}{mA} pendant \SI{1}{s} & \textbf{Fibrillation ventriculaire} : le cœur bat de façon désordonnée et inefficace ; arrêt circulatoire, mort probable sans secours immédiat. \\
\hline
$>$ \SI{1}{A} & \textbf{Brûlures internes graves} : destruction des tissus par effet Joule, arrêt cardiaque. \\
\hline
\end{tabularx}
\end{center}

\begin{popfigure}
\begin{tikzpicture}
\begin{loglogaxis}[
  width=0.95\linewidth, height=7.2cm,
  xlabel={Durée du passage du courant $t$ (ms)},
  ylabel={Courant traversant le corps $I$ (mA)},
  xmin=10, xmax=10000, ymin=0.1, ymax=2000,
  grid=both, grid style={popline!40},
  legend style={at={(0.02,0.02)}, anchor=south west, font=\scriptsize},
]
% Zone AC-4 (fibrillation probable) : au-dessus de la courbe c1 approximee
\addplot[draw=none, fill=poporange!25, forget plot] coordinates {(100,50) (200,70) (500,200) (1000,500) (2000,900) (5000,2000) (10000,2000) (10000,50) (5000,50) (2000,50) (1000,50) (500,50) (200,50) (100,50)} -- cycle;
% Seuil de perception (0.5 mA)
\addplot[popblue, thick, domain=10:10000, samples=2] {0.5};
\addlegendentry{Seuil de perception \SI{0,5}{mA}}
% Seuil de tetanisation / lacher prise (~10 mA)
\addplot[popgreen, thick, domain=10:10000, samples=2] {10};
\addlegendentry{Seuil de « lâcher prise » $\approx$ \SI{10}{mA}}
% Courbe de fibrillation (approximation c1)
\addplot[poporange, very thick, smooth] coordinates {(100,50) (200,70) (500,200) (1000,500) (2000,900) (5000,2000)};
\addlegendentry{Seuil de fibrillation (courbe c1)}
% Point differentiel 30 mA / 300 ms
\addplot[only marks, mark=*, mark size=2.5pt, poppink] coordinates {(300,30)};
\addlegendentry{Point de fonctionnement du différentiel \SI{30}{mA} / \SI{300}{ms}}
\node[font=\scriptsize\bfseries, text=poporange!70!black] at (axis cs:1500,300) {Zone AC-4 : fibrillation probable};
\node[font=\scriptsize, text=popdark, align=center] at (axis cs:60,3) {Zone AC-2 :\\ sans danger};
\end{loglogaxis}
\end{tikzpicture}
\legende{Courbe « intensité -- durée » des effets du courant alternatif \SI{50}{Hz} sur le corps humain (d'après la norme CEI 60479-1). En dessous de la courbe c1 (zone AC-2/AC-3), le danger de fibrillation est négligeable ; au-dessus (zone AC-4), la fibrillation ventriculaire devient probable. Le point rose montre que le différentiel \SI{30}{mA}, qui coupe en moins de \SI{300}{ms}, se situe en zone sûre (AC-3).}
\end{popfigure}

\begin{aretenir}[L'essentiel sur les effets du courant]
Trois facteurs déterminent la gravité d'une électrisation : l'\cle{intensité} du courant, la \cle{durée} du contact et le \cle{trajet} du courant dans le corps (le trajet main gauche -- pieds, qui traverse le cœur, est le plus dangereux). Dès \SI{50}{mA} pendant une seconde, la vie est en danger.
\end{aretenir}

\courssub{Tension de contact et tension de pas}

Quand un appareil à carcasse métallique (frigo, machine à laver, fer à repasser) présente un défaut d'isolement, sa carcasse peut se retrouver portée au potentiel de la phase. Une personne qui la touche est alors soumise à la \cle{tension de contact}.

\begin{definition}[Tension de contact et tension de pas]
La \cle{tension de contact} $U_c$ est la tension appliquée au corps d'une personne touchant une masse métallique en défaut. Si le courant de défaut $I_d$ s'écoule vers la terre par la prise de terre de résistance $R_A$ :
\[ U_c = I_d \times R_A \]
La \cle{tension de pas} $U_p$ est la tension qui apparaît entre les deux pieds d'une personne marchant près d'un point où le courant s'écoule dans le sol (par exemple près d'un poteau ENEO foudroyé ou d'un câble tombé au sol). Le courant traverse alors le corps par le trajet jambe -- jambe.
\end{definition}

\begin{propriete}[Tensions limites de sécurité]
En régime permanent, la tension de contact ne doit pas dépasser :
\begin{itemize}
\item $U_L = \SI{50}{V}$ en local \textbf{sec} ;
\item $U_L = \SI{25}{V}$ en local \textbf{mouillé} (salle de bain, piscine, extérieur).
\end{itemize}
Le disjoncteur différentiel \SI{30}{mA} garantit en plus que la durée du choc reste inférieure à environ \SI{300}{ms}, ce qui maintient le corps dans la zone sans fibrillation.
\end{propriete}

\begin{attention}[Piège classique : le courant de fuite « invisible »]
Un défaut d'isolement ne fait pas nécessairement disjoncter l'installation : un courant de fuite de \SI{200}{mA} vers la terre ne déclenche \emph{pas} un disjoncteur \SI{16}{A}, mais il peut tuer une personne qui touche la carcasse. Seuls la prise de terre (qui canalise le défaut) et le différentiel (qui détecte la fuite) assurent la protection.
\end{attention}

\courssub{Les protections : contre les contacts directs et indirects}

\begin{definition}[Contacts directs et indirects]
Un \cle{contact direct} est le contact d'une personne avec une pièce normalement \emph{sous tension} (borne phase d'une prise, fil dénudé). Un \cle{contact indirect} est le contact avec une masse métallique \emph{accidentellement} mise sous tension par un défaut d'isolement (carcasse de machine à laver devenue « vivante »).
\end{definition}

\textbf{Protection contre les contacts directs :}
\begin{itemize}
\item l'\cle{isolement} des conducteurs (gaine isolante des fils) ;
\item les \cle{enveloppes et barrières} (boîtiers des prises, capots du tableau électrique) ;
\item la \cle{TBTS} (Très Basse Tension de Sécurité, \SI{12}{V} ou \SI{24}{V}) : à cette tension, le courant traversant le corps reste infime même sur peau mouillée ; c'est elle qui alimente les lampes de piscine.
\end{itemize}

\textbf{Protection contre les contacts indirects :}
\begin{itemize}
\item la \cle{mise à la terre} des masses métalliques associée au \cle{disjoncteur différentiel} (ou interrupteur différentiel en tête d'installation) : le défaut s'écoule vers la terre, le différentiel détecte la fuite et coupe le circuit en une fraction de seconde ;
\item la \cle{double isolation} (matériel de \textbf{classe II}, symbole : double carré) : l'appareil possède deux niveaux d'isolement, sa carcasse ne peut pas devenir dangereuse ; il n'a pas besoin de terre (ex. : sèche-cheveux, perceuse à poignée plastique) ;
\item la \cle{séparation électrique} (transformateur d'isolement), utilisée pour certains équipements.
\end{itemize}

\courssub{Les dispositifs de protection de l'installation}

\begin{center}
\begin{tabularx}{\linewidth}{|l|X|X|}
\hline
\rowcolor{popblueL} \textbf{Dispositif} & \textbf{Danger combattu} & \textbf{Principe} \\
\hline
Disjoncteur divisionnaire (magnéto-thermique) & \textbf{Surcharge} (trop d'appareils) et \textbf{court-circuit} & Déclenche quand $I$ dépasse son calibre (ex. \SI{16}{A}) : effet thermique (surcharge lente) et magnétique (court-circuit brutal). Protège les \textbf{fils}. \\
\hline
Interrupteur différentiel \SI{30}{mA} (en tête d'installation) & \textbf{Fuite de courant} vers la terre, contact indirect ou direct & Compare le courant qui part (phase) et celui qui revient (neutre) ; si la différence dépasse \SI{30}{mA}, il coupe en moins de \SI{300}{ms}. Protège les \textbf{personnes}. \\
\hline
Parafoudre & \textbf{Surtension} due à la foudre & Écoule vers la terre l'onde de choc de la foudre avant qu'elle n'atteigne les appareils. \\
\hline
Prise de terre & Contact indirect & Offre au courant de défaut un chemin facile vers le sol. \\
\hline
\end{tabularx}
\end{center}

\begin{attention}[ERREUR FRÉQUENTE : « mon disjoncteur \SI{16}{A} me protège »]
Faux ! Le disjoncteur \SI{16}{A} protège les \textbf{fils} de l'installation contre l'échauffement : il ne déclenche qu'au-delà de $16\,000$ mA. Or \SI{50}{mA} suffisent à tuer. Entre les deux, le disjoncteur reste silencieux pendant que le corps est en danger. Seul le \cle{différentiel \SI{30}{mA}}, qui détecte la \emph{différence} phase -- neutre, protège la vie humaine : $\SI{30}{mA} \ll \SI{16000}{mA}$, soit plus de 500 fois plus sensible !
\end{attention}

\begin{savaistu}[La foudre, ennemie numéro un des installations au Cameroun]
Le Cameroun connaît une kéraunicité (nombre de jours d'orage par an) parmi les plus élevées du monde, surtout dans les régions de l'Ouest, du Centre et du Littoral. La foudre est une cause majeure d'incendies électriques et de destruction d'appareils (téléviseurs, décodeurs, réfrigérateurs). La norme NF C 15-100 (\S~443) rend le \cle{parafoudre} (de type 1 ou 2) obligatoire lorsque le bâtiment possède un paratonnerre ou se situe en zone à forte kéraunicité : au Cameroun, son installation au tableau principal est donc \emph{vivement recommandée}. Pense aussi à débrancher les appareils sensibles pendant les gros orages.
\end{savaistu}

\courssub{Exemple chiffré : le drame du carrelage mouillé}

\begin{exoresolu}[Électrisation sur sol humide]
\textbf{Énoncé.} Une personne pieds nus sur un carrelage humide (résistance du corps $R = \SI{2000}{\Omega}$) touche la phase \SI{230}{V} d'un fil dénudé. 1) Calculer l'intensité $I$ qui traverse son corps. 2) Quel est l'effet physiologique prévisible ? 3) Que change la présence d'un différentiel \SI{30}{mA} ?
\tcblower
\textbf{Solution.}
1) D'après la loi d'Ohm :
\[ I = \frac{U}{R} = \frac{230}{2000} = \SI{0,115}{A} = \SI{115}{mA} \]
2) D'après le tableau des seuils, \SI{115}{mA} pendant une seconde provoque une \textbf{fibrillation ventriculaire quasi certaine} : sans protection, l'accident est mortel en moins d'une seconde.
3) Avec un différentiel \SI{30}{mA} : dès que la fuite atteint \SI{30}{mA}, le dispositif coupe le circuit en quelques dizaines de millisecondes (généralement $< \SI{40}{ms}$ à $5 I_{\Delta n}$). Le point $(\SI{30}{mA}\,;\,\SI{40}{ms})$ se situe très largement dans la zone AC-3 de la courbe, bien en deçà du seuil de fibrillation : la victime ressent une violente secousse mais \textbf{survit probablement}. Le différentiel sauve des vies.
\end{exoresolu}

\courssub{Conduite à tenir face à une victime d'électrisation}

\begin{methode}[Les gestes qui sauvent, dans l'ordre]
\begin{enumerate}
\item \textbf{COUPER le courant} immédiatement : disjoncteur général, ou débrancher la prise. Si c'est impossible, \textbf{écarter la victime} du conducteur avec un objet \emph{isolant} (manche en bois sec, bâton plastique, linge sec) sans jamais la toucher directement.
\item \textbf{NE JAMAIS TOUCHER} la victime tant qu'elle est sous tension : le courant te traverserait aussi, et il y aurait deux victimes au lieu d'une.
\item \textbf{Alerter les secours} : appeler le \textbf{112} (numéro d'urgence européen, valide sur GSM au Cameroun) ou le \textbf{119} (sapeurs-pompiers au Cameroun), en précisant le lieu et la nature de l'accident.
\item Si la victime est \textbf{inconsciente mais respire} : la placer en \cle{PLS} (Position Latérale de Sécurité), sur le côté, pour éviter l'étouffement, et surveiller sa respiration.
\item Si la victime \textbf{ne respire plus} (arrêt cardiaque) : commencer la \cle{RCP} (Réanimation Cardio-Pulmonaire) : compressions thoraciques au centre de la poitrine, rythme de 100 à 120 par minute, et utiliser un \cle{DAE} (Défibrillateur Automatisé Externe) s'il est disponible, en suivant ses instructions vocales.
\end{enumerate}
\end{methode}

\begin{popfigure}
\begin{tikzpicture}[
  etape/.style={rectangle, rounded corners=6pt, draw=popblue, fill=popblueL, thick, minimum width=3.1cm, minimum height=1.5cm, align=center, font=\scriptsize\bfseries, text=popdark},
  fleche/.style={-{Stealth[length=3mm]}, very thick, poporange},
]
\node[etape, align=center] (e1) at (0,0){1. COUPER\\le courant};
\node[etape, align=center] (e2) at (3.6,0){2. Écarter avec\\un isolant sec};
\node[etape, align=center] (e3) at (7.2,0){3. Appeler\\le 112 / 119};
\node[etape, align=center] (e4) at (10.8,0){4. PLS si\\inconscient};
\node[etape, draw=poppink, fill=poppinkL, align=center] (e5) at (7.2,-2.4){5. RCP + DAE\\si arrêt cardiaque};
\draw[fleche] (e1) -- (e2);
\draw[fleche] (e2) -- (e3);
\draw[fleche] (e3) -- (e4);
\draw[fleche] (e4.south) |- (e5.east);
% Symbole differentiel
\node[rectangle, rounded corners=6pt, draw=popgreen, fill=popgreenL, very thick, minimum width=4.4cm, minimum height=1.5cm, align=center, font=\scriptsize\bfseries, text=popdark] (diff) at (2.2,-2.4) {Différentiel \SI{30}{mA}\\« Sauve des vies »};
\draw[popgreen, very thick] (1.1,-2.15) circle (0.32cm);
\node[font=\footnotesize\bfseries, text=popgreen] at (1.1,-2.15) {30};
\end{tikzpicture}
\legende{Infographie « Gestes qui sauvent » : la chaîne des actions face à une victime d'électrisation. La première règle absolue : couper le courant avant tout contact. Le différentiel \SI{30}{mA}, en amont, est le gardien silencieux de la vie domestique.}
\end{popfigure}

\courssub{Bonnes pratiques quotidiennes et procédure VAT}

\textbf{Les bons réflexes à la maison :}
\begin{itemize}
\item Ne jamais \textbf{surcharger une multiprise} : la puissance totale branchée ne doit pas dépasser environ $P_{max} = \SI{3500}{W}$ (soit \SI{16}{A} sous \SI{230}{V}). Un four (\SI{2000}{W}) + une bouilloire (\SI{1800}{W}) sur la même multiprise, c'est l'échauffement garanti et le risque d'incendie.
\item \textbf{Interdiction formelle} de brancher des multiprises \emph{en série} (l'une sur l'autre).
\item Remplacer immédiatement tout fil \textbf{effiloché} ou dénudé ; ne jamais réparer « provisoirement » avec du tissu ou du papier.
\item Installer des \textbf{cache-prises} dans les maisons avec de jeunes enfants.
\item \textbf{Éloigner l'eau de l'électricité} : jamais de sèche-cheveux près du lavabo plein, jamais de rallonge dans une flaque, ne jamais toucher un appareil électrique avec les mains mouillées.
\end{itemize}

\begin{methode}[Procédure VAT : Vérification d'Absence de Tension]
Avant toute intervention sur un circuit (changer une prise, réparer un plafonnier) :
\begin{enumerate}
\item \textbf{Identifier} le circuit concerné au tableau électrique.
\item \textbf{Couper} le disjoncteur divisionnaire correspondant (et le général si possible).
\item \textbf{Verrouiller} (cadenas ou consignation) et poser une \textbf{signalétique} (« Ne pas réenclencher — intervention en cours »).
\item \textbf{Vérifier l'absence de tension} avec un contrôleur normé (type VAT), en trois temps : tester d'abord l'appareil sur une prise \emph{vivante} (il doit réagir), puis tester le circuit coupé (il ne doit rien indiquer), puis re-tester sur la prise vivante (pour s'assurer que l'appareil fonctionne toujours). C'est le test « avant -- pendant -- après ».
\item En haute tension (HTA/HTB), on ajoute une mise à la terre et en court-circuit ; cette étape \textbf{ne concerne pas} la basse tension domestique.
\end{enumerate}
\end{methode}

\begin{exercice}[À toi de jouer]
1) Une personne à la peau sèche ($R = \SI{100}{k\Omega}$) touche accidentellement une borne à \SI{230}{V}. Calculer l'intensité qui la traverse et conclure sur le danger ressenti. 2) Même question avec la peau mouillée ($R = \SI{1}{k\Omega}$). 3) Expliquer pourquoi le différentiel \SI{30}{mA} est indispensable dans le second cas.
\end{exercice}

\begin{corrige}[Exercice]
1) $I = \dfrac{U}{R} = \dfrac{230}{100\,000} = \SI{2,3}{mA}$ : sensation de picotement, au-dessus du seuil de perception (\SI{0,5}{mA}) mais sans gravité immédiate. 2) $I = \dfrac{230}{1000} = \SI{230}{mA}$ : fibrillation ventriculaire quasi certaine si le contact dure une seconde ; danger de mort. 3) Le différentiel détecte la fuite dès \SI{30}{mA} et coupe en moins de \SI{300}{ms}, avant que le courant n'atteigne une durée et une intensité mortelles : c'est la seule protection efficace pour la personne.
\end{corrige}

\begin{aretenir}[Ce qu'il faut retenir de cette section]
Le danger, c'est le \cle{courant}, pas la tension affichée. Dès \SI{50}{mA} pendant une seconde, le cœur peut se fibriller. La peau mouillée divise la résistance du corps par 100 et multiplie le danger par 100. Le \cle{disjoncteur} protège les fils, le \cle{différentiel \SI{30}{mA}} protège les personnes, la \cle{prise de terre} canalise les défauts, le \cle{parafoudre} protège contre la foudre. Face à une victime : couper, ne jamais toucher sous tension, alerter le 112/119, PLS ou RCP. Ces gestes, appris aujourd'hui, sauveront des vies demain.
\end{aretenir}

\coursec{Lecture de plans, schémas unifilaires et mise en situation}

Dans la section précédente, tu as appris à te protéger : fusibles, disjoncteurs, différentiel, terre. Mais toutes ces protections ne servent à rien si l'installation est mal conçue dès le départ. Avant de poser le moindre fil, l'électricien professionnel travaille sur des \cle{documents} : un plan d'implantation, un schéma unifilaire, un schéma multifilaire. C'est exactement comme un chauffeur de moto-taxi qui prépare son itinéraire avant de partir de Douala vers Bafoussam : on ne roule pas au hasard. Dans cette dernière section, tu vas apprendre à \emph{lire} ces documents, puis à les \emph{produire} toi-même sur un cas concret complet : une pièce avec une lampe commandée par deux va-et-vient et deux prises de courant.

\courssub{Les symboles normalisés : l'alphabet de l'électricien}

\begin{experience}[Observer un plan de maison]
\textbf{Document observé :} le plan électrique d'une maison en construction à Yaoundé. \textbf{Observation :} le plan est couvert de petits symboles identiques d'une maison à l'autre : un cercle barré pour une lampe, un cercle avec un trait pour un interrupteur, un demi-cercle pour une prise. \textbf{Question :} pourquoi des symboles et pas des dessins réalistes ? \textbf{Hypothèse :} pour que tous les électriciens, partout, comprennent le même plan sans ambiguïté. \textbf{Vérification :} on montre le plan à deux électriciens différents ; les deux décrivent exactement la même installation. \textbf{Conclusion :} les symboles sont un \emph{langage universel} normalisé (norme NF EN 60617).
\end{experience}

\begin{definition}[Symbole normalisé]
Un \cle{symbole normalisé} est un dessin conventionnel, défini par la norme \cle{NF EN 60617}, qui représente un appareil électrique (interrupteur, prise, lampe, disjoncteur...) sur un plan ou un schéma, indépendamment de sa marque ou de son aspect réel.
\end{definition}

Voici les symboles à connaître pour le BEPC :

\begin{center}
\begin{tabularx}{\linewidth}{|l|X|l|X|}
\hline
\rowcolor{popblueL} \textbf{Appareil} & \textbf{Symbole (plan)} & \textbf{Appareil} & \textbf{Symbole (plan)} \\
\hline
Lampe / point lumineux & cercle barré d'une croix & Prise 2P+T & demi-cercle avec un petit trait (terre) \\
\hline
Interrupteur simple & cercle + un trait oblique & Va-et-vient & cercle + deux traits obliques \\
\hline
Poussoir (bouton) & cercle + trait fléché & Prise de terre & symbole terre (trois traits) \\
\hline
Disjoncteur & rectangle avec un crochet & Interrupteur différentiel & disjoncteur + boucle (toroïde) \\
\hline
Compteur & rectangle inscrit \og kWh \fg{} & Tableau (TGBT) & rectangle marqué \og T \fg{} \\
\hline
\end{tabularx}
\end{center}

\begin{attention}[Ne pas confondre les deux familles de schémas]
Sur un \textbf{plan d'implantation} (vue du dessus de la maison), les symboles indiquent la \emph{position} des appareils. Sur un \textbf{schéma électrique}, les symboles indiquent les \emph{liaisons électriques}. Le même symbole de prise ne veut pas dire la même chose selon le document ! Regarde toujours le titre et la légende du document avant d'interpréter.
\end{attention}

\courssub{Le schéma unifilaire : la vue d'ensemble}

\begin{definition}[Schéma unifilaire]
Le \cle{schéma unifilaire} est une représentation simplifiée d'une installation électrique dans laquelle \textbf{tous les conducteurs d'un même circuit} (Phase, Neutre, Terre) sont représentés par \textbf{un seul trait}. Il sert à l'étude des protections et au dimensionnement (calibres, sections de câbles).
\end{definition}

Sur ce trait unique, l'électricien écrit des indications codées. Exemple de mention rencontrée sur un schéma unifilaire :

\begin{center}
\textbf{3G2,5 --- Disjoncteur 16 A --- 12 m --- ICTA Ø20 encastré}
\end{center}

Décodons ce langage, mot à mot :
\begin{itemize}
\item \cle{3G2,5} : câble de \textbf{3 conducteurs} dont un conducteur de protection (\textbf{G} = Ground, vert/jaune), chaque conducteur ayant une \textbf{section de \SI{2,5}{mm^2}} (section normalisée des circuits de prises).
\item \cle{Disjoncteur 16 A} : calibre de la protection contre les surintensités, adapté à la section \SI{2,5}{mm^2}.
\item \cle{12 m} : longueur du circuit depuis le tableau.
\item \cle{ICTA Ø20 encastré} : mode de pose, ici une gaine plastique ICTA de diamètre \SI{20}{mm} noyée dans le mur.
\end{itemize}

\begin{aretenir}[Sections et calibres à connaître]
\begin{itemize}
\item Circuit \textbf{éclairage} : conducteurs \SI{1,5}{mm^2}, disjoncteur \textbf{10 A} (ou 16 A), maximum 8 points lumineux.
\item Circuit \textbf{prises} : conducteurs \SI{2,5}{mm^2}, disjoncteur \textbf{16 A} (ou 20 A), maximum 8 prises (5 avec 20 A selon les cas).
\item La section du câble et le calibre du disjoncteur vont \textbf{toujours ensemble} : augmenter le calibre sans augmenter la section fait chauffer le câble : danger d'incendie !
\end{itemize}
\end{aretenir}

\courssub{Le schéma multifilaire : la recette de câblage}

Le schéma unifilaire est une carte routière ; le \cle{schéma multifilaire} (ou schéma développé) est le mode d'emploi détaillé. Il montre \textbf{tous les fils} (Phase en rouge, Neutre en bleu, Terre en vert/jaune) et \textbf{toutes les connexions exactes} dans chaque boîte : boîte de dérivation, boîte d'interrupteur, boîte de prise. C'est le document indispensable pour réaliser le câblage réel sans erreur.

\begin{propriete}[Câblage d'un va-et-vient]
Dans un montage \cle{va-et-vient}, les deux interrupteurs sont reliés entre eux par \textbf{deux fils navettes} (fils voyageurs, souvent orange et violet). La \textbf{phase} arrive sur la borne commune (borne L, repère rouge) du \emph{premier} va-et-vient ; la \textbf{lampe} est branchée sur la borne commune du \emph{second} va-et-vient. Le \textbf{neutre} va directement à la lampe, sans passer par les interrupteurs. Ainsi, quel que soit l'état de départ, chaque interrupteur peut allumer ou éteindre la lampe.
\end{propriete}

\begin{popfigure}
\begin{center}
\begin{tikzpicture}[scale=0.9, every node/.style={font=\small}]
% Tableau
\draw[fill=popblueL, thick] (0,0) rectangle (1.6,2.2);
\node at (0.8,2.5) {\textbf{Tableau}};
\node at (0.8,1.5) {Disj.};
\node at (0.8,1.0) {10 A};
\node at (0.8,0.4) {Diff.};
% Va-et-vient 1
\draw[fill=poporangeL, thick] (3.4,0.4) rectangle (5.0,2.0);
\node at (4.2,2.3) {\textbf{Va-et-vient 1}};
\node at (3.7,1.6) {L};
\node at (3.7,0.8) {1};
\node at (4.7,0.8) {2};
% Va-et-vient 2
\draw[fill=poporangeL, thick] (7.4,0.4) rectangle (9.0,2.0);
\node at (8.2,2.3) {\textbf{Va-et-vient 2}};
\node at (8.7,1.6) {L};
\node at (7.7,0.8) {1};
\node at (8.7,0.8) {2};
% Lampe
\draw[thick] (11.2,1.2) circle (0.55);
\draw[thick] (10.81,0.81) -- (11.59,1.59);
\draw[thick] (10.81,1.59) -- (11.59,0.81);
\node at (11.2,2.1) {\textbf{Lampe}};
% Fils phase
\draw[red, thick] (1.6,1.9) -- (3.4,1.9) -- (3.7,1.9) -- (3.7,1.7);
\node[red, above] at (2.5,1.9) {Phase};
% Navettes
\draw[orange, thick] (3.7,0.9) -- (3.7,0.35) -- (7.7,0.35) -- (7.7,0.9);
\node[orange, below] at (5.7,0.35) {Navette 1};
\draw[poppurple, thick] (4.7,0.9) -- (4.7,0.05) -- (8.7,0.05) -- (8.7,0.9);
\node[poppurple, below] at (6.7,0.05) {Navette 2};
% Retour lampe
\draw[red, thick] (8.7,1.7) -- (8.7,1.9) -- (10.65,1.9) -- (10.65,1.5);
\node[red, above] at (9.7,1.9) {Retour lampe};
% Neutre
\draw[blue, thick] (1.6,1.2) -- (2.2,1.2) -- (2.2,-0.6) -- (10.65,-0.6) -- (10.65,0.9);
\node[blue, below] at (4.5,-0.6) {Neutre (direct)};
% Terre
\draw[popgreen, thick] (1.6,0.5) -- (1.9,0.5) -- (1.9,-1.0) -- (11.75,-1.0) -- (11.75,0.9);
\node[popgreen, below] at (5.0,-1.0) {Terre (vert/jaune)};
% Prises
\draw[fill=popgreenL, thick] (3.0,-2.6) rectangle (4.6,-1.6);
\node at (3.8,-1.35) {\textbf{Prise 1}};
\node at (3.8,-2.1) {L\; N\; T};
\draw[fill=popgreenL, thick] (7.0,-2.6) rectangle (8.6,-1.6);
\node at (7.8,-1.35) {\textbf{Prise 2}};
\node at (7.8,-2.1) {L\; N\; T};
\draw[red, thick] (1.6,1.9) -- (1.6,-2.1) -- (3.0,-2.1);
\draw[blue, thick] (2.2,-0.6) -- (2.2,-2.3) -- (3.0,-2.3);
\draw[popgreen, thick] (1.9,-1.0) -- (1.9,-1.9) -- (3.0,-1.9);
\draw[red, thick] (4.6,-2.1) -- (7.0,-2.1);
\draw[blue, thick] (4.6,-2.3) -- (7.0,-2.3);
\draw[popgreen, thick] (4.6,-1.9) -- (7.0,-1.9);
\node at (5.8,-2.9) {Repiquage Prise 1 $\to$ Prise 2};
\end{tikzpicture}
\end{center}
\legende{Schéma multifilaire simplifié : lampe commandée par deux va-et-vient (navettes orange et violette) et deux prises 2P+T repiquées. Phase en rouge, neutre en bleu, terre en vert. Les bornes L, 1, 2 des va-et-vient sont repérées.}
\end{popfigure}

\begin{attention}[L'erreur classique du va-et-vient]
L'erreur de câblage la plus fréquente : \textbf{inverser la phase et une navette} sur la borne commune L. Résultat : la lampe ne s'allume que lorsque les deux interrupteurs sont dans une position particulière (fonctionnement en \og série \fg), au lieu de pouvoir allumer et éteindre depuis \emph{chacun} des deux points (fonction \og ou exclusif \fg). Avant de refermer les boîtes, teste toujours les 4 combinaisons des deux interrupteurs !
\end{attention}

\courssub{Le plan d'implantation : où placer les appareils}

Le \cle{plan d'implantation} (ou plan architectural) est une vue du dessus de la pièce, à l'échelle, sur laquelle on positionne les points lumineux, les prises, les interrupteurs et le tableau, avec les symboles normalisés. Il respecte des \textbf{hauteurs standard} mesurées depuis le sol fini :

\begin{center}
\begin{tabularx}{\linewidth}{|l|X|}
\hline
\rowcolor{popblueL} \textbf{Appareil} & \textbf{Hauteur standard} \\
\hline
Interrupteurs, va-et-vient & \SI{1,10}{m} (accessible, hors portée des petits enfants pour la sécurité) \\
\hline
Prises de courant & \SI{0,30}{m} (pièces courantes) ; \SI{0,90}{m} au-dessus du plan de travail en cuisine \\
\hline
Tableau de répartition & entre \SI{1,30}{m} et \SI{1,60}{m} (lecture et manœuvre faciles) \\
\hline
Point lumineux & au plafond, centré ou selon l'usage de la pièce \\
\hline
\end{tabularx}
\end{center}

\begin{popfigure}
\begin{center}
\begin{tikzpicture}[scale=1.05, every node/.style={font=\small}]
% Pièce 4m x 5m
\draw[very thick, fill=popcream] (0,0) rectangle (5,4);
\node at (2.5,3.75) {\textbf{Pièce 5 m $\times$ 4 m}};
% Porte
\draw[very thick, popcream] (0.8,0) -- (1.8,0);
\draw[thin, dashed] (1.8,0) arc (0:90:1.0);
\node at (1.3,-0.3) {Porte};
% Fenêtre
\draw[very thick, popcream] (3.2,4) -- (4.4,4);
\draw[thin] (3.2,3.92) -- (4.4,3.92);
\node at (3.8,4.3) {Fenêtre};
% Tableau
\draw[fill=popblueL, thick] (0.05,2.6) rectangle (0.35,3.2);
\node at (0.75,2.9) {T};
% Va-et-vient près porte
\draw[thick, fill=white] (1.0,0.25) circle (0.16);
\draw[thick] (1.0,0.25) -- (1.3,0.55);
\draw[thick] (1.0,0.25) -- (1.35,0.35);
\node[right] at (1.35,0.4) {V1 (1,10 m)};
% Va-et-vient près fenêtre
\draw[thick, fill=white] (4.0,3.7) circle (0.16);
\draw[thick] (4.0,3.7) -- (4.3,3.4);
\draw[thick] (4.0,3.7) -- (4.35,3.5);
\node[left] at (3.95,3.45) {V2 (1,10 m)};
% Lampe centre
\draw[thick] (2.5,2.0) circle (0.3);
\draw[thick] (2.29,1.79) -- (2.71,2.21);
\draw[thick] (2.29,2.21) -- (2.71,1.79);
\node[right] at (2.85,2.0) {Lustre};
% Prises
\draw[thick, fill=popgreenL] (4.85,1.2) arc (-90:90:0.22);
\draw[thick] (4.85,1.42) -- (4.85,1.62);
\node[left] at (4.8,1.4) {P1 (0,30 m)};
\draw[thick, fill=popgreenL] (4.85,2.6) arc (-90:90:0.22);
\draw[thick] (4.85,2.82) -- (4.85,3.02);
\node[left] at (4.8,2.8) {P2 (0,30 m)};
% Cotes
\draw[<->, thin] (0,-0.6) -- (5,-0.6);
\node at (2.5,-0.85) {5 m};
\draw[<->, thin] (5.5,0) -- (5.5,4);
\node[rotate=90] at (5.75,2) {4 m};
\end{tikzpicture}
\end{center}
\legende{Extrait de plan d'implantation : tableau T, deux va-et-vient V1 (côté porte) et V2 (côté fenêtre), lustre au centre, deux prises P1 et P2 le long du mur des bureaux.}
\end{popfigure}

\begin{savaistu}[La norme NF C 15-100 veille sur ta maison]
La norme \cle{NF C 15-100} impose un nombre \textbf{minimum} de prises par pièce : un séjour de moins de \SI{20}{m^2} doit avoir au moins \textbf{5 prises}, plus une prise par tranche de \SI{4}{m^2} supplémentaire ; la cuisine exige \textbf{6 prises} dont 4 au-dessus du plan de travail ; dans la salle de bain, la prise est \textbf{interdite} dans les volumes 0, 1 et 2 (proches de l'eau) et tolérée seulement en volume 3 ou hors volume. Ces règles, inspirées des normes françaises, guident aussi les bons électriciens camerounais : moins de prises imposées, c'est plus de rallonges dangereuses !
\end{savaistu}

\courssub{Mise en situation : du plan au devis}

\begin{methode}[Méthode complète de conception d'une installation]
\begin{enumerate}
\item \textbf{Lire le plan d'implantation} : repérer les appareils, leurs positions et les hauteurs.
\item \textbf{Tracer le schéma unifilaire} : un trait par circuit, avec sections, calibres, longueurs ; vérifier les protections.
\item \textbf{Dessiner le schéma multifilaire} : tous les fils (L rouge, N bleu, PE vert/jaune) et toutes les connexions dans les boîtes.
\item \textbf{Établir la liste de matériel} : longueurs de câbles, gaines, nombre de boîtes, appareillages.
\item \textbf{Rédiger le devis} : quantités $\times$ prix unitaires.
\end{enumerate}
\end{methode}

\begin{exoresolu}[Câblage virtuel d'une pièce 4 m $\times$ 5 m]
\textbf{Énoncé.} Une chambre-bureau de 4 m $\times$ 5 m comporte : un lustre au centre commandé par deux va-et-vient (V1 près de la porte, V2 près de la fenêtre) et deux prises 2P+T pour les bureaux. Les longueurs de gaines ICTA Ø20 mesurées sur le plan sont : Tableau $\to$ V1 : 6 m ; V1 $\to$ V2 (navettes) : 4 m ; V2 $\to$ lustre : 3 m ; Tableau $\to$ P1 : 8 m ; P1 $\to$ P2 : 3 m. 1) Établir la liste des câbles. 2) Calculer les longueurs totales par type de câble.
\tcblower
\textbf{Solution.}
1) \textbf{Circuit éclairage} (câble 3G1,5, disjoncteur 10 A) :
\begin{itemize}
\item Tableau $\to$ V1 : 6 m de 3G1,5 (phase + neutre + terre de passage) ;
\item V1 $\to$ V2 : 4 m de 3G1,5 (deux navettes + terre de passage) ;
\item V2 $\to$ lustre : 3 m de 3G1,5 (retour lampe + neutre + terre).
\end{itemize}
\textbf{Circuit prises} (câble 3G2,5, disjoncteur 16 A) :
\begin{itemize}
\item Tableau $\to$ P1 : 8 m de 3G2,5 ;
\item P1 $\to$ P2 (repiquage) : 3 m de 3G2,5.
\end{itemize}
2) \textbf{Totaux} :
\begin{align*}
L_{3G2,5} &= 8 + 3 = \SI{11}{m} \quad (\text{arrondi chantier : } \approx \SI{12}{m})\\
L_{3G1,5} &= 6 + 4 + 3 = \SI{13}{m} \quad (\text{arrondi chantier : } \approx \SI{15}{m})
\end{align*}
On ajoute toujours une marge d'environ 10 \% pour les chutes et les réserves dans les boîtes. Le devis comprend aussi : 2 va-et-vient, 2 prises 2P+T, 1 douille + lustre, 4 boîtes d'encastrement, 1 disjoncteur 10 A, 1 disjoncteur 16 A, gaines ICTA Ø20 et borniers de connexion.
\end{exoresolu}

\courssub{Vérification avant mise sous tension}

Un électricien sérieux ne branche jamais le disjoncteur général sans avoir vérifié son travail. Voici la check-list professionnelle :

\begin{methode}[Les 5 vérifications avant la mise sous tension]
\begin{enumerate}
\item \textbf{Continuité de la terre} : au sondeur (ou multimètre en mode continuité), vérifier que le fil vert/jaune relie bien chaque prise à la barrette de terre du tableau : l'appareil doit sonner.
\item \textbf{Isolement} : au \cle{mégaohmmètre} sous \SI{500}{V}, mesurer la résistance d'isolement entre conducteurs actifs et terre : elle doit être \textbf{supérieure à \SI{0,5}{M\ohm}}. En dessous, il y a un défaut d'isolement (fil pincé, dénudage mal protégé).
\item \textbf{Polarité phase/neutre} : au tournevis testeur (section 4), vérifier que la phase arrive bien sur la bonne borne des prises et sur le commun du premier va-et-vient.
\item \textbf{Serrage des bornes} : resserrer chaque vis de borne ; un contact desserré chauffe et peut provoquer un incendie.
\item \textbf{Étiquetage des circuits} : coller sur le tableau une étiquette par disjoncteur (\og Prises chambre \fg, \og Éclairage séjour \fg) : en cas de panne, on coupe le bon circuit immédiatement.
\end{enumerate}
\end{methode}

\begin{exercice}[À toi de jouer]
Sur le plan d'implantation de la figure précédente, on ajoute une troisième prise P3 repiquée sur P2, à 2 m de celle-ci. 1) Quel type de câble utilise-t-on pour ce prolongement ? 2) Quelle longueur totale de 3G2,5 faut-il désormais prévoir (sans la marge) ? 3) Le circuit prises comporte maintenant 3 prises : le disjoncteur 16 A convient-il toujours ? Justifie.
\end{exercice}

\begin{corrige}[Exercice 1]
1) On prolonge un circuit de prises : on utilise du \textbf{3G2,5} (section \SI{2,5}{mm^2}, avec conducteur de terre). 2) Nouvelle longueur : $11 + 2 = \SI{13}{m}$ de 3G2,5. 3) Oui : la norme autorise jusqu'à 8 prises sur un circuit en \SI{2,5}{mm^2} protégé par un disjoncteur 16 A ; avec 3 prises, on reste largement dans les limites.
\end{corrige}

\begin{aretenir}[L'essentiel de la section]
\begin{itemize}
\item Les symboles normalisés (NF EN 60617) sont le langage commun de tous les électriciens.
\item Le \textbf{schéma unifilaire} (un trait par circuit) sert au dimensionnement : sections, calibres, longueurs, mode de pose.
\item Le \textbf{schéma multifilaire} montre tous les fils et toutes les connexions : c'est le document du câblage réel.
\item Le \textbf{plan d'implantation} fixe les positions et les hauteurs : interrupteurs à \SI{1,10}{m}, prises à \SI{0,30}{m} (\SI{0,90}{m} en cuisine), tableau entre \SI{1,30}{m} et \SI{1,60}{m}.
\item Va-et-vient : phase sur le commun du premier, lampe sur le commun du second, deux navettes entre eux, neutre direct à la lampe.
\item Avant toute mise sous tension : continuité de terre, isolement $> \SI{0,5}{M\ohm}$, polarité, serrage, étiquetage.
\end{itemize}
\end{aretenir}

\newpage

\coursec{Activité d'intégration}

\courssub{Objectifs de l'activité}

À la fin de cette activité d'intégration, tu seras capable de :

\begin{itemize}
\item identifier les conducteurs d'une installation domestique réelle (phase, neutre, terre) à l'aide du code couleur et du tournevis testeur ;
\item établir un constat de danger sur une installation vétuste et proposer une solution conforme ;
\item choisir un disjoncteur divisionnaire et une section de câble adaptés à un circuit de prises ;
\item dessiner le schéma multifilaire d'un circuit domestique simple en respectant le code couleur ;
\item appliquer la procédure de Vérification d'Absence de Tension (VAT) avant toute intervention ;
\item justifier, à l'oral, des choix techniques par des arguments de sécurité.
\end{itemize}

\courssub{Situation}

\textbf{Rénovation d'une salle multimédia au collège de Nkol-Eton.}

Monsieur Abena, enseignant d'informatique au collège public de Nkol-Eton (région du Centre), veut installer un poste de travail dans une ancienne salle de classe. Le matériel à brancher est le suivant :

\begin{itemize}
\item une unité centrale d'ordinateur, appareil de \cle{classe I} (carcasse métallique, nécessite la terre), puissance $P_1 = \SI{350}{W}$ ;
\item un écran, appareil de \cle{classe II} (double isolation), puissance $P_2 = \SI{60}{W}$ ;
\item une imprimante, appareil de \cle{classe I}, puissance $P_3 = \SI{400}{W}$ ;
\item une lampe de bureau, puissance $P_4 = \SI{40}{W}$, que l'on souhaite commander par un interrupteur.
\end{itemize}

Dans la salle, il n'existe qu'\textbf{une seule prise murale}, installée dans les années 1980. En ouvrant cette prise, on découvre trois fils : un fil \textbf{rouge}, un fil \textbf{noir} et un fil \textbf{blanc}. Aucun fil vert/jaune n'est visible. Le tableau de répartition du bâtiment, alimenté par le compteur ENEO ($U = \SI{220}{V}$, $f = \SI{50}{Hz}$), possède un \textbf{interrupteur différentiel général de \SI{30}{mA}} et \textbf{deux emplacements libres} pour des disjoncteurs divisionnaires.

\textbf{Documents fournis :}
\begin{itemize}
\item \textbf{Doc 1 :} photo de la prise existante ouverte (fils rouge, noir, blanc ; aucune borne de terre).
\item \textbf{Doc 2 :} extrait de catalogue de disjoncteurs divisionnaires : calibres disponibles 10 A, 16 A, 20 A ; courbes B, C, D ; prix de 3 500 à 6 000 FCFA.
\item \textbf{Doc 3 :} tableau de correspondance section/calibre : circuit prises en câble 3G2,5 (trois conducteurs dont la terre, section \SI{2,5}{mm^2}) protégé par disjoncteur 16 A maximum ; circuit éclairage en 3G1,5 protégé par 10 A.
\end{itemize}

\textbf{Problème à résoudre :} comment alimenter ce poste de travail en toute sécurité, sachant que la prise existante ne possède pas de terre ?

\courssub{Consignes}

Travail en binôme, durée 30 minutes. Rédige tes réponses sur la fiche de synthèse.

\begin{enumerate}
\item \textbf{Analyse.} À l'aide du tournevis testeur, explique comment identifier le fil de phase parmi les trois fils de la prise ouverte. À quels conducteurs correspondent les fils rouge, noir et blanc selon l'ancien code couleur ? Que déduis-tu de l'absence de fil vert/jaune ? Rédige un constat de danger en trois phrases maximum.
\item \textbf{Conception.} Calcule la puissance totale du poste de travail, puis l'intensité totale absorbée sous \SI{220}{V}. Choisis dans le Doc 2 le calibre du disjoncteur divisionnaire du nouveau circuit prises et justifie. Pourquoi choisit-on le câble 3G2,5 ?
\item \textbf{Schéma.} Dessine le schéma multifilaire du nouveau circuit : départ tableau (différentiel \SI{30}{mA} puis disjoncteur divisionnaire 16 A), boîte de dérivation, deux prises 2P+T et une prise commandée par un interrupteur simple pour la lampe. Respecte le code couleur normalisé.
\item \textbf{Sécurité.} Liste les cinq étapes de la VAT avant le raccordement au tableau. Explique, par au moins trois arguments (section, protection, terre), pourquoi il est interdit de simplement brancher une rallonge ou une prise en dérivation sur l'ancienne prise.
\item \textbf{Restitution.} Prépare une présentation orale de 2 minutes : « Notre solution technique et nos arguments de sécurité ».
\end{enumerate}

\courssub{Ressources à mobiliser}

\begin{aretenir}[Boîte à outils de l'activité]
\begin{itemize}
\item \textbf{Code couleur normalisé :} phase = rouge (ou marron, noir) ; neutre = bleu clair ; terre = vert/jaune. Dans les anciennes installations, le blanc ou le noir pouvait servir de neutre : d'où l'intérêt du testeur.
\item \textbf{Tournevis testeur :} il s'allume au contact de la phase (le courant très faible traverse le corps vers la terre) ; il reste éteint sur le neutre et sur la terre.
\item \textbf{Puissance et intensité :} $P = U \times I$, donc $I = \dfrac{P}{U}$.
\item \textbf{Prise 2P+T :} phase à droite (vue de face), neutre à gauche, terre sur la broche ; la terre évacue les courants de défaut des appareils de classe I.
\item \textbf{Rôle du différentiel \SI{30}{mA} :} il coupe le circuit dès qu'une fuite de courant supérieure à \SI{30}{mA} est détectée (protection des personnes).
\item \textbf{VAT :} Vérification d'Absence de Tension en 5 étapes.
\end{itemize}
\end{aretenir}

\courssub{Démarche et corrigé commenté}

\begin{exoresolu}[Corrigé commenté de l'activité]
\textbf{1. Analyse de la prise existante.} On met le circuit sous tension, puis on touche successivement chaque fil avec la pointe du tournevis testeur, le doigt sur le capuchon métallique. Le fil \textbf{rouge} fait briller le néon du testeur : c'est la \cle{phase}. Les fils noir et blanc ne l'allument pas : l'un est le \cle{neutre}, l'autre éventuellement un fil de retour ou un neutre repéré. Aucun fil vert/jaune n'est présent et la prise ne comporte pas de broche de terre. \textbf{Constat de danger :} l'installation est vétuste, sans conducteur de protection ; brancher un appareil de classe I (unité centrale, imprimante) y expose l'utilisateur au risque d'électrisation en cas de défaut d'isolement, car la carcasse métallique ne peut pas être mise à la terre.
\tcblower
\textbf{2. Conception.} Puissance totale :
\[ P = P_1 + P_2 + P_3 + P_4 = 350 + 60 + 400 + 40 = \SI{850}{W}. \]
Intensité totale :
\[ I = \frac{P}{U} = \frac{850}{220} \approx \SI{3,9}{A}. \]
Même avec une marge d'évolution (ajout futur d'un deuxième poste), un disjoncteur divisionnaire de \textbf{16 A, courbe C} convient : $16 > 3{,}9$ et c'est le calibre normalisé d'un circuit de prises. Le câble \textbf{3G2,5} (phase, neutre, terre, section \SI{2,5}{mm^2}) est imposé par le Doc 3 pour un circuit prises protégé à 16 A : il échauffe peu et véhicule le conducteur de terre indispensable aux appareils de classe I.

\textbf{3. Schéma multifilaire.} Depuis le tableau : le différentiel \SI{30}{mA} alimente le disjoncteur 16 A ; de là partent trois fils — phase (rouge), neutre (bleu), terre (vert/jaune) — vers une boîte de dérivation. De la boîte : deux départs en 3G2,5 vers les deux prises 2P+T (phase à droite, neutre à gauche, terre au centre) ; un départ vers l'interrupteur simple qui \textbf{coupe la phase} (jamais le neutre) avant la prise commandée de la lampe ; neutre et terre vont directement à cette prise.

\textbf{4. Sécurité.} Les 5 étapes de la VAT : (1) couper le disjoncteur divisionnaire et consigner ; (2) vérifier que le testeur fonctionne sur une source sous tension ; (3) vérifier l'absence de tension entre phase et neutre puis phase et terre sur le circuit ; (4) revérifier le bon fonctionnement du testeur ; (5) constater l'absence de tension et commencer le travail. On ne peut pas dériver l'ancienne prise car : sa section de fil est inconnue et probablement insuffisante ; elle n'est pas protégée par un disjoncteur divisionnaire adapté ; elle ne possède pas de terre, donc aucune protection des appareils de classe I.

\textbf{5. Restitution.} Exemple d'argumentaire : « Nous créons un circuit neuf dédié, protégé par un disjoncteur 16 A et le différentiel 30 mA, câblé en 3G2,5 avec terre sur toutes les prises ; l'ancienne prise est déposée. Ainsi, les personnes sont protégées contre les contacts indirects et le matériel contre les surintensités. »
\end{exoresolu}

\begin{attention}[Pièges fréquents]
Ne confonds jamais le rôle du disjoncteur (protège contre les surintensités) et celui du différentiel (protège les personnes contre les fuites de courant). L'interrupteur se place toujours sur la phase. Le tournevis testeur indique la présence de la phase, mais ne mesure pas la tension : ce n'est pas un voltmètre.
\end{attention}

\courssub{Bilan de l'activité}

Cette activité t'a permis de mobiliser en situation réelle l'ensemble des savoirs de la leçon : reconnaître les conducteurs par le code couleur et le testeur, comprendre le rôle vital de la terre pour les appareils de classe I, dimensionner un circuit (calcul de $I = P/U$, choix du calibre 16 A et de la section \SI{2,5}{mm^2}), traduire une installation en schéma multifilaire et appliquer les règles de sécurité (VAT, différentiel \SI{30}{mA}). Tu as aussi constaté qu'une installation vétuste sans terre, fréquente dans les bâtiments anciens au Cameroun, impose une rénovation complète du circuit plutôt qu'un bricolage dangereux : la sécurité électrique se conçoit, elle ne s'improvise pas.

\newpage

\coursec{Méthodes et exercices résolus}

\coursec{Leçon 21 : Généralités sur les circuits électriques domestiques}

\courssub{Méthode 1 : Identifier les éléments d'une installation domestique}

\begin{methode}[Identifier et décrire une installation électrique domestique]
Pour analyser l'installation électrique d'une maison, suis toujours le trajet du courant, de l'arrivée ENEO jusqu'aux appareils :
\begin{enumerate}
\item \textbf{Repère le point d'entrée} : le branchement (câble venant du poteau ou du compteur) et le \cle{compteur} qui mesure l'énergie consommée (en \SI{}{kWh}).
\item \textbf{Identifie l'organe de coupure générale} : le \cle{disjoncteur général} (ou disjoncteur de branchement), qui protège toute l'installation et permet de couper tout le courant.
\item \textbf{Décris le tableau de répartition} : il contient les \cle{disjoncteurs divisionnaires} (ou fusibles) qui répartissent le courant en plusieurs circuits (prises, éclairage, chauffe-eau...).
\item \textbf{Nomme les circuits terminaux} : circuits d'éclairage, circuits de prises, circuits spécialisés (cuisinière, climatiseur).
\item \textbf{Vérifie la présence de la prise de terre} : piquet de terre relié au tableau par le conducteur vert/jaune.
\end{enumerate}
\textbf{Réflexe rédaction} : cite toujours le rôle de chaque élément (mesurer, couper, protéger, répartir), pas seulement son nom.
\end{methode}

\begin{exoresolu}[L'installation de la famille Ngo Bell]
La famille Ngo Bell vient de construire une maison à Yaoundé. En partant du poteau ENEO, on observe dans l'ordre : un câble de branchement, un compteur, un disjoncteur général de \SI{30}{A}, puis un tableau comportant 4 disjoncteurs divisionnaires alimentant respectivement les prises du salon, l'éclairage, la cuisine et la douche. Un piquet de terre est planté dans la cour.
\begin{enumerate}
\item Fais le schéma simplifié (unifilaire) de cette installation.
\item Donne le rôle du compteur et du disjoncteur général.
\item Pourquoi le circuit de la douche possède-t-il son propre disjoncteur divisionnaire ?
\end{enumerate}
\tcblower
\textbf{1. Schéma unifilaire.} On représente les éléments en chaîne, dans l'ordre du trajet du courant :
\begin{popfigure}
\begin{tikzpicture}[font=\small, >=Stealth]
\node[draw, rounded corners, fill=popblueL, minimum width=2.2cm, minimum height=0.9cm] (poteau) at (0,0) {Poteau ENEO};
\node[draw, rounded corners, fill=popblueL, minimum width=1.8cm, minimum height=0.9cm, right=0.7cm of poteau] (compteur) {Compteur};
\node[draw, rounded corners, fill=poporangeL, minimum width=2.2cm, minimum height=0.9cm, right=0.7cm of compteur] (dg) {Disjoncteur général};
\node[draw, rounded corners, fill=popgreenL, minimum width=2.6cm, minimum height=0.9cm, right=0.7cm of dg] (tableau) {Tableau de répartition};
\draw[->, thick] (poteau) -- (compteur);
\draw[->, thick] (compteur) -- (dg);
\draw[->, thick] (dg) -- (tableau);
\foreach \k/\y/\nom in {1/1.2/{Prises salon}, 2/0.4/{Éclairage}, 3/-0.4/{Cuisine}, 4/-1.2/{Douche}}{
\node[draw, fill=popcream, minimum width=2cm, minimum height=0.5cm, right=1.2cm of tableau, yshift=\y cm, inner sep=2pt] (cc\k) {\footnotesize \nom};
\draw[->, thick] (tableau.east) -- (cc\k.west);}
\node[below=0.15cm of tableau, text=popgreen] {\footnotesize Relié au piquet de terre};
\end{tikzpicture}
\legende{Schéma unifilaire de l'installation de la famille Ngo Bell.}
\end{popfigure}
\textbf{2. Rôles.}
\begin{itemize}
\item Le \cle{compteur} \textbf{mesure l'énergie électrique} consommée par l'habitation, en kilowattheures (\SI{}{kWh}) : c'est sur cette mesure qu'ENEO établit la facture.
\item Le \cle{disjoncteur général} \textbf{coupe automatiquement} tout le courant de la maison en cas de surcharge ou de court-circuit ; il permet aussi de couper volontairement l'alimentation avant toute intervention.
\end{itemize}
\textbf{3. Circuit de la douche.} La douche est une pièce humide : le risque d'électrisation y est plus élevé car l'eau diminue la résistance du corps humain. Un disjoncteur divisionnaire propre à ce circuit permet de le protéger et de le couper séparément, sans priver toute la maison de courant.
\textbf{Conclusion} : une installation domestique s'organise toujours selon la chaîne : branchement $\to$ compteur $\to$ disjoncteur général $\to$ tableau de répartition $\to$ circuits divisionnaires, avec une prise de terre pour la sécurité des personnes.
\end{exoresolu}

\courssub{Méthode 2 : Identifier les conducteurs grâce au code couleur}

\begin{methode}[Reconnaître phase, neutre et terre]
Pour identifier les conducteurs d'une installation :
\begin{enumerate}
\item \textbf{Mémorise le code couleur normalisé} :
\begin{itemize}
\item \cle{Phase} (fil qui « apporte » le courant, dangereux) : rouge, marron ou noir ;
\item \cle{Neutre} (fil de retour du courant) : bleu clair, obligatoirement ;
\item \cle{Terre} (fil de protection) : bicolore vert/jaune, obligatoirement et exclusivement.
\end{itemize}
\item \textbf{En cas de doute, vérifie avec un tournevis testeur} : il s'allume sur la phase, pas sur le neutre ni sur la terre.
\item \textbf{Retiens les tensions} : entre phase et neutre : $U = \SI{220}{V}$ ; entre phase et terre : $U = \SI{220}{V}$ ; entre neutre et terre : $U \approx \SI{0}{V}$.
\item \textbf{Sections des fils} (à connaître) : éclairage : \SI{1,5}{mm^2} ; prises de courant : \SI{2,5}{mm^2} ; gros appareils (cuisinière) : \SI{4}{mm^2} ou \SI{6}{mm^2}.
\end{enumerate}
\textbf{Interdit absolu} : ne jamais utiliser le fil vert/jaune comme phase ou neutre, même en dépannage.
\end{methode}

\begin{exoresolu}[Le chantier de l'électricien]
Un électricien ouvre une boîte de dérivation dans une maison à Douala. Il y trouve trois fils : un fil bleu, un fil rouge et un fil bicolore vert/jaune.
\begin{enumerate}
\item Identifie chaque fil et précise son rôle.
\item Il mesure la tension entre le fil rouge et le fil bleu : il trouve \SI{220}{V}. Ce résultat est-il cohérent ? Justifie.
\item Son apprenti propose de remplacer un fil de phase cassé par un morceau de fil vert/jaune « puisque c'est le même cuivre ». Que répondre ?
\end{enumerate}
\tcblower
\textbf{1. Identification des fils.}
\begin{itemize}
\item Le fil \textbf{rouge} est la \cle{phase} : c'est le conducteur qui transporte le courant depuis le réseau ENEO ; il présente une tension de \SI{220}{V} par rapport à la terre et est dangereux au toucher.
\item Le fil \textbf{bleu} est le \cle{neutre} : c'est le conducteur de retour du courant ; sa tension par rapport à la terre est voisine de \SI{0}{V}.
\item Le fil \textbf{vert/jaune} est le \cle{conducteur de terre} : il relie les carcasses métalliques des appareils au piquet de terre et évacue les courants de défaut pour protéger les personnes.
\end{itemize}
\textbf{2. Cohérence de la mesure.} La tension mesurée entre la phase (rouge) et le neutre (bleu) doit être égale à la tension du réseau domestique, soit environ \SI{220}{V}. La valeur trouvée, \SI{220}{V}, est donc \textbf{parfaitement cohérente} avec l'identification faite à la question 1.
\textbf{3. Réponse à l'apprenti.} Cette proposition est \textbf{interdite et dangereuse}. Le vert/jaune est réservé exclusivement à la terre : si on l'utilise comme phase, toute personne qui verra ce fil croira qu'il s'agit d'un fil de terre sans danger et pourra le toucher ou le raccorder à une carcasse métallique. La carcasse se retrouverait alors portée à \SI{220}{V} : risque mortel d'électrocution.
\textbf{Conclusion} : le code couleur (phase : rouge/marron/noir ; neutre : bleu ; terre : vert/jaune) est une règle de sécurité qu'aucun dépannage ne justifie de transgresser.
\end{exoresolu}

\courssub{Méthode 3 : Brancher correctement une prise de courant 2P+T}

\begin{methode}[Câbler une prise 2P+T]
Pour brancher ou vérifier une prise de courant à deux pôles plus terre :
\begin{enumerate}
\item \textbf{Coupe d'abord le courant} au disjoncteur général : on ne travaille jamais sur une prise sous tension.
\item \textbf{Repère les trois bornes} : borne de phase, borne de neutre, et borne de terre (reliée aux deux languettes métalliques latérales ou à la broche de terre).
\item \textbf{Raccorde} : le fil de phase (rouge/marron/noir) sur la borne de phase ; le fil bleu sur la borne de neutre ; le fil vert/jaune sur la borne de terre.
\item \textbf{Vérifie le serrage} des vis : un fil mal serré chauffe et peut provoquer un incendie.
\item \textbf{Explique le rôle de la terre} : en cas de défaut d'isolement (un fil de phase touche la carcasse métallique d'un appareil), le courant de défaut s'écoule vers la terre au lieu de traverser le corps de l'utilisateur, et le disjoncteur différentiel coupe le circuit.
\end{enumerate}
\textbf{Formule à retenir} : la prise 2P+T comporte 2 pôles (phase + neutre) pour l'alimentation et 1 contact de \cle{terre} pour la protection.
\end{methode}

\begin{exoresolu}[La prise de la cuisinière]
Maman Achu achète une cuisinière électrique dont la carcasse est métallique. La fiche du câble d'alimentation comporte trois conducteurs.
\begin{enumerate}
\item Pourquoi la fiche comporte-t-elle trois conducteurs alors que deux suffisent pour faire fonctionner la cuisinière ?
\item Par accident, le fil de phase touche la carcasse métallique à l'intérieur de l'appareil. Décris ce qui se passe si la prise de terre est correctement branchée, puis si elle ne l'est pas.
\end{enumerate}
\tcblower
\textbf{1. Utilité du troisième conducteur.} Deux conducteurs (phase et neutre) suffisent effectivement à alimenter la cuisinière en courant. Le troisième conducteur, le fil de \cle{terre} (vert/jaune), ne sert pas au fonctionnement : il relie la carcasse métallique au piquet de terre. C'est un organe de \textbf{sécurité} indispensable pour tout appareil à carcasse métallique.
\textbf{2. Conséquences du défaut.}
\begin{itemize}
\item \textbf{Terre branchée} : le courant de défaut s'écoule de la carcasse vers la terre par le fil vert/jaune, qui offre un chemin de faible résistance. Ce courant anormal est détecté par le disjoncteur (différentiel), qui \textbf{coupe automatiquement} le circuit. L'utilisateur qui touche la carcasse ne risque rien.
\item \textbf{Terre non branchée} : la carcasse reste portée au potentiel de la phase, soit \SI{220}{V}. L'utilisateur qui la touche ferme un circuit entre la phase et le sol à travers son propre corps : il est traversé par un courant dangereux, c'est l'\cle{électrisation}, qui peut être mortelle (électrocution).
\end{itemize}
\textbf{Conclusion} : la prise de terre ne sert pas au fonctionnement de l'appareil mais à la protection des personnes : elle évacue les courants de défaut et provoque la coupure du circuit.
\end{exoresolu}

\courssub{Méthode 4 : Utiliser et interpréter le tournevis testeur}

\begin{methode}[Tester une borne avec le tournevis testeur]
\begin{enumerate}
\item \textbf{Comprends le principe} : le tournevis testeur contient une lampe néon et une forte résistance en série. Quand la pointe touche la phase et que le doigt touche le capuchon métallique, un courant très faible traverse le corps vers la terre : la lampe s'allume.
\item \textbf{Pose la pointe} sur la borne à tester et \textbf{touche le capuchon métallique} avec un doigt (sinon le circuit n'est pas fermé et la lampe ne s'allume jamais).
\item \textbf{Interprète} : lampe allumée $\Rightarrow$ la borne est la \cle{phase} ; lampe éteinte $\Rightarrow$ borne de neutre, de terre, ou circuit hors tension.
\item \textbf{Connais les limites} : le testeur ne mesure pas la tension exacte ; s'il ne s'allume pas, on ne peut pas affirmer qu'il n'y a pas de tension (lampe grillée, mauvais contact du doigt). Avant un travail important, on coupe toujours au disjoncteur général.
\end{enumerate}
\textbf{Sécurité} : la résistance interne limite le courant à une valeur inoffensive ; n'utilise jamais un testeur endommagé.
\end{methode}

\begin{exoresolu}[Le test d'une prise]
Pour identifier les bornes d'une prise murale, Etienne utilise un tournevis testeur. En touchant le trou de gauche, la lampe s'allume ; en touchant le trou de droite, elle reste éteinte.
\begin{enumerate}
\item Identifie la borne de phase et la borne de neutre. Justifie.
\item Etienne conclut : « La lampe est éteinte sur le trou de droite, donc je peux toucher ce fil sans danger. » Cette conclusion est-elle acceptable ? Justifie.
\end{enumerate}
\tcblower
\textbf{1. Identification.} Le tournevis testeur s'allume lorsque sa pointe touche un conducteur porté à une tension élevée par rapport à la terre, car un faible courant traverse alors la lampe néon, la résistance interne et le corps de l'utilisateur vers la terre. Le trou de gauche, qui allume la lampe, est donc relié à la \cle{phase}. Le trou de droite, qui n'allume pas la lampe, est relié au \cle{neutre} (tension voisine de \SI{0}{V} par rapport à la terre).
\textbf{2. Critique de la conclusion.} Cette conclusion est \textbf{inacceptable}. Le fait que la lampe reste éteinte ne prouve pas l'absence de danger : la lampe néon peut être grillée, le doigt peut mal toucher le capuchon, ou le neutre peut être accidentellement sous tension à la suite d'un défaut de l'installation. De plus, toucher un conducteur, même le neutre, avec les mains ou en milieu humide reste imprudent.
\textbf{Conclusion} : le tournevis testeur permet d'identifier la phase, mais une lampe éteinte n'est jamais une garantie d'absence de tension ; la seule consigne sûre est de couper le disjoncteur général avant toute intervention.
\end{exoresolu}

\courssub{Méthode 5 : Appliquer les règles de sécurité électrique}

\begin{methode}[Prévenir les dangers et secourir]
\textbf{A. Dangers du courant} (à citer précisément) :
\begin{itemize}
\item \cle{Électrisation} : passage du courant à travers le corps (brûlures, arrêt cardiaque) ; le danger augmente avec l'humidité.
\item \cle{Court-circuit} : contact direct phase-neutre $\Rightarrow$ échauffement brutal, étincelles, incendie.
\item \cle{Surcharge} : trop d'appareils sur une même prise ou rallonge $\Rightarrow$ échauffement des fils.
\end{itemize}
\textbf{B. Protections} : disjoncteurs et fusibles (contre surcharges et courts-circuits), prise de terre et disjoncteur différentiel (contre l'électrisation), isolation des fils, prises avec obturateurs.
\textbf{C. Premiers secours} (ordre impératif) :
\begin{enumerate}
\item \textbf{Couper le courant} (disjoncteur général) avant tout contact avec la victime ;
\item si impossible, repousser la victime avec un objet \textbf{isolant et sec} (bois sec, bâton plastique) — jamais à mains nues ;
\item alerter les secours (pompiers : 118 au Cameroun) ;
\item si la victime ne respire pas, commencer les gestes de réanimation si l'on est formé.
\end{enumerate}
\end{methode}

\begin{exoresolu}[Accident dans la cour]
Pendant la saison des pluies à Bafoussam, Junior branche une rallonge sur une prise extérieure pour utiliser une meuleuse. Le sol est mouillé. Soudain, son petit frère, pieds nus dans l'eau, touche la carcasse métallique de la meuleuse dont l'isolement est défectueux : il reste « collé » à l'appareil et ne peut plus le lâcher.
\begin{enumerate}
\item Cite deux facteurs qui ont aggravé le danger dans cette situation.
\item Décris, dans l'ordre, les gestes que Junior doit faire pour secourir son frère.
\item Quelle protection aurait évité cet accident ?
\end{enumerate}
\tcblower
\textbf{1. Facteurs aggravants.}
\begin{itemize}
\item Le \textbf{sol mouillé et les pieds nus} : l'eau diminue fortement la résistance du corps humain, donc le courant qui traverse la victime est plus intense ;
\item le \textbf{défaut d'isolement} de la meuleuse : sa carcasse métallique est portée au potentiel de la phase (\SI{220}{V}) ;
\item l'utilisation d'un appareil à carcasse métallique \textbf{en extérieur humide}, milieu à haut risque.
\end{itemize}
\textbf{2. Gestes de secours (dans l'ordre).}
\begin{enumerate}
\item Junior doit d'abord \textbf{couper le courant} au disjoncteur général ou débrancher la rallonge sans toucher la victime ;
\item s'il ne peut pas couper, il repousse son frère avec un objet \textbf{sec et isolant} (manche en bois sec), jamais à mains nues, sinon il sera électrisé à son tour ;
\item il \textbf{alerte les secours} (118) ;
\item il surveille la respiration de la victime et, s'il est formé, pratique les gestes de premiers secours en attendant les pompiers.
\end{enumerate}
\textbf{3. Protection préventive.} Une \textbf{prise de terre correctement raccordée}, associée à un \textbf{disjoncteur différentiel}, aurait évacué le courant de défaut vers la terre et coupé le circuit dès l'apparition du défaut, avant même que l'enfant ne touche la carcasse.
\textbf{Conclusion} : l'eau et les défauts d'isolement multiplient les dangers du courant ; la règle d'or du secours est : couper le courant avant de toucher la victime.
\end{exoresolu}

\courssub{Méthode 6 : Lire un schéma unifilaire et calculer la puissance d'un circuit}

\begin{methode}[Exploiter un schéma et vérifier une surcharge]
\begin{enumerate}
\item \textbf{Lis le schéma unifilaire} : chaque trait représente un circuit (phase + neutre + terre) partant d'un disjoncteur divisionnaire ; identifie ce que chaque circuit alimente.
\item \textbf{Calcule la puissance totale} branchée sur un circuit : $P_{\text{totale}} = P_1 + P_2 + \dots$
\item \textbf{Calcule l'intensité} correspondante avec la formule de la puissance :
\[ P = U \times I \quad \Longleftrightarrow \quad I = \dfrac{P}{U} \]
avec $P$ en watts (\SI{}{W}), $U$ en volts (\SI{}{V}), $I$ en ampères (\SI{}{A}).
\item \textbf{Compare} au calibre du disjoncteur : si $I > I_{\text{disjoncteur}}$, il y a surcharge : le disjoncteur saute. Il faut alors débrancher des appareils ou répartir sur un autre circuit.
\end{enumerate}
\textbf{Réflexe} : convertis toujours les kilowatts en watts ($\SI{1}{kW} = \SI{1000}{W}$) avant le calcul.
\end{methode}

\begin{exoresolu}[La prise multiple du salon]
Sur une prise du salon protégée par un disjoncteur divisionnaire de \SI{16}{A} (tension du réseau : $U = \SI{220}{V}$), la famille Fotso branche simultanément sur une multiprise : un téléviseur de \SI{150}{W}, un décodeur de \SI{30}{W}, un ventilateur de \SI{70}{W} et un fer à repasser de \SI{1200}{W}.
\begin{enumerate}
\item Calcule la puissance totale utilisée.
\item Calcule l'intensité totale du courant qui traverse la multiprise.
\item Le disjoncteur va-t-il se déclencher ? On ajoute ensuite une bouilloire de \SI{2000}{W}. Que se passe-t-il ? Conclus.
\end{enumerate}
\tcblower
\textbf{1. Puissance totale.} Les appareils branchés en parallèle sur la même prise additionnent leurs puissances :
\[ P_{\text{totale}} = 150 + 30 + 70 + 1200 = \SI{1450}{W}. \]
\textbf{2. Intensité totale.} D'après la relation $P = U \times I$, on tire :
\[ I = \dfrac{P_{\text{totale}}}{U} = \dfrac{1450}{220} \approx \SI{6,6}{A}. \]
\textbf{3. Comparaison au calibre du disjoncteur.} On a $I \approx \SI{6,6}{A} < \SI{16}{A}$ : l'intensité est inférieure au calibre du disjoncteur, donc celui-ci \textbf{ne se déclenche pas} ; l'installation fonctionne normalement.
\textbf{Ajout de la bouilloire.} La nouvelle puissance totale vaut :
\[ P'_{\text{totale}} = 1450 + 2000 = \SI{3450}{W}. \]
La nouvelle intensité est :
\[ I' = \dfrac{3450}{220} \approx \SI{15,7}{A}. \]
On a $I' \approx \SI{15,7}{A} < \SI{16}{A}$ : le disjoncteur tient encore, mais on est \textbf{très proche de la limite}. Si l'on ajoutait encore un appareil (par exemple un radiateur de \SI{500}{W}, soit $I'' = \dfrac{3950}{220} \approx \SI{18}{A} > \SI{16}{A}$), le disjoncteur \textbf{sauterait} : c'est une surcharge.
\textbf{Conclusion} : avant de brancher plusieurs appareils sur une multiprise, on additionne les puissances, on calcule $I = \dfrac{P}{U}$ et on vérifie que l'intensité reste inférieure au calibre du disjoncteur divisionnaire. La surcharge fait sauter le disjoncteur et peut échauffer dangereusement les fils de la rallonge.
\end{exoresolu}

\begin{attention}[Les 5 erreurs qui coûtent des points au BEPC]
\begin{enumerate}
\item \textbf{Confondre les rôles du compteur et du disjoncteur.} Le compteur \emph{mesure} l'énergie (en \SI{}{kWh}) ; le disjoncteur \emph{coupe et protège}. Écrire « le compteur protège l'installation » est une erreur éliminatoire.
\item \textbf{Inverser les couleurs des conducteurs.} Neutre = bleu, terre = vert/jaune, phase = rouge/marron/noir. Dire que le fil de terre est bleu, ou utiliser le vert/jaune comme phase, est une faute grave car c'est une règle de sécurité.
\item \textbf{Affirmer que le tournevis testeur « mesure la tension ».} Il ne fait que \emph{détecter} la présence de la phase. Et une lampe éteinte ne prouve jamais l'absence de tension : ne jamais conclure « on peut toucher sans danger ».
\item \textbf{Oublier de convertir les unités dans $P = U \times I$.} Les kilowatts doivent être convertis en watts ($\SI{1}{kW} = \SI{1000}{W}$) avant de calculer $I = \dfrac{P}{U}$ ; sinon l'intensité est 1000 fois trop faible et la conclusion sur la surcharge est fausse.
\item \textbf{Décrire les premiers secours dans le mauvais ordre.} Toucher la victime avant d'avoir coupé le courant est la première chose à proscrire : le secouriste serait électrisé à son tour. La réponse attendue commence toujours par « couper le courant au disjoncteur général ».
\end{enumerate}
\end{attention}

\newpage

\coursec{Exercices}

\courssub{Je vérifie mes connaissances}

\begin{exercice}[QCM : le branchement dangereux]
Une prise 2P+T est branchée avec les fils suivants : fil bleu sur la borne L, fil rouge sur la borne N, fil vert/jaune sur la borne Terre. L'appareil raccordé (Classe I) fonctionne normalement. Quel est le danger principal ?
\begin{enumerate}
\item[A)] L'appareil ne peut pas fonctionner.
\item[B)] L'interrupteur de l'appareil coupe le Neutre : la Phase reste présente sur les composants internes de l'appareil.
\item[C)] La Terre n'est pas connectée.
\item[D)] Le différentiel ne déclenchera jamais.
\end{enumerate}
\emph{Une seule réponse exacte. Justifier brièvement votre choix.}
\end{exercice}

\begin{exercice}[Vrai ou faux]
Pour chaque affirmation, répondre par VRAI ou FAUX, puis justifier en une phrase.
\begin{enumerate}
\item Le fil de phase d'une installation domestique est toujours de couleur bleue.
\item Le tournevis testeur permet de distinguer avec certitude le Neutre de la Terre.
\item Le disjoncteur différentiel de \SI{30}{mA} protège les personnes contre les contacts indirects.
\item Couper l'interrupteur d'une lampe suffit pour travailler en sécurité sur une prise de la même pièce.
\end{enumerate}
\end{exercice}

\begin{exercice}[Définitions et code couleur]
\begin{enumerate}
\item Définir en une phrase chacun des termes suivants : \cle{disjoncteur divisionnaire}, \cle{conducteur de protection (terre)}, \cle{schéma unifilaire}.
\item Compléter le tableau ci-dessous avec la couleur normalisée de chaque conducteur.
\end{enumerate}
\begin{center}
\begin{tabularx}{\linewidth}{|l|X|}
\hline
\rowcolor{popblueL} \textbf{Conducteur} & \textbf{Couleur normalisée} \\
\hline
Phase (L) & \ldots \\
\hline
Neutre (N) & \ldots \\
\hline
Terre (PE) & \ldots \\
\hline
\end{tabularx}
\end{center}
\end{exercice}

\begin{exercice}[Calcul direct]
Une cuisinière électrique de puissance $P = \SI{4600}{W}$ est branchée sous la tension du réseau $U = \SI{230}{V}$.
\begin{enumerate}
\item Calculer l'intensité $I$ du courant absorbé.
\item Cette cuisinière peut-elle être protégée par un disjoncteur divisionnaire de \SI{20}{A} ? Justifier.
\end{enumerate}
\end{exercice}

\courssub{Je m'entraîne}

\begin{exercice}[Identifier les éléments d'une installation]
Sur la facture ENEO et le tableau électrique de sa maison, Martial relève les mentions suivantes : compteur \SI{45}{A}, disjoncteur d'abonné, interrupteur différentiel \SI{30}{mA}, disjoncteurs divisionnaires \SI{10}{A}, \SI{16}{A} et \SI{32}{A}.
\begin{enumerate}
\item Classer ces éléments dans l'ordre où le courant les traverse, de l'arrivée du réseau vers les prises.
\item Donner le rôle précis de chacun de ces quatre types d'éléments.
\item Quel élément protège les personnes ? Quels éléments protègent les circuits et les appareils ?
\end{enumerate}
\end{exercice}

\begin{exercice}[Sections de câbles et calibres]
Recopier et compléter le tableau d'association suivant, en utilisant les valeurs du cours.
\begin{center}
\begin{tabularx}{\linewidth}{|X|l|l|}
\hline
\rowcolor{popblueL} \textbf{Circuit} & \textbf{Calibre du disjoncteur} & \textbf{Section minimale} \\
\hline
Éclairage & \ldots & \SI{1,5}{mm^2} \\
\hline
Prises de courant ordinaires & \SI{16}{A} & \ldots \\
\hline
Four ou cuisinière & \ldots & \SI{6}{mm^2} \\
\hline
\end{tabularx}
\end{center}
Expliquer pourquoi on écrit « 3G1,5 » pour un câble d'éclairage : que signifient le chiffre 3, la lettre G et le nombre 1,5 ?
\end{exercice}

\begin{exercice}[Utilisation du tournevis testeur]
Aminatou veut vérifier qu'une prise de sa chambre est correctement alimentée avant d'y brancher son groupe électrogène de secours (hors tension).
\begin{enumerate}
\item Décrire, étape par étape, la manipulation correcte du tournevis testeur sur cette prise (position du doigt, bornes à tester).
\item Sur la borne de Phase, le néon s'allume ; sur la borne de Neutre, il reste éteint. Conclure.
\item Le néon s'allume faiblement sur la borne de Terre. Peut-on en conclure que la Terre est en bon état ? Justifier la réponse.
\end{enumerate}
\end{exercice}

\begin{exercice}[Choix d'une protection]
Un atelier de menuiserie à Nkongsamba possède un circuit dédié à une scie électrique de puissance $P = \SI{3500}{W}$ sous $U = \SI{230}{V}$.
\begin{enumerate}
\item Calculer l'intensité absorbée en fonctionnement normal.
\item Choisir le calibre du disjoncteur divisionnaire parmi : \SI{10}{A}, \SI{16}{A}, \SI{20}{A}, \SI{32}{A}. Justifier le choix.
\item Le patron propose de « remonter » le calibre à \SI{32}{A} car le disjoncteur de \SI{16}{A} déclenche souvent. Expliquer pourquoi ce raisonnement est dangereux et proposer la bonne démarche.
\end{enumerate}
\end{exercice}

\courssub{Je me prépare au Bac}

\begin{exercice}[Type BEPC : le circuit four de la cuisine]
Dans une villa à Yaoundé, on installe un four électrique de puissance $P = \SI{7200}{W}$ sous la tension $U = \SI{230}{V}$. Le circuit est protégé par un disjoncteur divisionnaire de calibre \SI{32}{A} et câblé en \SI{6}{mm^2} sur une longueur de $L = \SI{20}{m}$ entre le tableau et le four. On donne la résistivité du cuivre : $\rho = \SI{1,7e-8}{\ohm.m}$ et la tension admissible de chute : $3\,\%$ de la tension réseau.
\begin{enumerate}
\item Calculer l'intensité $I$ absorbée par le four en fonctionnement.
\item Le disjoncteur de \SI{32}{A} est-il correctement choisi ? Justifier par une comparaison chiffrée.
\item Calculer la résistance $R$ du câble aller-retour (longueur totale $2L$), puis la chute de tension $\Delta U = R \times I$ dans le câble.
\item Exprimer $\Delta U$ en pourcentage de la tension réseau. L'installation est-elle conforme ? Conclure.
\item Expliquer pourquoi ce circuit doit obligatoirement comporter un conducteur de terre relié à la carcasse du four.
\end{enumerate}
\end{exercice}

\begin{exercice}[Type BEPC : schéma unifilaire d'une maison]
On donne le schéma de principe incomplet d'une installation domestique : Compteur \SI{45}{A} $\rightarrow$ Disjoncteur d'abonné $\rightarrow$ Interrupteur différentiel \SI{30}{mA} $\rightarrow$ trois départs divisionnaires (\SI{10}{A} éclairage, \SI{16}{A} prises, \SI{32}{A} four).
\begin{enumerate}
\item Reproduire et compléter ce schéma unifilaire en indiquant, pour chaque départ : le calibre du disjoncteur, la section du câble et la désignation normalisée du câble (par exemple 3G1,5).
\item Préciser, pour chaque départ, le nombre de conducteurs et leur code couleur.
\item Un défaut d'isolement apparaît sur le lave-linge branché sur le circuit prises : un courant de fuite de \SI{80}{mA} s'écoule vers la terre. Quel élément déclenche ? Pourquoi les disjoncteurs divisionnaires ne réagissent-ils pas ?
\item Le propriétaire affirme : « le différentiel coupe tout seul, la terre ne sert à rien ». Construire un argumentaire scientifique en trois points (boucle de défaut, impédance de la prise de terre, tension de contact) pour démontrer que la terre est indispensable au bon fonctionnement de la protection différentielle.
\end{enumerate}
\end{exercice}

\begin{exercice}[Type BEPC : intervention sur une prise suspecte]
Dans un atelier de Douala, on dispose d'un tournevis testeur, d'une prise de courant fonctionnelle (Phase, Neutre et Terre identifiés) et d'une prise suspecte dont les trois fils ne sont pas repérés. Un élève, Bienvenue, veut changer cette prise suspecte. Il coupe simplement l'interrupteur de la lampe de la pièce (alimentée par le même disjoncteur) et commence à dévisser la prise.
\begin{enumerate}
\item Expliquer pourquoi l'action de Bienvenue est dangereuse, alors que la lampe est éteinte.
\item Citer et expliquer les trois règles d'or oubliées : coupure à la source amont, vérification d'absence de tension (VAT), verrouillage et signalisation.
\item Décrire la procédure pas à pas permettant d'identifier la Phase, le Neutre et la Terre sur la prise suspecte, en utilisant uniquement le tournevis testeur et la prise fonctionnelle comme référence.
\item Préciser les limites de cette méthode concernant l'identification de la Terre, et proposer une précaution complémentaire avant toute remise en service.
\end{enumerate}
\end{exercice}

\newpage

\coursec{Corrigés des exercices}

\coursec{Corrigés des exercices — Leçon 21 : Généralités sur les circuits électriques domestiques}

\begin{corrige}[Exercice 1 — QCM : le branchement dangereux]
\textbf{Réponse exacte : B).}

\medskip
\textbf{Justification.} Le fil bleu (couleur normalisée du Neutre) est branché sur la borne L et le fil rouge (Phase) sur la borne N : Phase et Neutre sont \emph{inversés}. L'appareil fonctionne car il « voit » toujours une tension de \SI{230}{V} entre ses deux bornes d'alimentation, ce qui élimine la réponse A. Mais l'interrupteur interne de l'appareil, conçu pour couper la Phase, coupe ici le Neutre : même appareil « éteint », ses composants internes restent portés au potentiel de la Phase (\SI{230}{V} par rapport à la terre). Toute intervention ou contact interne devient alors dangereux.

\medskip
\textbf{Pourquoi les autres réponses sont fausses :}
\begin{itemize}
\item \textbf{C)} est faux : l'énoncé précise que le fil vert/jaune est bien sur la borne Terre.
\item \textbf{D)} est faux : le différentiel compare le courant de Phase et le courant de Neutre ; l'inversion des fils ne l'empêche pas de détecter une fuite à la terre.
\end{itemize}
\end{corrige}

\begin{attention}[Point de vigilance]
Au BEPC, ne jamais écrire « l'appareil ne marche plus » : une inversion Phase/Neutre laisse l'appareil fonctionnel, c'est justement ce qui rend le défaut \emph{invisible} et dangereux.
\end{attention}

\begin{corrige}[Exercice 2 — Vrai ou faux]
\begin{enumerate}
\item \textbf{FAUX.} La couleur normalisée de la Phase est le rouge, le marron ou le noir ; le bleu est \emph{exclusivement} réservé au Neutre.
\item \textbf{FAUX.} Le tournevis testeur distingue la Phase (néon allumé) du Neutre (néon éteint), mais il ne permet pas de différencier de façon fiable le Neutre de la Terre, car aucun des deux ne fait normalement briller le néon.
\item \textbf{VRAI.} Le disjoncteur (ou interrupteur) différentiel de \SI{30}{mA} détecte les courants de fuite vers la terre, typiques des contacts indirects (carcasse métallique devenue sous tension), et coupe le circuit avant que le courant traversant le corps ne devienne mortel.
\item \textbf{FAUX.} L'interrupteur de la lampe ne coupe que le circuit d'éclairage (et parfois seulement un des deux fils) ; la prise reste alimentée en permanence par son propre départ. Il faut couper le disjoncteur divisionnaire du circuit concerné, puis faire une VAT (vérification d'absence de tension).
\end{enumerate}
\end{corrige}

\begin{corrige}[Exercice 3 — Définitions et code couleur]
\textbf{1. Définitions.}
\begin{itemize}
\item \cle{Disjoncteur divisionnaire} : appareil de protection placé en tête de chaque circuit (départ) du tableau électrique, qui coupe automatiquement le circuit en cas de surcharge ou de court-circuit.
\item \cle{Conducteur de protection (terre)} : conducteur vert/jaune reliant les masses métalliques des appareils à la prise de terre, qui évacue les courants de défaut et permet au différentiel de déclencher.
\item \cle{Schéma unifilaire} : représentation simplifiée d'une installation électrique où chaque ensemble de conducteurs (Phase, Neutre, Terre) est figuré par un seul trait, avec les appareils de protection et leurs calibres.
\end{itemize}

\textbf{2. Code couleur normalisé.}
\begin{center}
\begin{tabularx}{\linewidth}{|l|X|}
\hline
\rowcolor{popblueL} \textbf{Conducteur} & \textbf{Couleur normalisée} \\
\hline
Phase (L) & Rouge, marron ou noir \\
\hline
Neutre (N) & Bleu clair (obligatoirement) \\
\hline
Terre (PE) & Bicolore vert/jaune (obligatoirement) \\
\hline
\end{tabularx}
\end{center}
\end{corrige}

\begin{attention}[Erreur fréquente]
Le vert/jaune est réservé \emph{exclusivement} à la Terre : il est interdit de l'utiliser pour la Phase ou le Neutre, même « provisoirement ».
\end{attention}

\begin{corrige}[Exercice 4 — Calcul direct]
\textbf{1. Intensité absorbée.}

Pour un appareil monophasé, la puissance s'écrit $P = U \times I$, donc :
\[
I = \frac{P}{U} = \frac{\SI{4600}{W}}{\SI{230}{V}} = \SI{20}{A}.
\]

\textbf{2. Choix du disjoncteur de \SI{20}{A}.}

Le calibre du disjoncteur doit être \emph{au moins égal} au courant absorbé en service normal. Ici $I = \SI{20}{A}$ est exactement égal au calibre : le disjoncteur fonctionne en permanence à sa limite de déclenchement. La moindre surintensité (démarrage, tolérance du réseau, échauffement) provoquera des déclenchements intempestifs.

\textbf{Conclusion :} non, un disjoncteur de \SI{20}{A} est inadapté. Conformément au cours, une cuisinière se raccorde sur un circuit spécialisé protégé par un disjoncteur de \SI{32}{A} avec un câble de section \SI{6}{mm^2}.
\end{corrige}

\begin{attention}[Point de vigilance]
On ne choisit jamais un calibre \emph{inférieur} au courant d'emploi, et on évite un calibre \emph{exactement égal} : il faut une marge. Mais on ne choisit pas non plus un calibre excessif (voir exercice 8).
\end{attention}

\begin{corrige}[Exercice 5 — Identifier les éléments d'une installation]
\textbf{1. Ordre de traversée du courant (du réseau vers les prises) :}
\[
\text{Compteur \SI{45}{A}} \;\longrightarrow\; \text{Disjoncteur d'abonné} \;\longrightarrow\; \text{Interrupteur différentiel \SI{30}{mA}} \;\longrightarrow\; \text{Disjoncteurs divisionnaires} \;\longrightarrow\; \text{prises et luminaires.}
\]

\textbf{2. Rôle de chaque élément :}
\begin{itemize}
\item \textbf{Compteur (\SI{45}{A})} : appareil appartenant au distributeur (ENEO) qui mesure l'énergie électrique consommée (en kWh) pour la facturation ; son calibre indique la puissance maximale souscrite.
\item \textbf{Disjoncteur d'abonné} : protection générale de l'installation, placée juste après le compteur ; il permet de couper toute l'installation et protège contre les surcharges et courts-circuits généraux.
\item \textbf{Interrupteur différentiel \SI{30}{mA}} : compare le courant qui part par la Phase et celui qui revient par le Neutre ; dès qu'une fuite vers la terre dépasse \SI{30}{mA}, il coupe l'alimentation en une fraction de seconde.
\item \textbf{Disjoncteurs divisionnaires (\SI{10}{A}, \SI{16}{A}, \SI{32}{A})} : protègent chaque circuit (éclairage, prises, four) contre les surcharges et les courts-circuits, en fonction de la section des câbles.
\end{itemize}

\textbf{3. Qui protège quoi ?}
\begin{itemize}
\item \textbf{Protection des personnes} : l'interrupteur différentiel \SI{30}{mA} (associé à la prise de terre).
\item \textbf{Protection des circuits et des appareils (biens)} : le disjoncteur d'abonné et les disjoncteurs divisionnaires.
\end{itemize}
\end{corrige}

\begin{corrige}[Exercice 6 — Sections de câbles et calibres]
\textbf{Tableau complété :}
\begin{center}
\begin{tabularx}{\linewidth}{|X|l|l|}
\hline
\rowcolor{popblueL} \textbf{Circuit} & \textbf{Calibre du disjoncteur} & \textbf{Section minimale} \\
\hline
Éclairage & \SI{10}{A} & \SI{1,5}{mm^2} \\
\hline
Prises de courant ordinaires & \SI{16}{A} & \SI{2,5}{mm^2} \\
\hline
Four ou cuisinière & \SI{32}{A} & \SI{6}{mm^2} \\
\hline
\end{tabularx}
\end{center}

\textbf{Signification de la désignation « 3G1,5 » :}
\begin{itemize}
\item le chiffre \textbf{3} : le câble contient \textbf{trois conducteurs} (Phase, Neutre, Terre) ;
\item la lettre \textbf{G} : l'un des conducteurs est le conducteur de protection \textbf{vert/jaune} (terre) ;
\item le nombre \textbf{1,5} : la \textbf{section} de chaque conducteur est de \SI{1,5}{mm^2}.
\end{itemize}
Ainsi un câble 3G2,5 pour les prises contient trois conducteurs de \SI{2,5}{mm^2}, dont un vert/jaune.
\end{corrige}

\begin{corrige}[Exercice 7 — Utilisation du tournevis testeur]
\textbf{1. Manipulation correcte, étape par étape :}
\begin{enumerate}
\item Vérifier d'abord le bon état du testeur sur une prise dont on est sûr qu'elle est sous tension.
\item Saisir le tournevis testeur par son manche isolant, \textbf{sans jamais toucher la lame métallique}.
\item Poser \textbf{un doigt sur le capuchon métallique} de l'extrémité du manche : le corps de l'utilisateur ferme le circuit vers la terre à travers la résistance interne du testeur.
\item Introduire la lame dans l'alvéole de la borne à tester (Phase puis Neutre), sans toucher les deux alvéoles en même temps.
\item Observer le néon : allumé = présence de la Phase ; éteint = absence de Phase (Neutre, Terre ou fil hors tension).
\end{enumerate}

\textbf{2. Conclusion.} Le néon s'allume sur la borne de Phase et reste éteint sur la borne de Neutre : la prise est \textbf{correctement alimentée} et le repérage Phase/Neutre semble conforme. On peut y brancher le groupe électrogène de secours (hors tension au moment du branchement).

\textbf{3. Néon faiblement allumé sur la borne de Terre.} \textbf{Non}, on ne peut pas conclure que la Terre est en bon état. Une lueur faible peut provenir d'une tension induite (couplage capacitif entre fils voisins) ou d'un défaut d'isolement, et le testeur ne mesure ni la continuité ni la valeur de la résistance de la prise de terre. Seul un appareil de mesure adapté (contrôleur de terre, multimètre utilisé par un professionnel) permet de vérifier que la prise de terre est effectivement reliée et de valeur suffisamment faible.
\end{corrige}

\begin{attention}[Limite du tournevis testeur]
Le testeur détecte une \emph{présence de tension} par rapport à la terre ; il ne mesure ni la valeur de la tension, ni la continuité des conducteurs, ni la qualité de la prise de terre.
\end{attention}

\begin{corrige}[Exercice 8 — Choix d'une protection]
\textbf{1. Intensité absorbée en fonctionnement normal :}
\[
I = \frac{P}{U} = \frac{\SI{3500}{W}}{\SI{230}{V}} \approx \SI{15,2}{A}.
\]

\textbf{2. Choix du calibre.} Le calibre doit être supérieur au courant d'emploi ($\SI{15,2}{A}$) tout en restant compatible avec la section du câble :
\begin{itemize}
\item \SI{10}{A} : trop faible ($10 < 15,2$), déclenchement permanent ;
\item \SI{16}{A} : correct, c'est le calibre immédiatement supérieur, adapté à un câble de \SI{2,5}{mm^2} ;
\item \SI{20}{A} et \SI{32}{A} : trop élevés pour un câble de \SI{2,5}{mm^2}, ils laisseraient le câble chauffer dangereusement avant de déclencher.
\end{itemize}
\textbf{Choix : disjoncteur de \SI{16}{A}} avec un câble 3G2,5.

\textbf{3. Pourquoi « remonter » à \SI{32}{A} est dangereux.} Si le disjoncteur de \SI{16}{A} déclenche souvent, c'est un \emph{signal d'alarme} : surcharge réelle du circuit, défaut d'isolement naissant ou appareil défectueux. Augmenter le calibre à \SI{32}{A} supprime le symptôme mais pas la cause : le câble de \SI{2,5}{mm^2}, prévu pour \SI{16}{A}, pourra alors être traversé durablement par un courant excessif, s'échauffer, dégrader son isolant et provoquer un \textbf{incendie} — le disjoncteur ne protégeant plus rien.

\textbf{Bonne démarche :} rechercher la cause des déclenchements (mesurer le courant réel absorbé, vérifier l'état de la scie et du câblage), répartir les appareils sur plusieurs circuits, ou créer un circuit spécialisé correctement dimensionné (câble \SI{6}{mm^2} si un calibre de \SI{32}{A} est réellement nécessaire).
\end{corrige}

\begin{corrige}[Exercice 9 — Type BEPC : le circuit four de la cuisine]
\textbf{1. Intensité absorbée par le four :}
\[
I = \frac{P}{U} = \frac{\SI{7200}{W}}{\SI{230}{V}} \approx \SI{31,3}{A}.
\]

\textbf{2. Choix du disjoncteur de \SI{32}{A}.} On compare : $I \approx \SI{31,3}{A} < \SI{32}{A}$. Le calibre est immédiatement supérieur au courant d'emploi et correspond à la section de \SI{6}{mm^2} : le disjoncteur de \SI{32}{A} est \textbf{correctement choisi}.

\textbf{3. Résistance du câble et chute de tension.}

Longueur totale aller-retour : $2L = \SI{40}{m}$ ; section $S = \SI{6}{mm^2} = \SI{6e-6}{m^2}$.
\begin{align*}
R &= \rho \times \frac{2L}{S} = \SI{1,7e-8}{\ohm.m} \times \frac{\SI{40}{m}}{\SI{6e-6}{m^2}} \approx \SI{0,113}{\ohm}.\\
\Delta U &= R \times I = \num{0,113} \times \num{31,3} \approx \SI{3,5}{V}.
\end{align*}

\textbf{4. Pourcentage et conformité :}
\[
\frac{\Delta U}{U} = \frac{\num{3,5}}{\num{230}} \approx \num{0,0154} \approx 1,5\,\%.
\]
Or $1,5\,\% < 3\,\%$ : la chute de tension est inférieure à la limite admissible. \textbf{L'installation est conforme.}

\textbf{5. Nécessité du conducteur de terre.} Le four possède une carcasse métallique (appareil de Classe I). En cas de défaut d'isolement, un fil sous tension peut toucher cette carcasse. Sans terre, la carcasse resterait portée à \SI{230}{V} : toute personne la touchant serait électrisée. Le conducteur de terre (vert/jaune) crée un chemin de faible résistance vers la prise de terre : le courant de défaut s'y écoule, le différentiel de \SI{30}{mA} le détecte immédiatement et coupe le circuit avant tout danger mortel. La terre est donc \textbf{obligatoire} sur ce circuit.

\medskip
\textbf{Barème indicatif (sur 10) :} question 1 : 1,5 pt (formule + application) ; question 2 : 1,5 pt (comparaison chiffrée exigée) ; question 3 : 3 pts (conversion de la section 1 pt, calcul de $R$ 1 pt, $\Delta U$ 1 pt) ; question 4 : 2 pts (pourcentage + conclusion de conformité) ; question 5 : 2 pts (défaut d'isolement + rôle conjoint terre/différentiel).
\end{corrige}

\begin{attention}[Points de vigilance]
\begin{itemize}
\item Ne pas oublier de \textbf{doubler la longueur} (aller-retour Phase/Neutre) : oublier le facteur 2 divise $\Delta U$ par deux et fausse la conclusion.
\item Convertir la section : $\SI{6}{mm^2} = \SI{6e-6}{m^2}$ ; laisser $S$ en mm\textsuperscript{2} donne un résultat $10^6$ fois trop grand.
\item La justification de la question 2 doit être \emph{chiffrée} : écrire « c'est bon » sans comparer \num{31,3} et \num{32} ne rapporte pas les points.
\end{itemize}
\end{attention}

\begin{corrige}[Exercice 10 — Type BEPC : schéma unifilaire d'une maison]
\textbf{1. Schéma unifilaire complété :}
\begin{center}
\begin{tabularx}{\linewidth}{|X|l|l|l|}
\hline
\rowcolor{popblueL} \textbf{Départ} & \textbf{Calibre} & \textbf{Section} & \textbf{Câble} \\
\hline
Éclairage & \SI{10}{A} & \SI{1,5}{mm^2} & 3G1,5 \\
\hline
Prises de courant & \SI{16}{A} & \SI{2,5}{mm^2} & 3G2,5 \\
\hline
Four / cuisinière & \SI{32}{A} & \SI{6}{mm^2} & 3G6 \\
\hline
\end{tabularx}
\end{center}
Chaîne complète : Compteur \SI{45}{A} $\rightarrow$ Disjoncteur d'abonné $\rightarrow$ Interrupteur différentiel \SI{30}{mA} $\rightarrow$ [Départ \SI{10}{A} en 3G1,5 ; Départ \SI{16}{A} en 3G2,5 ; Départ \SI{32}{A} en 3G6].

\textbf{2. Conducteurs de chaque départ :} chaque départ comporte \textbf{trois conducteurs} : la Phase (rouge, marron ou noir), le Neutre (bleu clair) et la Terre (vert/jaune). D'où la lettre G dans la désignation des câbles.

\textbf{3. Défaut d'isolement de \SI{80}{mA} sur le lave-linge.}

C'est l'\textbf{interrupteur différentiel \SI{30}{mA}} qui déclenche : le courant de fuite ($\SI{80}{mA} > \SI{30}{mA}$) crée une différence entre le courant sortant par la Phase et le courant revenant par le Neutre, ce qui provoque l'ouverture immédiate du circuit.

Les disjoncteurs divisionnaires ne réagissent pas car ils ne sont sensibles qu'aux \textbf{surcharges et courts-circuits}, c'est-à-dire à des courants supérieurs à leur calibre (\SI{16}{A} pour le circuit prises). Or $\SI{80}{mA} = \SI{0,08}{A} \ll \SI{16}{A}$ : ce courant de fuite est 200 fois plus faible que le seuil de déclenchement du divisionnaire. Il est pourtant déjà mortel pour une personne — c'est précisément pour cela que le différentiel existe.

\textbf{4. Argumentaire : la terre est indispensable au différentiel.}
\begin{itemize}
\item \textbf{Boucle de défaut :} le différentiel ne peut détecter une fuite que si un courant de défaut \emph{circule réellement}. Sans conducteur de terre relié à la carcasse, aucun courant ne s'écoule lors du défaut : la carcasse monte silencieusement à \SI{230}{V} et le différentiel ne voit rien. C'est la \emph{personne} qui touchera la carcasse qui fermera la boucle — le différentiel coupera alors, mais le courant aura déjà traversé le corps.
\item \textbf{Impédance de la prise de terre :} avec une prise de terre de faible résistance (bonne liaison vert/jaune), le courant de défaut est important et franchit instantanément le seuil de \SI{30}{mA} : le déclenchement est rapide et sûr, dès l'apparition du défaut, avant tout contact humain.
\item \textbf{Tension de contact :} la terre maintient la carcasse à un potentiel proche de zéro ; la tension de contact $U_c = R_{\text{terre}} \times I_{\text{défaut}}$ reste alors inférieure à la tension de sécurité (\SI{50}{V} en locaux secs) pendant le bref instant précédant la coupure. Sans terre, la tension de contact vaut \SI{230}{V} : mortelle.
\end{itemize}
\textbf{Conclusion :} terre et différentiel forment un \emph{couple indissociable} : la terre canalise le défaut, le différentiel le détecte et coupe. Affirmer que la terre « ne sert à rien » est scientifiquement faux et mortellement dangereux.

\medskip
\textbf{Barème indicatif (sur 10) :} question 1 : 2,5 pts ; question 2 : 1,5 pt ; question 3 : 3 pts (élément qui déclenche 1 pt, comparaison \num{0,08} A et \SI{16}{A} 2 pts) ; question 4 : 3 pts (1 pt par argument correctement formulé).
\end{corrige}

\begin{attention}[Points de vigilance]
\begin{itemize}
\item Question 3 : la comparaison \emph{chiffrée} entre le courant de fuite et le calibre du divisionnaire est la justification attendue ; convertir \SI{80}{mA} en \SI{0,08}{A}.
\item Question 4 : les trois points imposés par l'énoncé (boucle de défaut, impédance, tension de contact) doivent apparaître explicitement.
\end{itemize}
\end{attention}

\begin{corrige}[Exercice 11 — Type BEPC : intervention sur une prise suspecte]
\textbf{1. Pourquoi l'action de Bienvenue est dangereuse.}

L'interrupteur de la lampe n'agit que sur le fil de commande de l'éclairage ; il ne coupe en rien l'alimentation des prises de la pièce, qui possèdent leur propre départ (ou restent alimentées en amont de l'interrupteur). De plus, un interrupteur simple ne coupe parfois que le Neutre si le montage est défectueux. La prise suspecte reste donc \textbf{sous tension} : en dévissant, Bienvenue risque le contact direct avec la Phase — électrisation, voire électrocution.

\textbf{2. Les trois règles d'or oubliées :}
\begin{itemize}
\item \textbf{Coupure à la source amont :} ouvrir le disjoncteur divisionnaire du circuit concerné (ou le disjoncteur général d'abonné), et non un simple interrupteur d'usage.
\item \textbf{Vérification d'absence de tension (VAT) :} après coupure, s'assurer avec un appareil adapté (testeur vérifié avant et après sur une source sûre) qu'aucune borne de la prise n'est sous tension. On ne travaille jamais « à l'aveugle ».
\item \textbf{Verrouillage et signalisation :} empêcher toute remise en tension intempestive par un tiers (cadenas ou consignation du disjoncteur, affichage « intervention en cours — ne pas réenclencher »).
\end{itemize}

\textbf{3. Procédure d'identification des trois fils (circuit remis sous tension, prise démontée, extrémités accessibles et écartées, sans contact) :}
\begin{enumerate}
\item Vérifier le bon fonctionnement du tournevis testeur sur la \textbf{prise fonctionnelle} : le néon doit s'allumer sur sa borne de Phase (référence).
\item Sur la prise suspecte, doigt sur le capuchon métallique, toucher successivement chacun des trois fils avec la lame.
\item Le fil qui fait \textbf{franchement briller le néon} est la \textbf{Phase} (à raccorder sur la borne L).
\item Les deux fils qui laissent le néon éteint sont le Neutre et la Terre, indiscernables au testeur : on les repère provisoirement.
\item Couper à nouveau le circuit (VAT), raccorder la Phase identifiée sur la borne L, et l'un des deux fils restants sur la borne N en respectant si possible la couleur d'origine (bleu = Neutre, vert/jaune = Terre) ; remettre sous tension et contrôler le fonctionnement.
\end{enumerate}

\textbf{4. Limites et précaution complémentaire.}

Le tournevis testeur ne distingue \textbf{pas} le Neutre de la Terre : tous deux laissent le néon éteint, et une lueur faible (tension induite) peut induire en erreur. Si les couleurs des fils sont douteuses ou effacées, l'identification de la Terre reste incertaine — or une inversion Neutre/Terre rend la protection différentielle inefficace sur les contacts indirects.

\textbf{Précaution complémentaire avant remise en service :} faire vérifier la continuité et la valeur de la prise de terre par un appareil de mesure approprié (multimètre en mode continuité, contrôleur de terre) ou par un électricien qualifié, puis contrôler le déclenchement du différentiel (bouton test). On ne remet l'installation en service définitif qu'après ces vérifications.

\medskip
\textbf{Barème indicatif (sur 10) :} question 1 : 2 pts ; question 2 : 3 pts (1 pt par règle nommée et expliquée) ; question 3 : 3 pts (référence sur prise saine 1 pt, identification de la Phase 1 pt, organisation logique 1 pt) ; question 4 : 2 pts (limite Neutre/Terre 1 pt, précaution 1 pt).
\end{corrige}

\begin{attention}[Points de vigilance]
\begin{itemize}
\item La VAT doit être faite \emph{après} coupure et l'appareil de contrôle \emph{vérifié avant et après} sur une source sûre : citer cette double vérification.
\item Ne jamais écrire qu'on peut identifier la Terre « parce que le néon brille faiblement » : c'est la limite même de la méthode.
\item Toute intervention se fait \textbf{hors tension} ; le test sous tension de la question 3 est une mesure de \emph{repérage} exceptionnelle, avec précautions.
\end{itemize}
\end{attention}

\newpage

\coursec{Fiche bilan}

\begin{popfigure}
\begin{center}\tcbox[colback=popblue,colframe=popblue,coltext=white,arc=9pt,boxsep=4pt,left=6pt,right=6pt]{\bfseries\small Circuits électriques domestiques}\end{center}
\vspace{-2pt}
\begin{tcbraster}[raster columns=3,raster equal height=rows,raster column skip=6pt,raster row skip=6pt,enhanced,blankest]
\begin{tcolorbox}[enhanced,colback=poporange,colframe=poporange,coltext=white,boxrule=0.9pt,arc=3pt,left=4pt,right=4pt,top=3pt,bottom=3pt,boxsep=1pt,halign=left]\bfseries\scriptsize Installation : compteur, disjoncteur, tableau, circuits\end{tcolorbox}
\begin{tcolorbox}[enhanced,colback=popgreen,colframe=popgreen,coltext=white,boxrule=0.9pt,arc=3pt,left=4pt,right=4pt,top=3pt,bottom=3pt,boxsep=1pt,halign=left]\bfseries\scriptsize Conducteurs : phase, neutre, terre ; code couleur\end{tcolorbox}
\begin{tcolorbox}[enhanced,colback=poppurple,colframe=poppurple,coltext=white,boxrule=0.9pt,arc=3pt,left=4pt,right=4pt,top=3pt,bottom=3pt,boxsep=1pt,halign=left]\bfseries\scriptsize Prise 2P+T : phase, neutre, terre\end{tcolorbox}
\begin{tcolorbox}[enhanced,colback=poppink,colframe=poppink,coltext=white,boxrule=0.9pt,arc=3pt,left=4pt,right=4pt,top=3pt,bottom=3pt,boxsep=1pt,halign=left]\bfseries\scriptsize Tournevis testeur : repérer la phase\end{tcolorbox}
\begin{tcolorbox}[enhanced,colback=popgold,colframe=popgold,coltext=white,boxrule=0.9pt,arc=3pt,left=4pt,right=4pt,top=3pt,bottom=3pt,boxsep=1pt,halign=left]\bfseries\scriptsize Sécurité : fusible, disjoncteur, terre, gestes\end{tcolorbox}
\begin{tcolorbox}[enhanced,colback=popblue!60,colframe=popblue!60,coltext=white,boxrule=0.9pt,arc=3pt,left=4pt,right=4pt,top=3pt,bottom=3pt,boxsep=1pt,halign=left]\bfseries\scriptsize Schémas : unifilaire, lecture de plan\end{tcolorbox}
\end{tcbraster}

\legende{Carte mentale de la leçon : les six idées essentielles sur l'installation électrique domestique.}
\end{popfigure}

\begin{definition}[Définitions clés]
\begin{itemize}
  \item \cle{Installation domestique} : ensemble des équipements qui amènent et distribuent l'énergie électrique dans une habitation (branchement ENEO, compteur, disjoncteur général, tableau de répartition, circuits, prises et points lumineux).
  \item \cle{Phase} (L) : conducteur qui transporte le courant depuis le réseau ; c'est le conducteur \emph{dangereux}, à la tension de \SI{220}{V} par rapport à la terre.
  \item \cle{Neutre} (N) : conducteur de retour du courant ; il est relié à la terre côté transformateur, sa tension par rapport à la terre est voisine de \SI{0}{V}.
  \item \cle{Terre} (T) : conducteur de protection relié au sol ; il évacue les courants de défaut pour protéger les personnes.
  \item \cle{Prise 2P+T} : prise à deux pôles (phase + neutre) complétée par une borne de terre.
  \item \cle{Tournevis testeur} : petit outil contenant une lampe néon qui s'allume lorsque sa pointe touche la phase et que l'utilisateur touche le capuchon métallique.
\end{itemize}
\end{definition}

\begin{aretenir}[Code couleur et rôles des conducteurs]
\begin{center}
\begin{tabularx}{\linewidth}{|l|X|X|}\hline
\rowcolor{popblueL} \textbf{Conducteur} & \textbf{Couleur du fil} & \textbf{Rôle} \\ \hline
Phase (L) & rouge, marron ou noir & transporte le courant vers l'appareil \\ \hline
Neutre (N) & bleu clair & assure le retour du courant \\ \hline
Terre (T) & vert/jaune & protection des personnes \\ \hline
\end{tabularx}
\end{center}
La tension délivrée entre phase et neutre dans les habitations au Cameroun est de \SI{220}{V} (courant alternatif, fréquence \SI{50}{Hz}).
\end{aretenir}

\begin{methode}[Méthodes express]
\begin{enumerate}
  \item \textbf{Brancher une prise 2P+T} : par convention, la phase se place sur la borne de droite (prise vue de face), le neutre sur la borne de gauche, et le fil de terre sur la borne centrale (broche ou ergots). Serrer les vis fermement : aucun brin dénudé ne doit dépasser.
  \item \textbf{Utiliser le tournevis testeur} : circuit sous tension, poser la pointe sur la borne à tester et un doigt sur le capuchon métallique. Si la lampe s'allume, c'est la phase ; si elle reste éteinte, c'est le neutre (ou la terre).
  \item \textbf{En cas d'accident électrique} : couper d'abord le courant au disjoncteur général, ne jamais toucher directement la victime tant qu'elle est en contact avec le circuit, puis alerter les secours.
  \item \textbf{Règle d'or} : on ne travaille jamais sur un circuit sous tension ; on coupe au disjoncteur, puis on vérifie l'absence de tension au testeur.
\end{enumerate}
\end{methode}

\begin{attention}[Limites du tournevis testeur]
Le testeur indique seulement la \emph{présence} de la phase : il ne \textbf{mesure pas} la valeur de la tension. Une lampe éteinte ne garantit pas l'absence de danger (testeur défectueux, mauvais contact du doigt). Il ne remplace ni le multimètre ni les précautions : on coupe toujours le circuit avant d'intervenir.
\end{attention}

\begin{aretenir}[Checklist avant le BEPC]
\begin{itemize}
  \item[$\square$] Je sais citer les éléments d'une installation : compteur, disjoncteur général, tableau de répartition, circuits.
  \item[$\square$] Je connais le code couleur : phase (rouge/marron/noir), neutre (bleu), terre (vert/jaune).
  \item[$\square$] Je sais que la tension domestique au Cameroun est de \SI{220}{V}.
  \item[$\square$] Je sais nommer les trois bornes d'une prise 2P+T et leur rôle.
  \item[$\square$] Je sais expliquer le rôle protecteur de la prise de terre.
  \item[$\square$] Je sais utiliser un tournevis testeur et interpréter le résultat.
  \item[$\square$] Je connais les limites du testeur : il ne mesure pas la tension et peut être trompeur.
  \item[$\square$] Je sais distinguer fusible et disjoncteur divisionnaire.
  \item[$\square$] Je connais les dangers : électrisation, incendie, court-circuit.
  \item[$\square$] Je sais les premiers gestes : couper le courant avant tout secours.
  \item[$\square$] Je sais lire un schéma unifilaire simple (symboles prise, lampe, interrupteur).
  \item[$\square$] Je n'oublie jamais : jamais de travail sur un circuit sous tension.
\end{itemize}
\end{aretenir}

\findecours

\end{document}
